\documentclass[10pt,a4paper]{article}

% ----------------- Paquetes -------------------%
\usepackage[T1]{fontenc}
\usepackage[utf8]{inputenc}
\usepackage[a4paper,margin=2.5cm]{geometry}
\usepackage{mathptmx}             
\usepackage{changepage} 
\usepackage{setspace}
\usepackage[spanish]{babel}
\usepackage[labelfont=bf,labelsep=space]{caption}
\usepackage{amssymb}
\usepackage{amsmath}
\usepackage{ebgaramond}
\usepackage{placeins}
\usepackage{etoolbox}
\usepackage{setspace}
\usepackage{parskip} 
\usepackage{tcolorbox}
\usepackage{needspace}
\usepackage{xparse}
\usepackage{ocgx2}
\usepackage{hyperref}
\usepackage{hypcap}



%%%%%%%%%%%%%%%%%%%%%%%%%%%%%%%%%%%%%%%%%%%%%%%%%%%%%%%%%%%%%%%%
% --- Fija primero tus valores globales "buenos" ---
\setstretch{1.2}                 % Interlineado global
\setlength{\parskip}{0.7\baselineskip}
\setlength{\parindent}{0pt}
\newlength{\GlobalParskip}
\setlength{\GlobalParskip}{\parskip}

\newcounter{eqnum}[subsection]
\renewcommand{\theeqnum}{\arabic{eqnum}}
\newcounter{alphaitem}
\newcounter{numitem}

%% ------------------- Recuadros----------------------%%
\newcommand{\puntosclave}[1]{%
  \begin{tcolorbox}[
    colframe=black,
    title={Puntos clave},
    before upper={\setlength{\parindent}{0.5cm}\setlength{\parskip}{0.5\baselineskip}}
  ]
  #1
  \end{tcolorbox}
}

\newcommand{\modificaciones}[1]{
  \begin{tcolorbox}[
    colback=white,
    colframe=red,
    title={Modificaciones a realizar},
    before upper={\setlength{\parindent}{0.5cm}\setlength{\parskip}{0.5\baselineskip}}
  ]
  #1
  \end{tcolorbox}
}


%% ------------------- Preguntas TEST----------------------%%
% Opciones: radio (single-choice)
\NewDocumentCommand{\OptRadio}{m m m}{%
  \noindent\CheckBox[radio,name=#1,value=#2,width=1.2ex,height=1.2ex]{}%
  \;\textbf{#2.}~#3\par
}

% Opciones: checkbox (multi-choice)
\NewDocumentCommand{\OptCheck}{m m}{%
  \noindent\CheckBox[name=#1,width=1.2ex,height=1.2ex,bordercolor={0 0 0}]{}%
  \;#2\par
}

% Pregunta single-choice (radio)
% Args: 1=id, 2=enunciado, 3=opciones, 4=solución+justificación, 5=imagen (o {}), 6=origen
\NewDocumentCommand{\PreguntaTestSingle}{m m m m m m}{%
  \par\medskip
  \noindent\textbf{#1.}~#2\par\smallskip

    % Imagen (si no es {})
    \IfBlankTF{#5}{}{%
      \begin{center}
        \includegraphics[width=0.75\linewidth]{#5}
      \end{center}
    }

    \begin{Form}
      #3
    \end{Form}

    \par\smallskip
    \switchocg{sol-#1}{\fbox{Mostrar / ocultar solución}}%
    \hfill{\footnotesize #6}\par\smallskip

    \begin{ocg}{Solución}{sol-#1}{0}
      \noindent #4
    \end{ocg}
  \par\medskip
}

% Pregunta multi-choice (checkboxes)
\NewDocumentCommand{\PreguntaTestMulti}{m m m m m m}{%
  \par\medskip
  \noindent\textbf{#1.}~#2\par\smallskip

    % Imagen (si no es {})
    \IfBlankTF{#5}{}{%
      \begin{center}
        \includegraphics[width=0.75\linewidth]{#5}
      \end{center}
    }
    \begin{Form}
      #3
    \end{Form}

    \par\smallskip
    \switchocg{sol-#1}{\fbox{Mostrar / ocultar solución}}%
    \hfill{\footnotesize #6}\par\smallskip

    \begin{ocg}{Solución}{sol-#1}{0}
      \noindent #4
    \end{ocg}
  \par\medskip
}

%% ------------------- Listas----------------------%%
    % --- Lista alfabética ---
    \newenvironment{la}
      {\begin{list}{\alph{alphaitem})}{%
          \usecounter{alphaitem}%
          \setlength{\leftmargin}{1cm}%
          \setlength{\parskip}{3pt}% 
          \setlength{\itemsep}{0pt}% ← espacio entre ítems
          \setlength{\parsep}{0pt}%  ← párrafos dentro de un ítem
      }}
      {\end{list}}
      
    % --- Lista numérica ---
    \newenvironment{lnum}
      {\begin{list}{\arabic{numitem}.}{%
          \usecounter{numitem}%
          \setlength{\leftmargin}{1cm}%
          \setlength{\parskip}{3pt}% 
          \setlength{\itemsep}{0pt}% ← espacio entre ítems
          \setlength{\parsep}{0pt}%  ← párrafos dentro de un ítem
      }}
      {\end{list}}
      
% ------------------- Imágenes----------------------%
    \graphicspath{{Img/}}% Rutas por defecto 
    % --- Imagen adaptativa ---
    % Registros de dimensión definidos una sola vez
    \newdimen\imgcap
    \newdimen\imgmin
    \newdimen\imgmax
    \newdimen\imgremain
    \newdimen\imgh
    \newdimen\imgneed
    
    \newcommand{\imagenfit}[3]{%
      \addvspace{10pt}% separación antes
      \par\begingroup
        % Parámetros de composición (asignaciones, no \newdimen)
        \imgcap=1\baselineskip            % alto estimado de título+fuente
        \imgmin=0.15\textheight
        \imgmax=0.15\textheight
        \imgremain=\pagegoal
        \ifdim\pagegoal=\maxdimen \imgremain=0pt\fi
        \advance\imgremain by -\pagetotal
        \advance\imgremain by -\imgcap
        % Decisión de altura destino
        \ifdim\imgremain<\imgmin
          \newpage
          \imgh=0.20\textheight
        \else
          \imgh=\imgremain
          \ifdim\imgh>\imgmax \imgh=\imgmax \fi
        \fi
        % Reservar espacio para evitar cortes del bloque completo
        \imgneed=\imgh \advance\imgneed by \imgcap
        \Needspace{\imgneed}
        % Cuerpo
        \noindent\begin{minipage}{\textwidth}
          \centering
          \captionof{figure}{\textemdash\ #2}\par\vspace{0.3em}% título arriba
          \includegraphics[width=\linewidth,height=\imgh,keepaspectratio]{\detokenize{#1}}\par
          \vspace{0.3em}\small\textit{Fuente: }#3% fuente abajo
        \end{minipage}%
      \endgroup
      \par\addvspace{10pt}%
    }

%------------ Bloque de RECUADRO ESPECIAL -------------%
    \newenvironment{specialblock}[1]{%
      \begin{list}{}{%
          \setlength\leftmargin{1em}
          \setlength\rightmargin{1em}
      }
      \item[]
      \small
      \noindent\textbf{#1}\\[-0.3em]
      \rule{\linewidth}{0.4pt}\\[0.6em]
    }{%
      \end{list}
    }
      
%-------------------------- Portada -----------------------%
\newcommand{\oldtitle}[1]{\def\OldTitle{#1}}
\newcommand{\changerationale}[1]{\def\ChangeWhy{#1}}

\newcommand{\changebox}{%
\begin{tcolorbox}[title={Historial de cambios del tema}]
\textbf{Título anterior:} \OldTitle\par
\textbf{Motivo del cambio:} \ChangeWhy
\end{tcolorbox}
}

%-------------------------- Author Font -----------------------%
    \newcommand{\authorfont}[1]{{\fontfamily{EBGaramond-TLF}\selectfont #1}}

%-------------------------- Anexos -----------------------%
    \makeatletter
    \newcounter{anexo}
    \renewcommand{\theanexo}{\Roman{anexo}}
    \newcommand{\anexo}[1]{%
      \clearpage
      \phantomsection
      \refstepcounter{anexo}%
      \noindent\textbf{Anexo \theanexo.- }#1\par\medskip
      \addcontentsline{stc}{section}{Anexo \theanexo.- #1}%
    }
    \makeatother

    %-------------------------- Lista de figuras (personal) -----------------------%
\makeatletter
\newcommand{\MyListOfFigures}{%
  \clearpage
  \phantomsection
  {\centering\bfseries\Large Índice de figuras\par}%
  \medskip
  % Entrada en nuestro índice de contenido (stc)
  \addcontentsline{stc}{section}{Índice de figuras}%
  % Recuperación del listado (sin cabecera propia)
  \begingroup
    \parindent=0pt\parskip=2pt
    \@starttoc{lof}%
  \endgroup
}
\makeatother

%-------------------------- Bibliografía -----------------------%
\makeatletter
% Contador para las referencias
\newcounter{myref}
% Cabecera de la bibliografía, entra en el índice 'stc'
\newcommand{\MyBibliography}{%
  \clearpage
  \phantomsection
  {\centering\bfseries\Large Bibliografía y referencias\par}%
  \medskip
  \addcontentsline{stc}{section}{Bibliografía y referencias}%
  \setcounter{myref}{0}%
}
\newcommand{\MyBibItem}[4]{%
  \refstepcounter{myref}%
  \noindent[\themyref]\ #1.\ \emph{#2}. #3. #4\par
}
\makeatother

%--------- Establecer cuatro niveles de índice --------%
    \setcounter{tocdepth}{4}   
    \setcounter{secnumdepth}{4}% numera hasta \paragraph

%%%%%%%%%%%%%%%%%%%%%%%%%%%%%%%%

\makeatletter

    %--------- Interlineado Global --------%
    \edef\GlobalBaseLineStretch{\baselinestretch}
    
    %--------- Espaciado de seccinones --------%
    \renewcommand\section{\@startsection{section}
      {1}{\z@}
      {2.5ex \@plus 1ex \@minus .2ex} % espacio antes
      {1.5ex \@plus .2ex}             % espacio después
      {\normalfont\Large\bfseries}    % formato del título
    }
    \renewcommand\subsection{\@startsection{subsection}{2}{\z@}
      {3ex \@plus 0.8ex \@minus .2ex}
      {1.5ex \@plus .2ex}
      {\normalfont\large\bfseries}
    }
    \renewcommand\subsubsection{\@startsection{subsubsection}{3}{\z@}
      {3ex \@plus 0.5ex \@minus .2ex}
      {1ex \@plus .2ex}
      {\normalfont\normalsize\bfseries}
    }
    \renewcommand\paragraph{\@startsection{paragraph}{4}{\z@}%
    {-2ex \@plus -0.5ex \@minus -.2ex}% espacio antes
    {0.8ex \@plus .2ex}% espacio después
    {\normalfont\normalsize\itshape}}% ← cursiva en lugar de negrita

    %--------- Ecuaciones --------%   
    % Variables globales para almacenar títulos y notas
    \gdef\leq@current@section{}%
    \gdef\leq@current@subsection{}%
    \gdef\leq@last@printed@section{}%
    \gdef\leq@last@printed@subsection{}%
    
    % Comando para añadir nota (invisible en el cuerpo)
    \newcommand{\eqnote}[1]{%
      \addcontentsline{leq}{leqnote}{#1}%
    }
    
    % Comando para ecuaciones
    \newcommand{\eqblock}[2]{%
      % Verificar si necesitamos imprimir el título de sección
      \ifx\leq@current@section\leq@last@printed@section\else
        \addcontentsline{leq}{leqsection}{\leq@current@section}%
        \global\let\leq@last@printed@section\leq@current@section
        \gdef\leq@last@printed@subsection{}% Reset subsección al cambiar sección
      \fi
      % Verificar si necesitamos imprimir el título de subsección
      \ifx\leq@current@subsection\leq@last@printed@subsection\else
        \ifx\leq@current@subsection\@empty\else
          \addcontentsline{leq}{leqsubsection}{\leq@current@subsection}%
          \global\let\leq@last@printed@subsection\leq@current@subsection
        \fi
      \fi
      % Ecuación
      \refstepcounter{eqnum}%
      \begin{equation*}
        #1\tag{E.\theeqnum}
      \end{equation*}%
      \addcontentsline{leq}{leqt}{[E.\theeqnum] -- #2}%
      \addcontentsline{leq}{leqm}{#1}%
    }
    
    %%% ----- Índice de ecuaciones ----------%%%
    % Tamaño para []
    \newcommand{\leqIndexFont}{\normalsize}
    \newcommand{\leqTitleFont}{\tiny}
    \newcommand{\leqNoteFont}{\tiny}
    % Cabecera de sección 
    \newcommand*\l@leqsection[2]{%
      \addvspace{25pt}% Mayor separación antes de nueva sección
      \noindent\leqIndexFont
      \makebox[\linewidth]{\rule{.38\linewidth}{0.4pt}\ \ \textbf{  \MakeUppercase{#1}}\ \ \rule{.38\linewidth}{0.4pt}}%
      \par\vspace{-10pt}
    }
    % Cabecera de subsección 
    \newcommand*\l@leqsubsection[2]{%
      \par\addvspace{15pt}%
      \noindent\leqIndexFont #1
      \par\vspace{-7pt}
    }
    % Título de ecuación 
    \newcommand*\l@leqt[2]{%
    \par\addvspace{10pt}%
      \par\noindent\leqTitleFont #1\par
    }
    % Miniatura de ecuación
    \newcommand*\l@leqm[2]{%
      \vspace{0.3\baselineskip}%
      \par\scriptsize\noindent\hspace*{1.5em}\(#1\)\par%
    }
    % Nota de ecuación 
    \newcommand*\l@leqnote[2]{%
      \vspace{0.2\baselineskip}%
      \par\noindent\leqNoteFont\hspace*{1.5em}\textbf{Nota.--} #1\par\vspace{0.5\baselineskip}%
    }
    % Generación del índice
    \newcommand{\mylistofequations}{%
      \clearpage
      \phantomsection
      % Título visual de la lista (ajústalo a tu gusto):
      {\centering\bfseries\Large Índice de ecuaciones\par}
      % Entrada en nuestro índice 'stc':
      \addcontentsline{stc}{section}{Índice de ecuaciones}%
      % Llamada a tu lista real (usa el nombre exacto de tu macro):
      \@starttoc{leq}%
    }
    
    %--------- Captura de títulos de sección --------%
    \let\leq@original@section\section
    \let\leq@original@subsection\subsection
    % Flag para saber si estamos en una sección asterisco
    \newif\ifLeqStarSection
    \LeqStarSectionfalse
    
    % Redefinir \section CON CONTROL DE ASTERISCO
    \renewcommand{\section}{%
      \@ifstar{\leq@section@star}{\@dblarg\leq@section@nostar}%
    }
    \def\leq@section@star#1{%
      \global\LeqStarSectiontrue
      \gdef\leq@current@section{#1}%
      \gdef\leq@current@subsection{}%
      \leq@original@section*{#1}%
      \global\LeqStarSectionfalse
    }
    \def\leq@section@nostar[#1]#2{%
      \global\LeqStarSectionfalse
      \gdef\leq@current@section{#2}%
      \gdef\leq@current@subsection{}%
      \leq@original@section[#1]{#2}%
    }
    % Redefinir \subsection
    \renewcommand{\subsection}{%
      \@ifstar{\leq@subsection@star}{\@dblarg\leq@subsection@nostar}%
    }
    \def\leq@subsection@star#1{%
      \global\LeqStarSectiontrue
      \gdef\leq@current@subsection{#1}%
      \leq@original@subsection*{#1}%
      \global\LeqStarSectionfalse
    }
    \def\leq@subsection@nostar[#1]#2{%
      \global\LeqStarSectionfalse
      \gdef\leq@current@subsection{#2}%
      \leq@original@subsection[#1]{#2}%
    }

    % ----- Restauración de valores Globales de estilo ----------%
    % Flags
    \newif\ifInFootnote
    \InFootnotefalse
    \pretocmd{\@footnotetext}{\InFootnotetrue}{}{}
    \apptocmd{\@footnotetext}{\InFootnotefalse}{}{}
    % Profundidad de listas (soporta anidadas)
    \newcount\MyListDepth \MyListDepth=0
    \AtBeginEnvironment{la}{\advance\MyListDepth by 1}
    \AfterEndEnvironment{la}{\advance\MyListDepth by -1}
    \AtBeginEnvironment{lnum}{\advance\MyListDepth by 1}
    \AfterEndEnvironment{lnum}{\advance\MyListDepth by -1}
    \AtBeginEnvironment{itemize}{\advance\MyListDepth by 1}
    \AfterEndEnvironment{itemize}{\advance\MyListDepth by -1}
    \AtBeginEnvironment{enumerate}{\advance\MyListDepth by 1}
    \AfterEndEnvironment{enumerate}{\advance\MyListDepth by -1}
    % Restauración diferida: hazlo al INICIO del próximo párrafo real
    \newif\if@pendingRestore
    \@pendingRestorefalse
    \newcommand{\RequestRestore}{\global\@pendingRestoretrue}
    \everypar{%
      \if@pendingRestore
        \global\@pendingRestorefalse
        \ifInFootnote\else
          \ifnum\MyListDepth=0
            \setlength{\parskip}{\GlobalParskip}%
            \setstretch{\GlobalBaseLineStretch}%
          \fi
        \fi
      \fi
    }
    % Al final de bloques:
    \AfterEndEnvironment{equation}{\RequestRestore}
    \AfterEndEnvironment{equation*}{\RequestRestore}
    \AfterEndEnvironment{la}{\RequestRestore}
    \AfterEndEnvironment{lnum}{\RequestRestore}
    % Al final de macros:
    \apptocmd{\eqblock}{\RequestRestore}{}{}
    \apptocmd{\imagenfit}{\RequestRestore}{}{}
    
\makeatother

% ===================== ToC propio con 4 niveles (único bloque) =====================
\makeatletter

% --- Escritores: añadimos una línea al .stc en cada cabecera numerada ---
% Guardamos las versiones actuales para llamar al comportamiento normal
\let\MyTOC@orig@section\section
\let\MyTOC@orig@subsection\subsection
\let\MyTOC@orig@subsubsection\subsubsection
\let\MyTOC@orig@paragraph\paragraph

% SECTION
\renewcommand{\section}{%
  \@ifstar{\MyTOC@section@star}{\@dblarg\MyTOC@section@nostar}%
}
\def\MyTOC@section@star#1{%
  \MyTOC@orig@section*{#1}% sin entrada al .stc
}
\def\MyTOC@section@nostar[#1]#2{%
  \MyTOC@orig@section[{#1}]{#2}% comportamiento normal (escribe al .toc)
  % Escribimos también nuestra línea al .stc (texto con \numberline)
  \addcontentsline{stc}{section}{\protect\numberline{\thesection}#2}%
}

% SUBSECTION
\renewcommand{\subsection}{%
  \@ifstar{\MyTOC@subsection@star}{\@dblarg\MyTOC@subsection@nostar}%
}
\def\MyTOC@subsection@star#1{%
  \MyTOC@orig@subsection*{#1}%
}
\def\MyTOC@subsection@nostar[#1]#2{%
  \MyTOC@orig@subsection[{#1}]{#2}%
  \addcontentsline{stc}{subsection}{\protect\numberline{\thesubsection}#2}%
}

% SUBSUBSECTION
\renewcommand{\subsubsection}{%
  \@ifstar{\MyTOC@subsubsection@star}{\@dblarg\MyTOC@subsubsection@nostar}%
}
\def\MyTOC@subsubsection@star#1{%
  \MyTOC@orig@subsubsection*{#1}%
}
\def\MyTOC@subsubsection@nostar[#1]#2{%
  \MyTOC@orig@subsubsection[{#1}]{#2}%
  \addcontentsline{stc}{subsubsection}{\protect\numberline{\thesubsubsection}#2}%
}

% PARAGRAPH
\renewcommand{\paragraph}{%
  \@ifstar{\MyTOC@paragraph@star}{\@dblarg\MyTOC@paragraph@nostar}%
}
\def\MyTOC@paragraph@star#1{%
  \MyTOC@orig@paragraph*{#1}%
}
\def\MyTOC@paragraph@nostar[#1]#2{%
  \MyTOC@orig@paragraph[{#1}]{#2}%
  % Si no numeras los \paragraph, cambia \theparagraph por texto plano sin \numberline
  \addcontentsline{stc}{paragraph}{\protect\numberline{\theparagraph}#2}%
}

% --- Impresor del índice propio (lee el .stc) ---
\newcommand{\MyTOC}{%
  \par\bigskip
  {\centering\bfseries\Large Índice de contenido\par}%
  \medskip
  \begingroup
    \parindent=0pt\parskip=2pt
    \@starttoc{stc}%
  \endgroup
}

\makeatother
% ===================== fin ToC propio =====================


%%%%%%%%%%%%%%


% --- Datos del tema ---
\title{Tema 3.A.2 -- Los economistas clásicos y MARX\\
\large }
\author{Marcelo Ortega}
\date{Noviembre 2025}
\newcommand{\TemaCode}{3A02}

\oldtitle{Tema 3.A.2. Los economistas clásicos y MARX}
\changerationale{Sin cambio en el título. Este tema podría estructurarse por autores o, alternativamente, de manera temática (se presenta a continuación esta segunda aproximación).}

\begin{document}


\maketitle
\changebox

\newpage

% --- Página de resumen y modificaciones ---
\puntosclave{ %% Escribir puntos clave Apuntes CECO
    }
\vspace{1cm}

\modificaciones{
    Falta incluir el pensamiento económico de MALTHUS

    Falta incluir los antecedentes históricos, en particular: fisiócratas (QUESNAY) y mercantilistas (UK)
    }

\newpage
\MyTOC
\newpage

\mylistofequations
\clearpage

\MyListOfFigures
\clearpage


\pagenumbering{arabic}
\setcounter{page}{1}


\section{Introducción}\label{intro}
En primer lugar, podríamos preguntarnos porque es importante estudiar \textbf{X} desde un punto de vista histórico. 

Para responder a esta pregunta puede resultar muy ilustrativo el Premio Nobel de Economía de 2024, en particular, el concepto de \textit{destrucción creativa}. Un autor, o autores, puede encuadrarse en una época, resultar ser el contorno entre épocas o precisamente las ideas germinales de cambio que motivan el nacimiento de otra época, pero lo que es innegable es que ese autor es fruto del conocimiento previo que la ciencia ha alcanzado.

\imagenfit{DestruccionCreativa.png}{Ilustración: Destrucción creativa}{\href{https://www.nobelprize.org/prizes/economic-sciences/2025/press-release/}{Nobel Prize Outreach 2025. Nov 2025.}}

Si visualizamos el conocemiento económico como una escalera que se ha ido desarrollando a lo largo de su breve historia, es imprescindible conocer en que momento (peldaño) nos encontramos, es decir, cuál es la frontera del conocimiento, para poder avanzar, pero no por ello es menos importante entender y conocer el camino que nos ha traído hasta aquí, puesto que este conocimiento previo es condición necesaria para la innovación.

\imagenfit{PrescriptiveKnowledgeVSPropositionalKnowledge.png}{Ilustración: Conocimiento Prescriptivo y Conocimiento Proposicional}{\href{https://www.nobelprize.org/prizes/economic-sciences/2025/press-release/}{Nobel Prize Outreach 2025. Nov 2025.}}

\subsection{Contextualización}\label{context}
Estudiar la historia del pensamiento económico se revela una tarea clave para conocer el estado actual de la ciencia económica. Podemos definir la historia del pensamiento económico como la historia de las formas de adquirir conocimientos y razonar que usaron los hombres para comprender y explicar la naturaleza de los problemas económicos que se les presentaban.

En esta exposición, nos centraremos en la escuela clásica (1776-1850), responsable de que la economía se convierta en una disciplina científica e independiente. ¿Por qué consideramos que la ciencia económica nace a mediados del siglo XVIII? ¿No había pensamiento económico anterior?\footnote{Aunque la actividad económica ha sido una característica de la cultura humana desde los albores de la civilización, apenas se realizaron análisis formales de esa actividad hasta el desarrollo del capitalismo mercantilista en Europa occidental durante el siglo XV. En esa época, las principales sociedades europeas agrarias comenzaron a comerciar cada vez más entre sí, creando las condiciones necesarias para el nacimiento de la economía como estudio social. Los estudios económicos de esta época no eran sistemáticos: la teoría económica se desarrolló poco a poco a partir de las respuestas intelectuales individuales a problemas contemporáneos. No surgió ningún gran sistema analítico. No fue hasta mediados del siglo XVIII, con la aparición de la \textit{economía clásica} de ADAM SMITH, cuando la economía dio un importante paso para convertirse en una verdadera ciencia social.} De hecho, sí que había pensamiento económico anterior, pero los escritos económicos previos sufren en general 3 defectos:
\begin{lnum}
    \item Raramente son textos exclusivamente económicos. Su reflexión sobre la riqueza y el aumento de la producción viene de completar consideraciones más generales sobre la naturaleza del Estado, la justicia, la moral o se funda sobre los preceptos religiosos;
    \item Son más normativos que analíticos. Los autores no buscan comprender y deducir, sino prescribir. Estos textos comportan recomendaciones, a menudo muy precisas, de que debe hacer el Estado. En otros casos se centrar en descripciones de la sociedad ideal; y,
    \item Habitualmente son incoherentes y se centran en aspectos relativamente poco importantes de la economía.
\end{lnum}

En resumen, antes del siglo XVIII el pensamiento económico existe, pero no la ciencia económica, que comienza en este siglo con la primera formalización sistemática del pensamiento económico\footnote{Para un desarrollo histórico del periodo preclásico [Ver Anexo \ref{anx:preclasicos}]}.

\subsubsection{Relevancia}\label{importance}
SCHUMPETER (1954) aduce 4 motivos por los que considera fundamental el estudio de la historia del pensamiento económico:
\begin{lnum}
    \item \textit{Ventajas pedagógicas.} Resulta muy instructivo no contentarse con las enseñanzas del último tratado científico sobre una materia, puesto que éste se encuentra condicionado por el autor (por su método, su personalidad, su ideología y sus preferencias), especialmente en las denominadas ciencias sociales.
    \item \textit{Nuevas ideas.} Este estudio constituye una fuente de inspiración al encontrarse antiguos principios que en su momento no fructificaron y que en el presente pueden orientar nuevas investigaciones y avances en el conocimiento científico.
    \item \textit{Comprensión del proceder del espíritu humano.} Gracias al estudio de la historia del pensamiento económico, se entiende mejor el marco institucional y el espíritu de los tiempos que explica el por qué se producen unos determinados logros en unas épocas y no en otras.
    \item \textit{Desplazamiento progresivo de las fronteras de la ciencia económica.}\footnote{La evolución del pensamiento económico es un proceso complejo que resulta de la interacción de varios factores. Entre ellos, destacan el contexto económico del momento, el pensamiento económico precedente y los avances en otras disciplinas como la filosofía, las matemáticas, la física o la biología. El conocimiento de la historia del pensamiento económico permite entender las raíces intelectuales del pensamiento actual, aproximarse al análisis de los fenómenos económicos pasados y actuales, y valorar la importancia de los diferentes programas de investigación.} En cada época aparecen hechos y problemas distintos que requieren sus propias soluciones y su peculiar modo de afrontarlos, y por lo tanto, el estudio de la historia del pensamiento económico nos permite entender cómo ha evolucionado el objeto y el método de la ciencia económica.
\end{lnum}

\subsubsection{Problemática}\label{problem}
A lo largo de esta exposición intentaremos por tanto dar respuesta a las siguientes preguntas:
\begin{lnum}
    \item ¿Quiénes fueron los economistas clásicos?
    \item ¿Cuáles fueron sus aportaciones?
    \item ¿Qué trayectoria intelectual siguieron?
    \item ¿A quiénes influenciaron?
\end{lnum}

\subsection{Estructura de la exposición}\label{structure}


\clearpage
\section{Caracterización de la escuela}
La teoría económica clásica se refiere a una escuela de pensamiento económico cuyos principales exponentes son ADAM SMITH, JEAN-BAPTISTE SAY y DAVID RICARDO.\footnote{El término \textit{economía clásica} fue acuñado por MARX para referirse a la economía ricardiana, pero su uso se generalizó para describir también tanto a los seguidores de RICARDO y MILL como a todos los influidos por las percepciones generales de esos autores, incluido el propio marxismo.} Es considerada por muchos como la primera escuela económica moderna. Incluye también a autores como KARL MARX, THOMAS MALTHUS y FRÉDÉRIC BASTIAT. Habitualmente se considera que el último clásico fue JOHN STUART MILL.

Por todo ello, podemos definir el marco histórico en el que se desarrolla desde la segunda mitad del siglo XVIII hasta la segunda mitad del siglo XIX, con la muerte de JHON STUART MILL.

\imagenfit{escuelaclasica_autores.png}{Economía clásica, MALTHUS y MARX}{Original: LANDRETH y COLANDER (2010). \textit{Historia del pensamiento económico}. Extraído de: Apuntes Víctor Gutiérrez, Apuntes TCEE Tema 3.A.2}

De esta forma, y siguiendo la definición de SCHUMPETER\footnote{SCHUMPETER consideró que una escuela de pensamiento económico se podía definir por tres características: por tener un líder, unos seguidores y un cuerpo doctrinal propio}, la escuela clásica (1776-1850) cumple con las tres características para definirse como una escuela de pensamiento:
\begin{lnum}
    \item Podemos considerar a ADAM SMITH como el líder, en el sentido de que los temas que aborda son posteriormente tratados por otros autores clásicos;
    \item La escuela clásica cuenta con numerosos seguidores importantes como DAVID RICARDO, JEAN-BAPTISTE SAY o JOHN STUART MILL. De hecho, como veremos, se aprecia una evolución lineal de las contribuciones (en algunos casos), en el sentido de que JOHN STUART MILL busca matizar las conclusiones de RICARDO, que a su vez buscaba solucionar las fallas en los razonamientos de SMITH; y, finalmente,
    \item Se puede considerar que la escuela tiene un cuerpo doctrinal propio, que se podría sintetizar en un objeto y un método común o parecido.
    \begin{la}
        \item En relación al objeto, los economistas clásicos se centran en encontrar unos principios unificados para explicar algunos fenómenos como la causa de la riqueza de las naciones, la distribución de la renta, la teoría del valor o el comercio internacional. Además, hacen mucho énfasis en las relaciones de producción, en las clases sociales y contexto social, de modo que los sujetos objeto de estudio son países y clases sociales y no tanto el individuo.
        \item En relación al método, estos autores adoptarán una aproximación de economía política y realizan un tratamiento más sistematizado de los fenómenos económicos. Siguen principalmente un método deductivo, habitualmente mediante obras ensayísticas en las que abunda el razonamiento verbal.
    \end{la}
\end{lnum}

\subsection{Contexto histórico y autores}
En el año 1776, HUME murió, mientras que JACQUES TURGOT y el MARQUÉS DE CONDORCET abandonaron sus puestos en el gobierno. En ese mismo año, la revolución intelectual a la que habían contribuido, la Ilustración, comenzó a dar sus frutos. Fue un año de grandes tratados:
\begin{lnum}
    \item ADAM SMITH publicó \textit{An Inquiry into the Nature and Causes of the Wealth of Nations};
    \item ABBE DE CONDILLAC publicó su \textit{Commerce et le Gouvernement};
    \item JEREMY BENTHAM sus \textit{Fragments on Government}; y
    \item TOM PAINE su \textit{Common Sense}.
\end{lnum}

El 5 de mayo 1789,\footnote{Poco tiempo antes antes, KANT publica su \textit{opus magna}, \textit{Crítica de la razón pura} (“Kritik der reinen Vernunft”), 1781} los Estados Generales fueron llamados a sesión (por primera vez en casi dos siglos), estando el Estado en Bancarrota. Poco después, el 17 de junio, los miembros del Tercer Estado (la clase burguesa incipiente) forman la Asamblea Nacional, y tras varios meses de negociaciones fallidas la toma de la Bastilla da el pistoletazo de salida a la I Revolución Francesa. Poco después se aprobaría la Declaración de los Derechos del Hombre y del Ciudadano y se pone fin al regimen feudal característico, sistema de privilegios aristocráticos y clericales. Tras un interludio inicial, en 1792 se inician las Guerras Revolucionarias francesas, se abole definitivamente el régimen monárquico y en la Plaza de la Concordia, en 1973, el rey LOUIS XVI es decapitado por orden del Comité de Seguridad General (en manos de ROBESPIERRE). 

En 10 años más tarde, si alguien en Europa se había perdido estos eventos, pronto serían informados por las Guerras Napoleónicas. A medida que Europa descendía a un largo baño de sangre, las delicadas y optimistas tensiones de la Ilustración comenzaron a desvanecerse y la furia audaz y descarada de la era romántica retumbó con fuerza.

Mientras tanto, otra tormenta se estaba gestando bajo los pies, mucho más fantástica y destructiva del antiguo orden social que incluso los días más atareados de la guillotina o las bayonetas de Bonaparte. Llegó sin previo aviso, a espaldas de vagabundos que venían de la campiña inglesa. Algunos aficionados habían ideado una forma mejor de construir una rueca, y la civilización nunca volvió a ser la misma después de eso. La Revolución Industrial había comenzado.

Escritos durante la Ilustración, los mensajes centrales de \textit{La riqueza de las naciones} de ADAM SMITH (1776) —laissez-faire\footnote{La frase \textit{laissez faire}, \textit{laissez passer} es una expresión francesa que significa \textit{dejen hacer, dejen pasar} refiriéndose a una completa libertad en la economía: libre mercado, libre manufactura, bajos o nulos impuestos, libre mercado laboral y mínima intervención de los gobiernos. Fue usada por primera vez por VINCENT DE GOURNAY, fisiócrata del siglo XVIII, contra el intervencionismo del gobierno en la economía. De forma completa, la frase es: \textit{Laissez faire et laissez passer, le monde va de lui même}; \textit{Dejen hacer y dejen pasar, el mundo va solo}.}, las virtudes de la especialización, el libre comercio y la competencia, etc.— parecían una melodía apropiada para una era más apacible. No anticipó la magnitud de los trastornos económicos y sociales que la voraz era industrial estaba a punto de desencadenar. Además, era un libro confuso e inconsistente. Se necesitaba una nueva edición, corregida, ampliada y actualizada.

En la lucha por la sucesión apostólica surgieron tres nombres como fuertes contendientes: JEAN-BAPTISTE SAY, THOMAS MALTHUS y DAVID RICARDO. Todos tenían diferentes visiones de la economía política después de SMITH.
\begin{lnum}
    \item SAY (1803) quería llevarlo de regreso a la tradición francesa de oferta y demanda;
    \item MALTHUS (1798, 1820) quería agregar un énfasis completamente nuevo, alejándose de las complejidades obsesivas del valor y hacia una perspectiva más macroeconómica (y dinámica); y,
    \item RICARDO (1817) quería volver a las ideas SMITH, pero esta vez hacerlo correctamente, corrigiendo algunas de las fallas de las teorías del escocés.
\end{lnum}

De los tres, DAVID RICARDO resultó ser el más exitoso e influyente. Con su tratado de 1817, RICARDO elevó la economía a un grado sin precedentes de sofisticación teórica. La teoría de RICARDO, la más clara y sistemáticamente formalizada de todas, se convirtió en el sistema clásico.
Los bloques de construcción del sistema clásico, sin embargo, fueron establecidos inicialmente por THOMAS MALTHUS, el primer economista en Inglaterra. Su Ensayo sobre población de 1798 fue el primer cañonazo de esta época pesimista. Sosteniendo el axioma de que el crecimiento natural de la población excede el crecimiento económico, la teoría de MALTHUS implicaba una terrible dinámica de salario-fertilidad, es decir, cualquier aumento en el ingreso per cápita solo inducirá un aumento en la población que lo diluirá por completo. Los intentos de mejorar la suerte de los pobres, argumentó MALTHUS, son en última instancia contraproducentes. Esta se convertiría en la \textit{ley de hierro de los salarios} de la escuela ricardiana.

En 1810, en el curso de un debate público sobre la reanudación de la convertibilidad del oro (la controversia \textit{bullionista}), MALTHUS entró en contacto con el corredor de bolsa de Londres, DAVID RICARDO y entabló amistad (y competencia) con él. Fue durante el curso de este debate que se desarrolló la perspectiva clásica sobre el dinero. Invirtiendo la teoría cuantitativa, argumentaron que el valor del dinero estaba determinado por los costes inherentes de producir oro y plata, y que la cantidad de dinero se ajustaba para igualarla.

\subsection{Método: la Economía Política}


\clearpage
\section{Principales lineas de investigación}


\subsection{Teoría del Valor}
Una teoría del valor es cualquier teoría económica que intenta explicar el Valor implícito y explícito de los bienes.

\subsubsection{ADAM SMITH}
Para formular su teoría del valor, SMITH establece una distinción lógica previa y universal, entre dos tipos de valores:\footnote{Para ilustrar esto, plantea la siguiente paradoja:

\begin{adjustwidth}{2cm}{0pt}
\raggedright
\textit{Hay que destacar que la palabra VALOR tiene dos significados distintos. A veces expresa la utilidad de algún objeto en particular, y a veces el poder de compra de otros bienes que confiere la propiedad de dicho objeto. Se puede llamar a lo primero valor de uso y a lo segundo valor de cambio. Las cosas que tienen un gran valor de uso con frecuencia poseen poco o ningún valor de cambio. No hay nada más útil que el agua, pero con ella casi no se puede comprar nada; casi nada se obtendrá a cambio de agua. Un diamante, por el contrario, apenas tiene valor de uso, pero a cambio de él se puede conseguir generalmente una gran cantidad de otros bienes.}\\
\raggedleft{\authorfont{ADAM SMITH}, “\textit{La riqueza de las naciones}”, 1776}
\end{adjustwidth}

Esta paradoja muestra la incoherencia entre el valor de uso y el valor de intercambio y no será resuelta hasta el análisis marginalista que distingue la utilidad total de la marginal.
}
\begin{la}
    \item \textit{Valor de intercambio} (precio). Permite ver cómo se intercambia en el mercado una mercancía respecto de otra; y,
    \item \textit{Valor de uso} (utilidad). Es el poder que tiene un bien para satisfacer una necesidad.
\end{la}

A partir de esta distinción lógica, optará por dos aproximaciones distintas para explicar el valor de intercambio (valor explícito, en el que centrará su desarrollo).\footnote{Sin embargo, en todo momento deberá atenderse a la interacción entre ambas ya que la Teoría del precio (valor de intercambio) de Smith no tiene ninguna relación con el trabajo pasado gastado en producir una mercancía, es decir, habla solamente del trabajo que puede ser \textit{dado} o \textit{rescatado} en el presente (es un precursos de lo que posteriormente será la condición necesaria del valor: valor de uso $\to$ valor de intercambio, ningún bien que no tenga un valor de uso tendrá un valor de intercambio). 

\begin{adjustwidth}{2cm}{0pt}
\raggedright
\textit{Si un látigo de hule no sirve para nada, entonces el artículo no tiene valor económico en el comercio o en el uso, independientemente de todo el trabajo invertido en su creación.}\\
\raggedleft{\authorfont{ADAM SMITH}, “\textit{La riqueza de las naciones}”, 1776}
\end{adjustwidth}

} De hecho, es por esta doble aproximación la razón por la cual la obra de SMITH deja un vacío muy importante a la hora de formalizar una Teoría del Valor (de cambio) clara y definitiva:
\begin{la}
    \item \textit{Teoría del Valor-Trabajo}. En una sociedad primitiva donde no existe ni propiedad privada de la tierra ni beneficios, el único factor productivo es el trabajo. De esta forma, el valor de cambio de un bien depende únicamente del trabajo incorporado en su producción. Éste debe entonces remunerarse, no sólo por el número de horas sino, teniendo en cuenta otros aspectos como el ingenio y la cualificación;\footnote{ADAM SMITH se remontará en primer lugar a un estado primitivo, en el que analizará las bases del intercambio de bienes (independientemente del marco regulador a través del cual se realice), es decir, de la economía; y, a partir de ahí, deducirá las relaciones o motivaciones subyacentes en el marco de economías más avanzadas.\\ Esto no es más que una consecuencia lógica de ser fruto de su época. El estudio de la economía política que encauza SMITH sigue la misma metodología utilizada pos sus contemporáneas para explicar algunos fenómenos observados en las diferentes disciplinas de las ciencias sociales. En particular, debemos tener en cuenta que en esa misma época (el periodo de la Ilustración), surgen eminentes figuras que tratan de identificar el origen o las bases de las relaciones interpersonales, la finalidad de las congregaciones sociales (en forma de Estado) y los principios básicos que deben regir y respetar estas \textit{comunidades de vida}. Para ello, es común en la época remontarse a una situación primitiva a partir de la cual podamos identificar los fenómenos subyacentes que dan lugar a la organización social. Por ejemplo, la obra del \textit{Leviatán} de HOBBES; el concepto de \textit{buen salvaje} de ROUSSEAU; el \textit{Contractualismo} de LOCKE, etc.} y,
    \item \textit{Teoría de los Costes de Producción}. Más adelante, en una sociedad avanzada, donde existe acumulación de capital y propiedad privada de la tierra, deben pagarse rentas a todos los factores (trabajo, capital y tierra)\footnote{En este caso, ADAM SMITH llama al factor trabajo \textit{trabajo ordenado}: valor del trabajo cuando existe acumulación de capital) que en general es mayor porque, para él, mayor acumulación de capital implica mayor rendimiento del trabajo.}. Por lo tanto, el precio de un bien no dependerá exclusivamente de las rentas pagadas al factor trabajo, sino de los pagos a todos los factores productivos.
\end{la} 

\paragraph{Teoría del Valor-Trabajo: Economía primitiva}
SMITH se considera el primer formalizador de la Teoría del Valor-Trabajo (TvT):

\begin{adjustwidth}{2cm}{0pt}
\raggedright
\textit{El precio real de todas las cosas, lo que cada cosa cuesta realmente a la persona que desea adquirirla, es el esfuerzo y la fatiga que su adquisición supone. Lo que cada cosa verdaderamente vale para el hombre que la ha adquirido y que pretende desprenderse de ella o cambiarla por otra cosa, es el esfuerzo y la fatiga que se puede ahorrar y que puede imponer sobre otras personas. Aquello que se compra con dinero o con bienes se compra con trabajo, tanto como lo que compramos con el esfuerzo de nuestro propio cuerpo}\\
\raggedleft{\authorfont{ADAM SMITH}, “\textit{La riqueza de las naciones}”, 1776}
\end{adjustwidth}

Por tanto, el trabajo es la medida real del valor de un bien y, dada una movilidad laboral perfecta, el intercambio de las mercancías estaría determinada por la dificultad relativa de producirlas.\footnote{Los bienes económicos podían aumentar de valor en el mercado, pero lo que siempre permanece invariable es el trabajo, o sea el desgaste de energías física e intelectual del trabajador para producirlos, siendo entonces el trabajo el patrón definitivo e invariable del valor.

Todo bien producido necesariamente contiene trabajo, este trabajo es la fuerza de los hombres que han interactuado en el proceso de producción de dicho bien, o sea que en todo bien se vende la fuerza de trabajo (de cada hombre que interactuó en el proceso de producción). Esto no significa que a nivel individual un bien valga más por tener más trabajo en él. Aunque no era el factor determinante de los precios, estos oscilaban hacia su precio de producción gracias al juego de la oferta y la demanda.

\begin{adjustwidth}{2cm}{0pt}
\raggedright
\textit{Todo trato es: dame esto que deseo y obtendrás esto otro que deseas tú; y de esta manera conseguimos mutuamente la mayor parte de los bienes que necesitamos. No es la benevolencia del carnicero, el cervecero, o el panadero lo que nos procura nuestra cena, sino el cuidado que ponen ellos en su propio beneficio.}\\
\raggedleft{\authorfont{ADAM SMITH}, “\textit{La riqueza de las naciones}”, 1776}
\end{adjustwidth}
} Sin embargo, SMITH no logra explicar con la Teoría del Valor-Trabajo, los conceptos de beneficio, renta, interés y salarios, ni, principalmente, las evidentes diferencias que hay entre el trabajo incorporado a un bien (lo que sería su valor implícito) y el precio por el que se vende (lo que sería su valor explícito). 

\paragraph{Teoría del Valor-Trabajo: Economía moderna}
Esto lo lleva a desarrollar una segunda, la Teoría de los costes de producción, donde el valor de cambio de un bien (precio) depende del gasto invertido en el mismo (es decir, la suma de beneficios, rentas y salarios).\footnote{SMITH distingue tres clases sociales: terratenientes, trabajadores y capitalistas; en contraposición a los fisiócratas que distinguían dos clases sociales: clase estéril y clase productiva.} Ahora bien, SMITH se da cuenta de que esta Teoría del Valor (costes de producción) adolece de una falla terrible, pues no explica el precio de los productos en el corto plazo.

Para solventarlo, SMITH hace una distinción entre precio natural y precio de mercado:
\begin{la}
    \item El valor natural, o precio natural, es el suficiente para pagar a los factores productivos; mientras que,
    \item El valor relativo, o precio de mercado, fluctúa en función de la oferta y la demanda.
\end{la}

La competencia (\textit{mano invisible}) será entonces el proceso dinámico que lleva a los precios de mercado hacia sus precios naturales. Podría decirse pues, que el precio natural es un precio central en torno al cual gravita constantemente el precio de toda mercancía.\footnote{Por tanto, suponiendo una movilidad laboral perfecta, SMITH tiene una noción clara del precio (valor de intercambio) de equilibrio a largo plazo como resultado de un proceso de arbitraje competitivo que produce una tasa de ganancia igualada. Así los precios no oscilan ahora según los valores de los bienes, sino según sus precios naturales.}

\begin{adjustwidth}{2cm}{0pt}
\raggedright
\textit{Cuando el precio de cualquier producto no es ni más ni menos que lo que es suficiente para pagar la renta de la tierra, los salarios de la mano de obra, y los beneficios del capital empleado en su preparación, elaboración y transporte al mercado, de acuerdo a sus tasas naturales, el producto se vende por lo que puede llamarse su precio natural}\\
\raggedleft{\authorfont{ADAM SMITH}, “\textit{La riqueza de las naciones}”, 1776}
\end{adjustwidth}

Debemos notar, aunque es evidente, que esta Teoría del Valor está orientada en la oferta/producción de los bienes y, por tanto, la demanda no ocupa ningún lugar en la determinación del valor, dando lugar a la Teoría de los Costes de Producción, \footnote{El debate entorno a esta Teoría es la semilla intelectual de lo que es ahora conocido como el problema de la empresa: si todos los precios se definen por sus costes de producción, cualquier empresario recuperaría sus costes estableciendo un precio como mínimo igual a éstos} 

\subsubsection{DAVID RICARDO}
RICARDO reconoce dos conceptos como determinantes en el establecimiento del valor de cualquier bien:\footnote{Obviaremos la formalización que RICARDO hace de la condición previa para el valor de intercambia formulada por SMITH, que se desarrolla como sigue:\\
\begin{adjustwidth}{2cm}{0pt}
\raggedright
\textit{Poseyendo utilidad, las cosas derivan su valor en cambio de dos causas: de su escasez y de la cantidad de trabajo necesaria para obtenerla. [...] Así, pues, al hablar de las cosas, de su valor en cambio y de las leyes que regulan y especias respectivos, nos referimos siempre a aquellas cuya cantidad puede ser aumentada por el esfuerzo de la industria humana y en cuya producción la competencia actúa sin restricciones.}\\
\raggedleft{\authorfont{DAVID RICARDO}, “\textit{Principios de economía política y tributación}”, 1817}
\end{adjustwidth}

Por tanto, RICARDO atiende al valor, pero, ante todo, al valor de cambio, formalizando la condición previa del intercambio: sin utilidad no puede existir valor de cambio. Una cosa que no sea útil no se intercambia, pero la utilidad no determina su valor. Por otro lado RICARDO desarrolla otro concepto como fundamento del valor: la escasez. Para desentrañar la relación que RICARDO presupone entre escasez y utilidad, que de manera complementaria conducen al valor de intercambio, veamos otro razonamiento en el que disocia utilidad y valor: \footnote{Nótese que, para explicar el valor hoy emplearíamos este mismo ejemplo justamente en el sentido opuesto al de Ricardo. En efecto, al haber más medias (y aun aumentándose su utilidad total) su utilidad marginal desciende y por eso baja el valor de las medias.}

\begin{adjustwidth}{2cm}{0pt}
\raggedright
\textit{Si con una mejor máquina yo puedo, con la misma cantidad de mano de obra, producir dos pares de medias en vez de uno, de ninguna manera menoscabo la utilidad de un par de medias, aunque disminuye su valor.}\\
\raggedleft{\authorfont{DAVID RICARDO}, “\textit{Principios de economía política y tributación}”, 1817}
\end{adjustwidth}
} 
\begin{lnum}
    \item \textit{La escasez}. Puesta en relación con la demanda, o la intensidad de los deseos de quienes pretenden obtener un bien.\footnote{Este elemento interviene cuando los bienes no se pueden reproducir por el trabajo (como las obras artísticas, monedas raras, incunables, u otros bienes), o estén monopolizados, la escasez es el elemento determinante del valor.}
    \item \textit{El trabajo}. Éste es el elemento determinante del valor si los bienes pueden ser reproducidos por el trabajo humano. La mayoría de los bienes se encuentran en esta circunstancia cuando hay libre competencia. En este caso, la regla es: el valor de cambio es proporcional al trabajo incorporado al bien en su producción.
\end{lnum}  

Estos actuarán en mayor o menor medida en la determinación de dicho valor implícito del bien, en función de sus características. Sin embargo, como es normal, RICARDO pondrá el foco en los bienes que vienen caracterizados por la segunda condición, pues los primeros (considera él), son una parte residual de la economía.

La Teoría del Valor de RICARDO es simple, el valor implícito de los bienes viene determinado por el Trabajo aportado. Sin embargo, nos encontramos en una estado de desequilibrio en el que los precios (explícitos) no están ajustados y el mecanismo nivelador que conduce los precios a su precio natural (es decir, el que viene determinado por el trabajo aportado) se efectúa mediante la nivelación de los rendimientos del capital obtenido en todos los sectores.\footnote{Algunos autores interpretan la Teoría del Valor-Trabajo de RICARDO en el sentido de ser, en realidad, una Teoría basada en el coste real de producción, esto es, incluyendo beneficio y trabajo pero sin renta de la tierra, haciendo hincapié en el factor trabajo.}

\paragraph{Dinámica de ajuste RICARDIANA}
La condición de proporcionalidad adolece de una falla de la que RICARDO es consciente: la heterogeneidad de los trabajos y el empleo del capital en proporciones y cualidades diferentes según el tipo de producción.\footnote{También encuentra otras dificultades causadas por la desproporción entre el capital fijo y el circulante empleados en los diferentes procesos de producción; ora por la distinta durabilidad del capital fijo empleado (periodo largo o corto de amortización), ora por la desigual velocidad de rotación del capital circulante.} Para intentar solucionarlo, recurre al artificio de la abstracción y explicando las dinámicas que caracterizan al beneficio/capital.

En condiciones ideales, si toda clase de capital y sus combinaciones dispares se encuentran igualmente estructuradas en todas las industrias: el valor de cambio es proporcional al trabajo.\footnote{Como veremos, esta misma linea será luego utilizada por MARX para formular su Ley del Valor, en la que intenta reconciliar su Teoría del Valor-Trabajo con los valores explícitos observados (precio, valor de intercambio)} Ahora bien, RICARDO reconoce que esto no se da así en la realidad por lo que enuncia su Teoría que pretende ilustrar el camino a esta condición natural. Parte de la siguiente premisa:

\begin{adjustwidth}{2cm}{0pt}
\raggedright
\textit{El valor de cambio no varía con el aumento o reducción de los salarios si las demás circunstancias relativas al capital permanecen inalteradas.}\\
Y,\\
\textit{El capital es aquella parte de la riqueza de un país que se emplea con vistas a una producción futura, y puede ser aumentado de la misma manera que la riqueza.} (Es decir, se acumula por su productividad futura, y sigue:)\textit{[N]o puede existir acumulación sin motivo. [N]adie acumula sino con el propósito de hacer productiva su inversión.}
\raggedleft{\authorfont{DAVID RICARDO}, “\textit{Principios de economía política y tributación}”, 1817}
\end{adjustwidth}

A partir de estos dos puntos, RICARDO supone que los capitalistas actuaban movidos por el afán de lucro (principio de racionalidad económica) y, por lo tanto, que sus fondos se trasladarían sin dificultades hacia la industria que más beneficio brindara. El movimiento del capital hacia los sectores más rentables acabaría por igualar las tasas de beneficios en todos ellos.\footnote{Ricardo únicamente contempla un caso de rigidez en el traslado de los capitales, y es el que se debe a circunstancias de ausencia de riesgo, facilidad de fabricación u otros motivos eminentemente subjetivos, que inducen al propietario del capital a conformarse con una menor tasa de beneficio: 

\begin{adjustwidth}{2cm}{0pt}
\raggedright
\textit{[U]n capitalista que procura empleo provechoso para sus fondos, tomará naturalmente en cuenta todas las ventajas [e inconvenientes] que caracterizan a una ocupación con respecto a otra. Por tanto estará dispuesto a sacrificar parte de su utilidad monetaria en consideración a la garantía, sencillez o cualquier otra ventaja, real o imaginaria, que una colocación puede  tener sobre otra}\\
\raggedleft{\authorfont{DAVID RICARDO}, “\textit{Principios de economía política y tributación}”, 1817}\footnote{Ligeramente modificado y comentado para adecuar la lectura}
\end{adjustwidth}} Es decir, si la tasa de beneficio supera a la del interés abundan las peticiones de préstamo para invertir, y se iniciará un proceso de acumulación de capital productivo:\footnote{No hay dudas para interpretar que Ricardo se expresa en términos marginales; que considera los rendimientos netos del capital; que éstos los compara con el precio de coste del capital; y que la tasa de beneficio del capital (o eficiencia marginal del capital) la compara con el tipo de interés vigente.}

\begin{adjustwidth}{2cm}{0pt}
\raggedright
\textit{[E]sta competencia (es) la que ajusta el valor de cambio de los bienes, pues después de pagar los salarios del trabajo necesario para su producción, y todos los demás gastos requeridos para que el capital empleado vuelva a su primitivo estado de eficiencia, el valor restante o superávit será, en cada industria, proporcional al valor del capital empleado}\\
\raggedleft{\authorfont{DAVID RICARDO}, “\textit{Principios de economía política y tributación}”, 1817}\footnote{Ligeramente modificado y comentado para adecuar la lectura}
\end{adjustwidth}

\begin{adjustwidth}{2cm}{0pt}
\raggedright
\textit{La principal causa que impide la acumulación del capital es el incremento permanente de los salarios}. (Porque) \textit{[É]stos suben por la escasez de los artículos básicos que necesitan los asalariados, debido a la dificultad creciente de proporcionar alimentos y artículos de primera necesidad al creciente número de trabajadores. [E]n ese caso las utilidades necesariamente tendrán que disminuir.} (Porque, como ya se ha dicho) \textit{[E]l valor total de los bienes se divide solamente en dos porciones: la una constituye el beneficio; la otra, la retribución de la mano de obra.}\\
\textit{[D]e no existir esta circunstancia y mientras haya demanda suficiente, que es prácticamente ilimitada dada la gran diversidad de los deseos humanos por todo tipo de bienes y ornatos de la vida, el empleo del capital no tendría límite: y así como no existe límite para el deseo de conveniencias, aparato, mobiliario, ornato en la construcción, vestido y equipaje (...) [N]o puede existir límite al capital que puede emplearse para proporcionar esas cosas, excepto aquél que restringe nuestra aptitud para mantener a los trabajadores que las producen. [Y] además, su abundante acumulación no originaría ningún descenso de los beneficios.}
\raggedleft{\authorfont{DAVID RICARDO}, “\textit{Principios de economía política y tributación}”, 1817}
\end{adjustwidth}






\begin{adjustwidth}{2cm}{0pt}
\raggedright
\textit{[M]e he esforzado por demostrar que el valor real de una mercancía está regulado, no por las ventajas accidentales de que pueden disfrutar algunos de sus productores, sino por las dificultades reales que encuentra el productor menos favorecido [es decir, depende del coste marginal de la empresa marginal...], no está regulado por la tasa a que el Banco presta, sea 5, 4 ó 3\%, sino por la tasa de ganancias que puede obtenerse con el empleo del capital, lo que es totalmente independiente de la cantidad o valor del dinero [...]. Las solicitudes de dinero hechas al Banco dependerán, pues, de la comparación entre la tasa de ganancias que pueda lograrse con el empleo de éste, y la tasa a que está dispuesto a prestarlo. Si carga menos del tipo de interés del mercado no hay suma de dinero que no pueda prestar}\\
\raggedleft{\authorfont{DAVID RICARDO}, “\textit{Principios de economía política y tributación}”, 1817}\footnote{Ligeramente modificado y comentado para adecuar la lectura}
\end{adjustwidth}


De su regla del valor de cambio se extrae, como consecuencia inmediata, un corolario:\footnote{Para entender esto hay que tener en cuenta que el valor de cambio es la relación por la que se intercambian dos bienes: $Q_A/Q_B$. Esta relación, según RICARDO, es la de los trabajos incorporados en los dos productos a intercambiar ($L_A/L_B = Q_A/Q_B$); esta relación permanece invariable aunque previamente cada una de las cantidades de los dos bienes o de los trabajos incorporados se haya relacionado con un tercero que se toma como unidad de referencia, por ejemplo, el valor monetario de cada trabajo medido por los salarios pagados ($L_A/L_B = Q_A/Q_B = N_A·W /N_B·W$). Mientras esa relación sea la del trabajo empleado en cada bien, no se verá modificada con la variación del nivel general de los salarios ($W$), porque éste rige en todos los sectores; de modo que, al pagarse el mismo nivel salarial tanto en la producción de un bien como en la del otro, afecta por igual  tanto al numerador como al denominador de la relación.}



Así, resulta que un aumento salarial eleva relativamente el Valor de Cambio de los bienes que incorporan proporcionalmente menos capital fijo o un capital fijo menos duradero; y a la inversa, reduce el Valor relativo de los bienes producidos con más capital fijo o con un capital fijo más duradero.\footnote{El propio Ricardo nos hace una observación, en relación a la crítica que inicialmente suscitó su teoría. Recalca que él siempre trata sobre valores relativos; no dice que esto valga “1.000 £ y lo otro 2.000 £, porque el primero necesitó una cantidad tal de mano de obra, que costaría 1.000 £, y el otro una cantidad por valor de 2.000 £. Afirmé tan sólo que su valor relativo sería de dos a uno, y que se cambiarían uno por otro en esas proporciones" (ibídem, p. 35). De forma similar a la modificación de la regla, el corolario también se ve modificado según sea la diferente durabilidad del capital y la mayor o menor rapidez con la que vuelva a quien lo utiliza (ritmo de amortización).\\
Según HAYEK esto es el efecto Ricardo, admitido en la actualidad (aunque ya nadie se crea su teoría del valor-trabajo) porque las empresas que tienen proporcionalmente más mano de obra (es decir, están menos capitalizadas) son las más afectadas por una alteración de los salarios.} 

Por tanto, y he aquí la síntesis de sus postulados: lo que define el precio natural de los bienes es su poder de comprar otras mercancías, que sólo puede medirse de manera relativa y es únicamente dependiente del Valor-Trabajo incorporado:
\begin{lnum}
    \item El precio natural remunera salarios y beneficios, pero su determinante es el trabajo incorporado.; y,
    \item Los ingresos se ajustan a ese valor subyacente; no lo crean, y las dinámicas de mercado llevan el precio hacia el nivel que cubre el capital adelantado (fruto del Valor-Trabajo pasado) + el beneficio general + los salarios naturales, es decir, hacia su precio natural.
\end{lnum}

Es por tanto, una medida establecida en términos relativos y, en particular, en términos de Productividad del Trabajo en el proceso productivo de una mercancia


Sin embargo, tampoco supo conciliar su Teoría del Valor con los precios del mercado (corto plazo). Por lo general, estos precios no eran proporcionales al trabajo incorporado en la producción del bien. Zanjó esta cuestión distinguiendo entre precio natural, según el trabajo, y precio de mercado, o precio real.\footnote{Este último puede variar temporalmente del precio natural con el que tiene tendencia a converger; pues en caso contrario, si la diferencia de precios fuera persistente determinaría una mayor inversión de capital en la industria cuyo precio de mercado fuera mayor que el natural y una retirada de capital de las industrias que tuvieran un precio natural mayor que el de mercado.}




\subsubsection{MALTHUS}


\subsubsection{J.S. MILL}
En sus \textit{Principios de Economía Política} (1848), MILL (1806–1873) presenta el concepto de valor como:
\begin{lnum}
    \item Existen dos tipos de valor, el valor de uso y el valor de cambio, pero estos son conmensurables. El valor de uso es lo que uno estaría dispuesto a pagar por algo, y el valor de cambio es el valor medio de mercado; el valor de uso puede ser menor pero nunca mayor que el valor de cambio. Pero,
    \begin{la}
        \item El valor de uso no es objeto de estudio de la economía política;
        \item El valor (de cambio) es un concepto relativo, no absoluto.
    \end{la}
    \item El valor (implícito) se distingue del precio (valor explícito) debido a la variabilidad del poder adquisitivo del dinero, y puede medirse frente a un promedio general del resto de mercancías en lugar de frente a una sola (es decir, el dinero);
    \item El valor fluctúa según la oferta y la demanda alrededor de un \textit{valor natural}.
\end{lnum}

El valor de cambio de los bienes económicos, por tanto, está asociado al concepto de escasez y es, por tanto, una condición objetiva. De ahí que MILL nos haga observar que:\footnote{Malthus (1820, p. 66) ya había dicho algo parecido: \textit{ningún comprador estará dispuesto a pagar un precio elevado si puede pagar uno bajo}.}

\begin{adjustwidth}{2cm}{0pt}
\raggedright
\textit{(...) Mientras la cantidad disponible de un (bien) natural sea prácticamente ilimitada, no puede tener ningún valor en el mercado, a menos que sea susceptible de monopolizarse por medios artificiales, pues nadie pagará nada por lo que puede obtener gratis}\\
\raggedleft{\authorfont{J.S. MILL}, “\textit{Principios de Economía Política con algunas de sus aplicaciones a la Filosofía Social}”, 1848}
\end{adjustwidth} 

Esto nos lleva a las siguientes consideraciones: 
\begin{lnum}
    \item La finalidad de la economía es conseguir bienes escasos, de un modo u otro monopolizados, que así adquieren valor de cambio. Este enfoque pone el énfasis en la producción, en la obtención de bienes escasos y con valor de cambio.\footnote{Deja de tener su principal fundamento en el consumo; es decir, en la satisfacción de las necesidades humanas que aportan una vida grata al hombre.}
    \item El fundamento último del valor se encuentra en la institución de la propiedad privada, ya que sólo a través de ésta es posible impedir a los individuos el acceso libre (o sea, gratis) a los bienes. La escasez no se contempla en su sentido absoluto, tan sólo en términos relativos provocados artificialmente, mediante la monopolización, esto es, la apropiación por unos y la exclusión de otros del uso de los bienes.
    \item El campo de la economía se reduce a lo mercantil. Únicamente se contemplan los bienes con valor de cambio en el mercado.
\end{lnum}

Por tanto, MILL restringe el campo de (estudio) la economía a lo meramente mercantil, pues tan sólo son bienes económicos los que tienen valor de cambio;\footnote{Por tanto, MILL y sus seguidores (que hasta hoy son la mayoría de los economistas) no contemplan dentro del campo de la economía aquellos bienes producidos para la donación o el autoconsumo, por más que colaboren al bienestar y al aumento del nivel de vida de las personas.\\Curiosamente, MILL critica al sistema mercantilista, por haber confundido riqueza con acumulación de dinero en forma de metales preciosos. Pero su propuesta desemboca en el más genuino mercantilismo, pero disfrazado con la contemplación de una riqueza equivalente a los bienes materiales susceptibles de acumulación, que adquieren valor de cambio en el mercado. En realidad, sustituye el dinero por el conjunto de todos los bienes susceptibles de ser intercambiados; es decir, contempla la economía real y no la monetaria. En el sistema económico concebido por MILL el dinero es perfectamente suprimible; éste es un intermediario neutral que se usa por comodidad, para ahorrar tiempo y molestias. Por eso dice:\\
\begin{adjustwidth}{2cm}{0pt}
\raggedright
\textit{Se es rico cuando se tiene una provisión de artículos útiles, o los medios para adquirirlos. Todo aquello que sirve para comprar, todo aquello por lo que se dé a cambio algo útil o agradable forma parte de la riqueza. Aquellas cosas por las que no puede obtenerse nada a cambio, por muy útiles o necesarias que sean, no son riqueza en el sentido en que se emplea este término en economía política}\\
\raggedleft{\authorfont{J.S. MILL}, “\textit{Principios de Economía Política con algunas de sus aplicaciones a la Filosofía Social}”, 1848}
\end{adjustwidth}} y reduce, en consecuncia, el concepto de valor (precio natural) a nada (sin llegar realmente a descartarlo).\footnote{Algo así como un \textit{precio correcto}, ya que si el precio difiere del valor intrínseco (precio natural), es porque alguien ha sido estafado o porque existe una distorsión temporal en el mercado que se evidencia en el precio observado (explícito). Este autor acabará aclarando que: “No queda nada en las leyes del valor que el escritor presente o futuro deba aclarar.”}

Por otro lado, las aportaciones de MILL son destacables en la manera en que concilia tanto el enfoque de la demanda como el de la oferta, distinguiendo entre tres tipos de bienes:\footnote{El análisis de MILL de la teoría del valor es muy relevante al servir de base para dos cuestiones importantes:
\begin{lnum}
    \item La crítica marginalista de que sólo en el caso del segundo tipo de bienes el precio viene determinado íntegramente por la producción;
    \item La abstracción pionera que concibe las curvas de oferta y de demanda como funciones del precio.
\end{lnum}}
\begin{la}
    \item Bienes de oferta inelástica (p.ej. obras de arte), en los que el precio viene determinado íntegramente por la demanda;
    \item Bienes de oferta totalmente elástica (p.ej. dice MILL, las manufacturas), en los que el precio viene determinado por la oferta;
    \item Resto de bienes (i.e. aquellos cuya producción puede aumentar, pero a un coste creciente), en los que el precio viene determinado por la oferta y por la demanda.
\end{la}

\subsection{Crecimiento y distribución}

\subsubsection{ADAM SMITH} 
El capital es para Adam Smith un factor productivo, ya que es necesario disponer de él antes de empezar la producción para mantener al productor, a los obreros y para la adquisición de materias primas y herramientas durante el largo periodo de tiempo que media entre el inicio de la producción y la venta de las mercancías.\footnote{SMITH considera dos tipos de capital: fijo y circulante. El primero no puede rendir ningún tipo de ingreso sin la intervención del capital circulante que es el que suministra las materias primas necesarias y el que permite la manutención de los trabajadores, los cuales manejan la maquinaria, las herramientas y las instalaciones que constituyen el capital fijo.}

\paragraph{Crecimiento}
El progreso económico se logra mediante la acumulación de capital, mediante el ahorro, y la división del trabajo; pero nuevos grados de división del trabajo sólo es posible alcanzarlos con mayor acumulación de capital y el capital aumentará más pronto si se emplea en aquel ramo que proporciona la renta más considerable.\footnote{SMITH no consideró posible un ahorro improductivo, o sea, un atesoramiento (como el de los avaros), por el cual se retiran fondos de la circulación. Antes bien, para él todo ahorro siempre se convierte inmediatamente en capital, aun por personas distintas de las que ahorraron.\\Como para SMITH todo lo ahorrado termina por consumirse (es decir, podemos considerar que el ahorro es igual a la inversión), algunos autores ven en este principio un esbozo de la Ley de Say.} Además, para el autor sólo tiene relevancia la producción material, que físicamente es apreciable. 

Respecto a la distribución, SMITH considera, tres tipos de retribuciones:\footnote{A este respecto dice: 
\begin{adjustwidth}{2cm}{0pt}
\raggedright
\textit{El precio total o el valor en cambio de aquel producto anual no puede por menos de resolverse necesariamente en esas tres partes, y distribuirse entre los habitantes del país, como salarios del trabajo, o como beneficios del capital, o como renta de la tierra}\\
\raggedleft{\authorfont{ADAM SMITH}, “\textit{La riqueza de las naciones}”, 1776}
\end{adjustwidth}}
\begin{lnum}
    \item Los salarios de los trabajadores;
    \item Los beneficios del capital; y,
    \item Las rentas de la tierra.
\end{lnum}

\paragraph{Salarios}
Los salarios son la remuneración del trabajo por cuenta ajena.\footnote{En un Estado primitivo, el producto del trabajo pertenecía al trabajador; pero desde la apropiación de las tierras y la acumulación del capital, el trabajador tiene que compartir lo que produce con el propietario y el patrono.}
\begin{la}
    \item A corto plazo, los salarios vienen determinados por el \textit{fondo de salarios}.\footnote{Esta doctrina supone que hay un fondo fijo de capital destinado a pagar salarios, que depende de la tasa de acumulación de capital de los empresarios. La idea principal es que los salarios son adelantos en base a futuros beneficios de los empresarios derivados de las ventas de sus productos. De este modo, si se divide el fondo de salarios entre el número de trabajadores se obtiene el salario.\\
    Además, el salario puede subir si aumenta la demanda de trabajo, pero también existe un límite al alza salarial; los salarios sólo pueden subir en proporción al incremento de los capitales destinados al pago de dichas remuneraciones.} El fondo de salarios es el sobrante de la renta y del beneficio por encima de las necesidades de los propietarios y de los patronos para mantenerse ellos y sus actividades; dicho fondo se divide entre la mayor o menor cifra de población obrera. El aumento del fondo no depende de la riqueza absoluta de la nación, sino de la tasa con que crezca; es decir, el aumento de los salarios dependerá del ritmo con el que esté prosperando la nación.
    \item A largo plazo, los salarios vienen determinados por el salario mínimo de subsistencia.\footnote{La formalización de este concepto se conoce como \textit{la ley de hierro de los salarios} y se debe al economista prusiano FERDINAND LASALLE.\\
    LASSALLE fue un socialista comprometido desde temprana edad. Durante las revoluciones alemanas de 1848, habló en mítines públicos a favor de la causa democrática revolucionaria e instó a los ciudadanos de Düsseldorf a prepararse para la resistencia armada antes de la violencia que se esperaba tras la decisión del gobierno prusiano de disolver el Asamblea Nacional.\\
    Personaje muy interesante, fundador del partido socialista, amigo de MARX y otros miembros del círculo de Londres pero estrecho colaborador (\textit{laxo senso}) de las ideas de BISMARCK pues defendía la necesidad del sufragio universal frente al censitario característico del Partido Liberal Prusiano. \href{https://es.wikipedia.org/wiki/Ferdinand_Lassalle}{Ferdinand LASALLE}. Murió en un duelo a la edad de 39 años.} Los patronos pueden asociarse, con más facilidad que los obreros, para pagar salarios bajos, pero éstos no pueden descender de un mínimo: el de la subsistencia. En la medida en que el salario sea más elevado que ese mínimo de subsistencia, aumentará la población hasta que el nivel de salario por trabajador sea igual a ese mínimo (mecanismos poblacionistas). 
\end{la}

\paragraph{Los beneficios del capital}
Para SMITH, el interés y el beneficio son dos retribuciones distintas que entran dentro del mismo concepto de remuneración del capital.\footnote{El primero es debido al préstamo del capital y el segundo se debe a la utilización del capital en usos productivos. Ambos tipos de retribución van acompasados: si los beneficios suben o bajan también lo hacen los intereses. Los beneficios varían por las mismas causas que suben o bajan los salarios: el aumento o la disminución de los capitales (es decir, del ritmo de enriquecimiento de la nación). Pero la variación de los beneficios es de sentido inverso a la de los salarios.} 

$\Longrightarrow$ Con el incremento de la prosperidad de la nación y la mayor seguridad en las colocaciones del capital, el tipo de interés tiende a bajar; también disminuyen las oportunidades de efectuar inversiones rentables por lo que la tasa de beneficios desciende.\footnote{SMITH predijo que la tasa de beneficios caería con el paso del tiempo. Creía que había un número limitado de oportunidades de inversión y por el funcionamiento de la competencia y la reducción paulatina de las oportunidades disponibles de inversión, habría un descenso de la tasa de beneficio.} Los beneficios, evidentemente, dependen de los precios, pero los precios, con la prosperidad y el desarrollo económico, tienden a bajar. Este es un motivo adicional para el descenso del tipo de beneficio. 

\paragraph{Las rentas}
Para SMITH, la renta es el precio que se paga por el uso de la tierra con independencia de los gastos realizados en mejoras; por tanto, se basa en un privilegio o don gratuito de la naturaleza. Considera que la renta es producto de una situación de monopolio ya que la cantidad de tierras cultivables es limitada y por ello el propietario puede exigir la mayor parte de su rendimiento a los agricultores, dejándoles sólo lo indispensable para subsistir.\footnote{Para Smith las diferencias existentes entre los tipos de rentas son debidas a las diversas fertilidades de las tierras, y a los distintos emplazamientos, aunque la mejora en los transportes (al abaratar los costes) iguala la diferencia debida a la localización de las tierras.} 

$\Longrightarrow$ Las altas rentas dependen de los altos precios de los productos agrícolas. \footnote{Esto estaría en la base de los problemas circulares que presenta la Teoría del Valor (explícito) para SMITH, a través de la cual los precios giran en torno al valor natural capaz de remunerar los factores.\\ Esto es así porque en la formación del precio natural intervienen las rentas; es decir, la relación causal es al revés: los altos precios dependen de las altas rentas.\\
Al parecer, Smith no vió contradicción en esas aseveraciones, posiblemente porque se refirieran a situaciones distintas. A pesar de la observación que a este respecto le hizo su amigo HUME, no modificó su versión, teniendo la oportunidad de hacerlo, en las sucesivas ediciones.}

\subsubsection{DAVID RICARDO}  
RICARDO acepta las tésis de SMITH respecto a la acumulación de capital como motor del crecimiento económica, sin embargo,este crecimiento está constantemente amenazado por la caída tendencial de la tasa de beneficio resultado del aumento de la renta y del uso de tierras de peor calidad. Como se verá, RICARDO plantea que la economía tiende a un estado estacionario que no podremos sobrepasar.


El proceso de industrialización y la maquinaria fue un asunto especial y favorablemente tratado por RICARDO. Consideraba que las máquinas aumentaban grandemente la producción abaratando los precios, lo cual contribuiría al aumento generalizado del nivel de vida de toda la población, incluidas las clases trabajadoras, puesto que podrían comprar más mercancías con los mismos salarios. Al principio, no se percató del efecto tan negativo que provocaba la expulsión de la mano de obra, al ser ésta sustituida por las máquinas. En la tercera edición de su libro (para tratar este asunto de la maquinaria) se reafirmó en la tesis anterior, aunque reconoció que la sustitución del trabajo humano por la maquinaria es, a menudo, muy perjudicial a los intereses de la clase trabajadora. Pero supuso que sería transitoria tal situación, pues al poco tiempo los capitalistas, como ahorran más (pues gastan menos que antes en comprar artículos más baratos y en pagar menos salarios) y como también desean emplear productivamente sus fondos, invertirían sus ganancias en la elaboración de otra mercancía útil a la sociedad y de la que no pudiera faltar demanda. Así, el paro sería reabsorbido pronto por la demanda de mano de obra en la fabricación de nuevas mercancías. Una condición para llegar al pleno empleo (en el que creía RICARDO) era que al introducir la maquinaria se llegaría a aumentar la producción bruta (se trata de la valoración del producto neto incluyendo el valor añadido por el trabajo; o sea, los salarios); si ésta disminuyera ciertamente se originaría paro; aunque hay que tener en cuenta que es la producción neta (que no contabiliza los salarios) la que en realidad beneficia a los empresarios y, siendo el incremento de ésta compatible con la disminución de la producción bruta, podría resultar que un aumento de la producción neta beneficiaria a los empresarios y a la vez perjudicara a los trabajadores. 

Por tanto, las quejas en este sentido de los obreros estaban justificadas por la  leyes de la economía política. Otra condición (también supuesta por Ricardo y que para él avalaba su idea de la escasa influencia de la introducción de la maquinaria en el paro, al que eufemísticamente RICARDO denominaba población redundante) era que la invención de máquinas más productivas se realizaría en un proceso lento, de modo que habría tiempo para la reabsorción del paro que causara la sustitución de la mano de obra por las nuevas máquinas. Teoría de la compensación es la denominación genérica que se dio a las explicaciones que, como ésta de Ricardo, giran en torno a la posterior mejoría de las clases obreras en compensación a los sacrificios que temporalmente debían padecer los trabajadores a causa del progreso técnico.





Por otro lado, la distribución del valor del producto nacional entre los factores de la producción fue un asunto de preocupación preferente para RICARDO. 

\paragraph{Salarios} 
El estilo de Ricardo, claro, conciso y preciso, no deja lugar para explicar sus ideas más que con sus propias palabras:

"La mano de obra al igual que las demás cosas que se compran y se venden, y que pueden aumentar o disminuir en cantidad, tiene su precio natural y su precio de mercado" 

"El precio natural de la mano de obra es el precio necesario que permite a los trabajadores, uno con otro, subsistir y perpetuar su raza, sin incremento ni disminución" 

Como se ve, se trata del coste de producción de los obreros: 

"La aptitud del trabajador para sostenerse a sí mismo y a su familia [...] no depende de la cantidad de dinero que pueda percibir por concepto de salarios, sino de la cantidad de alimentos, productos necesarios y comodidades que por costumbre disfruta, adquiriéndola con dinero" (ibídem, p. 71). 

Es decir, se trata del coste de producción real, no del monetario, pues depende del salario real y no del nominal; por consiguiente si el precio de los alimentos y productos necesarios (de primera necesidad) sube, el precio natural de la mano de obra aumentará. En caso contrario bajará. 

La opinión de Ricardo (ibídem, p. 71) es que: 

"con el progreso de la sociedad, el precio natural de la mano de obra tiende siempre a aumentar, porque uno de los principales bienes que regula su precio natural tiene tendencia a encarecer, debido a la mayor dificultad para producirlo" (se refiere a los productos agrícolas y a la ley de los rendimientos decrecientes). 

"Sin embargo, [...] las mejoras agrícolas [y] el descubrimiento de nuevos mercados, de los cuales pueden importarse las provisiones, vienen a contrarrestar, por un tiempo, la tendencia ascendente del precio de los productos de primera necesidad, y a ocasionar a veces una reducción de su precio natural, así también las mismas causas producirán los efectos correspondientes sobre el precio natural de la mano de obra". 

"El precio natural de todos los bienes, salvo el de los productos primos y el de la mano de obra, tiende a disminuir al progresar la riqueza de la población", 

porque el incremento del precio de las materias primas y de la mano de obra se compensan con creces con las mejoras en la maquinaria y el incremento de productividad de la mano de obra debido a la división del trabajo, a la mayor habilidad y mejor distribución de la mano de obra. 

"El precio de mercado de la mano de obra es el precio que realmente se paga por ella, debido al juego natural de [...] la oferta y la demanda; la mano de obra es costosa cuando escasea y barata cuando abunda" (ibídem, p. 71-72).

No obstante, tal precio tiende al natural y sólo difiere de él temporalmente. Si dicho precio de mercado sube, "la condición del trabajador es floreciente [...] y puede [...] criar una familia sana y numerosa"(ibídem, p.72). 

Mas el consecuente aumento de la población y de la mano de obra hace caer los salarios hasta su nivel natural; y a veces hasta por debajo: "Cuando el precio de mercado de la mano de obra es inferior a su precio natural, la condición de los trabajadores es de lo más mísera [...] Sólo después de que sus privaciones han reducido su número [...] tendrá el trabajador las comodidades moderadas que le proporcionará la tasa natural de salarios". 

Sin embargo, a pesar de la tendencia a confluir el salario natural y el de mercado, en una sociedad mejorada (por el progreso económico) el salario de mercado puede estar por encima durante un período indefinido, porque un incremento constante y gradual del capital puede estimular continuamente un incremento de la demanda de mano de obra y de la población. 

Esto es, para Ricardo no hay paro a la larga:

"El precio natural de la mano de obra, aún estimado en alimentos y productos necesarios [no es] absolutamente fijo y constante. En un mismo país varía en distintas épocas y difiere cuantiosamente de un país a otro. Depende esencialmente de los hábitos y de las costumbres de la gente [...]. Muchas de las comodidades de que actualmente se goza en una casita inglesa se habrían considerado un lujo en un periodo anterior de nuestra historia" (ibídem, p. 73-74). 

En resumen, "los salarios suben o bajan por dos causas (ibídem, p. 74):
\begin{lnum}
    \item La oferta y la demanda de mano de obra; y,
    \item El precio de los bienes en que el obrero gasta su salario.
\end{lnum}

\paragraph{Los beneficios del capital}
El beneficio depende de los ingresos y del coste de la mano de obra. Concretamente, el beneficio es la diferencia entre los ingresos y los pagos de los salarios.\footnote{Por consiguiente, los beneficios aumentan con los ingresos, que dependen del precio de mercado; también aumentan al disminuir los salarios nominales, que dependen del precio de los alimentos y de los artículos de primera necesidad utilizados por los trabajadores.}

Ricardo estudia la tendencia, con el transcurso del tiempo, tanto de la tasa del beneficio como de los beneficios en sí. Para ello atiende al otro componente de los beneficios, el ingreso. Éste aumenta con la subida de los precios de mercado. Cuando el precio de mercado de un bien sube y se mantiene alto, los empresarios tienden a trasvasar sus capitales poco a poco hacia la industria que, en proporción, logra unos beneficios más altos, retirándolos de las industrias que obtienen unos beneficios proporcionalmente más bajos. De este modo, la tendencia a largo plazo es a igualarse las tasas de beneficio del capital en todas las clases de industrias.\footnote{En realidad, habría que decir en todos los sectores, porque también hay trasvases de capitales entre el agro y la industria y viceversa; es más, el beneficio normal obtenido en las explotaciones agrarias marca el límite al que tiende el beneficio normal en cualquier otro sector. Mientras no se introduzca maquinaria nueva que incorpore avances tecnológicos ahorradores de mano de obra la igualdad de la tasa de beneficios implica que el valor de los bienes es proporcional al trabajo empleado en producirlos. Esto quiere decir que el precio, siempre relativo, de las manufacturas no varía con los cambios de los precios de las materias primas y de los salarios, en el supuesto, claro está, de la constancia del precio del dinero.}

Cuando existe esta flexibilidad (en la que RICARDO creía) para trasladar, con relativa rapidez, los capitales de un sector a otro hasta conseguir la igualación de las tasas de beneficios, el valor de los bienes, tanto de los alimentos como de las manufacturas, tiende a su precio natural y se determina por el trabajo incorporado al producirlo en la empresa marginal, la que no paga renta por producir en la situación más desfavorable (y en el límite, la que sólo puede obtener la tasa normal de beneficio: aquella por debajo de la cual el empresario no está dispuesto a arriesgar su capital). En estas condiciones, el aumento de la población, y de la consiguiente producción de alimentos, causa un crecimiento de los salarios nominales y de la renta de la tierra, en tanto que disminuyen los beneficios y se mantienen constantes los salarios reales. Por tanto, la tendencia es que los beneficios bajan a medida que los salarios nominales suban (aun cuando el salario real se mantenga constante). Veamos esto con la ayuda del cuadro siguiente elaborado con las explicaciones y las cifras de RICARDO:\footnote{Se supone, en este modelo de RICARDO, que los trabajadores necesitan realmente para su sostenimiento 3Kg de cereales y el equivalente a 12 £ en otros bienes de primera necesidad, cuyos precios no varían al no verse afectados por un cambio en la cantidad de mano de obra necesaria para producirlos. Supongamos inicialmente que no es necesario utilizar otra tierra que no sea de 1ª calidad. Esta unidad de terreno requiere para su cultivo una cuadrilla de 10 trabajadores y rinde 180 Kg de cereal. El precio del cereal en el mercado equivale al natural que en proporción a los 10 trabajadores determina un valor de 720 £, lo cual arroja un precio unitario de 4,00 £/Kg. Por lo tanto, en la situación 1, inicial, el valor del producto por unidad de tierra, capital y trabajo utilizado en ella es: 180 Kg x 4 £/Kg = 720 £. El salario nominal de cada trabajador es: 3 Kg x 4,00 £/Kg + 12 £ = 24 £ El importe total de los salarios pagados a los 10 trabajadores por unidad de terreno y capital es: 10 x 24 £ = 240 £. Por consiguiente, el beneficio (diferencia entre ingresos y salarios) es: 720 - 240 = 480 £ Esta tierra no paga renta, ya que es la marginal. Si aumenta la población y se precisa mayor producción de alimentos y ya no hay disponibles tierras de primera calidad, habrá que recurrir a la explotación de tierras menos fértiles cuyo rendimiento es menor. Se supone que con el mismo número de trabajadores por unidad de tierra y con la misma cantidad de capital se obtiene una producción menor, digamos 170 Kg. Evidentemente, como la cantidad de trabajo no ha variado, el valor del producto tampoco cambia; por lo tanto, será 720 £, pero como ahora se han obtenido 170 Kg, el precio unitario sí varía: 720 : 170 = 4,24 £/Kg. Manteniéndose el salario real por obrero, el nominal se convierte en: 3 Kg x 4,24 £/Kg + 12 £ = 24,72 £. Los diez trabajadores totalizarán un gasto salarial de 247 £ (aproximadamente), y, en consecuencia, el beneficio será de: 720 £ - 247 £ = 473 £, sin que esta tierra tenga que pagar renta. Pero en las tierras de primera calidad se sigue produciendo su rendimiento normal, que es de 180 Kg de cereal por cada 10 trabajadores y unidad de tierra y capital. Puesto que el precio de mercado ha subido, por la presión de la demanda, hasta cubrir la inversión de la empresa marginal, quedará estabilizado según el cálculo anterior en 4,24 £/Kg. A este precio los ingresos proporcionados por las tierras de 1ª calidad serán de: 180 x 4,24 = 763 £; en estas tierras se paga el mismo salario real y nominal que en las de 2ª calidad, cuyo importe total será de 10 x 24,72 = 247 £ (aproximadamente); el beneficio también tiene que ser el mismo que en las tierras de 2ª calidad, por la hipótesis de la igualación de las tasas de beneficio. En caso contrario, la propia competencia, aumentando la demanda por arrendar tierras de 1ª calidad, haría subir la renta hasta que a un empresario agrícola le resulta indiferente cultivar tierras de 1ª o de 2ª calidad, por obtener la misma tasa de beneficio; así es que los terratenientes han conseguido extraer una renta por la diferencia entre los ingresos y la suma de los salarios y el beneficio normal: 763 -(247+473) = 43 £. Obsérvese que esta renta (si se hubieran realizado los cálculos con toda exactitud en lugar del redondeo aproximado) debe ser igual al producto de la diferencia de rendimiento en la producción ($\Delta$Q = 10 Kg) por el nuevo precio de mercado (P = 4,24 £/Kg). O bien, tal como se dijo antes, en el caso desalarios, la renta es la diferencia de valores de la producción después y antes de subir el precio en las explotaciones intramarginales. En el cuadro puede seguirse fácilmente las nuevas distribuciones de las retribuciones a los factores de la producción, cuando otro aumento de la población provoca la necesidad de cultivar tierras de 3ª, e incluso 4ª calidad. En el cuadro se aprecia cómo las tierras de 1ª calidad aumentan su renta y también las de 2ª, e incluso las de 3ª, a medida que dejan de ser los cultivos marginales. Mientras tanto los beneficios se van reduciendo. El límite absoluto de la reducción llegaría hasta el beneficio cero; sin embargo, previsiblemente el empresario agrícola dejaría antes el cultivo de una tierra tan marginal que no proporcionara un tasa de beneficio similar a otro tipo de inversión segura (como podría ser prestar el capital a un banco). En el cuadro también se ve que el importe total de los salarios nominales va creciendo (a pesar de mantenerse constante el salario real); el límite absoluto de su crecimiento sería el valor total del ingreso en la explotación marginal, pero previsiblemente (por las razones anteriormente dadas al considerar el límite de los beneficios) nunca se llegaría a ese límite absoluto.}

\imagenfit{EvDistribucionesRICARDO.png}{Evolución de la distribución de las retribuciones}{ESCARTÍN GONZÁLEZ (2008), \textit{Apuntes sobre historia del pensamiento económico.} Capítulo 16}

Como se ve, debido al incremento de la población, lo normal es que el precio de los alimentos suba (por necesitarse más mano de obra para obtener una igual cantidad adicional de subsistencias); pero este hecho no afecta al precio de los bienes manufacturados, mientras no requieran más cantidad de trabajo en su producción. Entonces, si los salarios continuasen iguales, los beneficios de los fabricantes permanecerían iguales; pero si, como con toda seguridad acontece, los salarios aumentasen a causa del alza del precio de los alimentos, en ese caso sus utilidades necesariamente tendrán que disminuir.\footnote{Tal idea es expresada con toda rotundidad por Ricardo: los beneficios disminuyen siempre que aumentan los salarios. Fue, por tanto, el responsable de esta falsa idea que los patronos, todavía en nuestros días, se creen irreflexivamente: los beneficios disminuyen siempre que aumentan los salarios. De esta idea, DE QUINCEY dijo socarronamente que los beneficios eran las migajas dejadas por los salarios. La falsedad de esta idea estriba en que pueden lograrse aumentos de los salarios y también de los beneficios (o utilidades). Se dolió Ricardo al comprobar que Malthus mal interpretaba su idea cuando decía \textit{que las utilidades y los salarios podían subir al mismo tiempo} y luego cuando Malthus escribía que era \textit{una perogrullada decir que si el valor de las mercancías se reparte entre el trabajo y las utilidades, cuanto mayor es la parte de uno, menor es la que queda para las otras}. A esto Ricardo replicó que él hablaba de proporciones sobre el producto total obtenido y que si la proporción de ese producto que se llevan los salarios aumenta necesariamente tendrá que disminuir la que reste para beneficios.\\
Efectivamente, si su cuadro de la distribución, anteriormente expuesto, lo modificamos y calculamos los porcentajes del valor del producto que van a salarios y a beneficios veremos que mientras sube el de los salarios va bajando el de los beneficios.} El progreso (no técnico) a largo plazo conduce, entonces, a los siguientes resultados:
\begin{lnum}
    \item La renta de la tierra va creciendo;
    \item Los salarios nominales también experimentan un alza, pero los salarios reales permanecerán en su nivel de subsistencia (según las condiciones sociales de cada país y época histórica); y,
    \item Los beneficios van disminuyendo.
\end{lnum}

Pero téngase en cuenta tres cosas sobre el modelo de Ricardo: 
\begin{lnum}
    \item Hay supuestos implícitos sin los cuales su teoría queda invalidada, a saber: que existe un incremento del precio del producto por haber aumentado la demanda; que dicho precio se regula por la empresa marginal operando a corto plazo; y que rige la ley de los rendimientos decrecientes;
    \item Pese a lo que luego diga Ricardo ad hoc, el trato que le da es para determinar los valores nominales y no los valores proporcionales; y,
    \item La relación de cambio que Ricardo tanto defiende (intercambios proporcionales al trabajo incorporado) no rige para un mismo artículo; porque si su regla se aplicara siempre, el cereal obtenido en la tierra menos fértil se debería intercambiar a la par con el obtenido en la tierra de primera calidad, pues en ambas se utiliza la misma cantidad de trabajo y capital; es decir, el producto de 10 hombres de la primera tierra (180 Kg) equivalen al producto de 10 hombres de la segunda tierra (170 Kg) y en consecuencia, 180Kg de cereal se deberían intercambiar por 170 Kg del mismo cereal.
\end{lnum} 

\paragraph{Las rentas}
La renta se paga porque la tierra: tiene dueños, su disponibilidad es limitada y a las diferencias de fertilidad y emplazamiento. El pago de la renta no es un incentivo para atraer sus servicios ni entra en el precio del producto, como ocurre con el trabajo y el capital. 

Según la definición de Ricardo (ibídem, p. 51) "es la parte del producto de la tierra que se paga al terrateniente por el uso de las energías originarias e indestructibles del suelo". 

La explicación de la naturaleza de la renta de la tierra dada por Ricardo ha venido a denominarse la teoría de la renta diferencial: el origen de la renta se encuentra en la diferente cantidad de producto obtenido mediante el empleo de dos cantidades iguales de capital y trabajo. 
Estas diferencias en la cantidad de producto se deben a dos motivos: 
\begin{lnum}
    \item Diferencias de fertilidad o emplazamiento de las tierras (teoría del margen extensivo);
    \item Diferencias de rendimiento en el mismo terreno ante sucesivos aumentos de capital y trabajo (teoría del margen intensivo).
\end{lnum}

En cualquiera de los dos casos el principio es el mismo; la renta proviene invariablemente del empleo de una cantidad adicional de trabajo con un ingreso proporcionalmente menor. La renta se basa en que: 

el valor de cambio de todos los bienes [...] está siempre regulado [...] por la mayor cantidad de trabajo necesariamente gastada en su producción, por quienes no disponen [de] circunstancias ampliamente favorables [o] por el capital que sigue produciendo esos bienes en las circunstancias más desfavorables.

Por estas apreciaciones, se nota que RICARDO usa el concepto de marginalidad, aunque luego no sepa aplicarlo como instrumento de análisis general. Es obvio que el aumento de la producción en esas circunstancias más desfavorables se debe al incremento de la demanda y del precio.

Para atenderla, se aumenta la producción, bien poniendo en cultivo nuevas tierras de menor grado de fertilidad (o más alejadas) con el consiguiente aumento de los costes, o intensificando la producción en las mismas tierras con menores rendimientos. La presión de la demanda más alta hace subir el precio hasta que se cubran los costes medios más elevados de los productores marginales, porque el sostenimiento de la producción en condiciones de extramarginalidad es transitorio: o cierra la empresa, reduciéndose la producción global, o suben los precios y la empresa se mantiene en la marginalidad. 

Los productores intramarginales se ven obligados a pagar la renta a cargo de sus excedentes so pena de perder los contratos a su finalización, pues no se les renovaría y se concertarían con quienes sí estuvieran dispuestos a pagarla. La renta en la empresa intramarginal es así la diferencia entre los valores de la producción después y antes de haberse incrementado el precio a causa del aumento de la demanda. 

Pongamos un ejemplo para aclarar el concepto: En una tierra A se recolectan 1.000 unidades de trigo que se venden por un valor de 1.800 £, lo que supone un precio unitario de 1,8 £/u. Luego, por aumentar la demanda, se necesita poner en cultivo la tierra B que rinde 900 unidades de trigo empleando la misma cantidad de mano de obra y capital que en A; entonces, por la teoría del Valor-Trabajo, el valor de lo producido en B tiene que ser igual al de antes en A: 1.800 £, lo que equivale a 2 £/u. Este nuevo precio de mercado, debido al incremento de la demanda, también rige para el trigo de la tierra A, cuyo valor sube ahora a 2.000 £. Por consiguiente la diferencia entre este nuevo valor y el antiguo (2.000 £ – 1.800 £ = 200 £) es la renta de la tierra A. Por lo tanto, al aumentar la producción los precios suben porque se emplea más trabajo en la producción de la última porción obtenida, y no porque se pague una renta al terrateniente. Es decir, el cereal no se encarece porque hay que pagar una renta, sino que debe pagarse una renta porque el cereal es caro; y no acaecería reducción alguna en el precio del cereal aunque los terratenientes condonasen la totalidad de las rentas. Ahora bien, los avances técnicos y el consiguiente incremento de la productividad en la agricultura (como con menos trabajo se produce más) originan un descenso de los precios agrícolas y no se necesitaría más tierra, ni más capital ni más trabajo, de suerte que se frenaría temporalmente el alza de las rentas, hasta que nuevos aumentos de la población requieran más alimentos y más tierras en cultivo.

\subsubsection{J.S. MILL} 
Su peculiar análisis estático lo combinó con el dinámico (que sí obedece a su definición) en la línea de Ricardo, pero sin visión pesimista. Así, contempló unas rentas en alza, unos beneficios en baja, unos salarios que van mejorando en poder adquisitivo y una disminución de las desigualdades sociales. La tendencia a la disminución de la tasa de beneficio sería la causa por la que irían desapareciendo las oportunidades rentables de inversión y, al final, provocaría la aparición del estado estacionario, en el que cesarían los avances tecnológicos y el progreso social. Para Mill este estado sería bastante confortable, porque en él no habría presión de la superpoblación ni problema de subconsumo. Su estado estacionario no produciría apresuramiento; en él no sería necesaria “una lucha constante por avanzar...[ni] pisotear, empujar, dar codazos y pisarle los talones al que va delante, que son características del tipo actual de vida social” (ibídem, p. 641). En dicho estado estacionario habría que tener en cuenta que “la mejor situación para la naturaleza humana es aquélla en la cual, mientras nadie es pobre, nadie desea tampoco ser más rico ni tiene ningún motivo para temer ser rechazado por los esfuerzos de otros que quieren adelantarse” (ibídem, p. 641). Para Mill (ibídem, p. 643), “una situación estacionaria del capital y de la población no implica una situación estacionaria del adelanto humano. Sería más amplio que nunca el campo para la cultura del entendimiento y para el progreso moral y social”. Evidentemente, en un estado tal los filósofos no tendrían el más mínimo inconveniente en vivir, disfrutando de los placeres intelectuales y morales; la sociedad gozaría de una adecuada prosperidad y la distribución de la renta sería más justa.




Según MILL, la distribución de la renta depende, principalmente, de las instituciones (no de leyes naturales como proponían ADAM SMITH y DAVID RICARDO). Sin embargo, sigue la misma clasificación para el reparto de la produción nacional entre los propietarios de los factores se corresponde con tres tipos de retribuciones: salarios, beneficios y rentas. 

\paragraph{Salarios} 
MILL se acogió inicialmente a la idea ricardiana y de su padre (procedente de SMITH) del fondo de salarios. Pero acabó rechazándola por ser muy poco laxa.\footnote{MILL tenía sentimientos favorables hacia los trabajadores y una preocupación por mejorar sus condiciones de vida y de ahí que concibiera con mucha flexibilidad la teoría del fondo de salarios, que hasta entonces había sido considerada con gran rigidez. Empezó por tener en cuenta que dicho fondo podía aumentar mucho con el progreso material y que, incluso, podía variar a corto plazo a costa del elevado tren de vida de los empresarios y propietarios de las empresas; de acuerdo con esto, el alza de los salarios a corto plazo sólo tropezaba con el límite que imponía el peligro de ruina de la empresa.} En estas circunstancias el salario efectivo se determinaba por la interacción de la oferta y la demanda, mientras que el salario del mínimo de subsistencia quedaba relegado a un límite extremo. Dada su idea de la posibilidad de variación del fondo a medio y largo plazo, debido al progreso material, MILL también contempló la negociación colectiva y los convenios laborales como medio de fijación del nivel salarial.\footnote{Por el interés que tiene esta ruptura de J. S. Mill con la concepción tópica del fondo de salarios, se transcriben a continuación sus ideas extractadas por Ashley (editor de su libro, Apéndice O pp. 850 y 851):

\begin{adjustwidth}{2cm}{0pt}
\raggedright
\textit{(...) ¿existe algo así como un fondo de salarios, en el sentido que implica esa teoría? ¿Existe una cantidad fija que sea, y ni más ni menos, la destinada a gastarse en salarios? Naturalmente la cantidad que puede gastarse con ese fin tiene un límite del cual no puede pasar: no puede exceder de la totalidad de los medios de que disponen las clases patronales. Y ni aun puede llegar a ese límite, pues los patrones tienen que mantenerse a sí mismos y a sus familias. Pero, por bajo de ese límite, no es, en ningún sentido de la palabra, una cantidad fija.\\[...] No hay ninguna ley natural que impida que los salarios suban hasta el punto de absorber no sólo los fondos que el capitalista pensaba dedicar a su negocio, sino hasta la totalidad de lo que éste se permitía para sus gastos privados, más allá de lo necesario para la vida. El límite real a la subida depende de la apreciación personal del capitalista sobre en qué forma le ocasionará la ruina o le obligará a abandonar el negocio; no de los límites inexorables del fondo de salarios.}\\
\raggedleft{\authorfont{J.S. MILL}, “\textit{Principios de Economía Política con algunas de sus aplicaciones a la Filosofía Social}”, 1848}
\end{adjustwidth}}

\paragraph{Los beneficios del capital} 
Mill expone una teoría de la abstinencia, según la cual a la gente, para inducirla a ahorrar, se le debe remunerar por su sacrificio ya que absteniéndose del consumo presente se obtiene el capital.

La ganancia bruta consta entonces de tres partes, que remuneran: la abstinencia, el riesgo, y el esfuerzo. Cada una de ellas recibe el nombre de interés, seguro y sueldos de dirección. Sin embargo, a la hora de determinar la tasa de beneficio, se encontró con una gran variedad de tipos de beneficio, en distintas clases de industrias y empresas. Esas diferencias de beneficios las atribuyó a diferencias en el riesgo o en las distintas formas de empleo del capital, y a monopolios naturales y artificiales.\footnote{Por otra parte, observó la tendencia a la nivelación de las tasas de beneficio, que preconizaba RICARDO; se trataba de una \textit{tasa ordinaria de ganancia} que equivalía a \textit{las ganancias medias que se obtienen en las diferentes formas de empleo del capital}. MILL supuso que el traspaso del capital desde una forma de empleo a otra no es una operación tan onerosa, lenta y casi impracticable como se la suele presentar.} MILL aceptó la idea de que los beneficios podían depender del nivel general de precios, según fuera la coyuntura social. Pero sacó poco partido de esta apreciación; únicamente la consideración de la existencia de unos beneficios extraordinarios debido a un alza general del nivel de precios y que tenían una naturaleza análoga a la renta de la tierra.\footnote{Que la renta no era algo exclusivo de la tierra ya lo había anticipado, pero atribuyéndolo a cualquier bien que pudiera ser monopolizado. En estas ideas se encuentra el origen del concepto de cuasi-renta y que prácticamente se deduce de un estudio un poco más a fondo de la teoría de la distribución de RICARDO. Además, MILL también nos ofrece otra explicación diferente sobre los beneficios del capital. A esta teoría, al estilo de la de Ricardo los socialistas las denominaron \textit{teorías de la explotación}, porque el beneficio del capital se atribuía a la apropiación en  provecho propio del capitalista de un parte de la gran productividad del trabajo; es decir, se pagaba al obrero menos del valor de lo que contribuía a la producción.\\
Esto lo dice muy claro MILL: “La causa de la ganancia es que el trabajo produce más de lo preciso para su sustento”, y por eso, después de lograr sus medios de vida, los trabajadores “dispondrán de una parte de su tiempo sobrante, durante el cual podrán trabajar para el capitalista. Vemos así que la ganancia surge, no por el accidente del intercambio, sino por la fuerza productiva del trabajo”. } 

\paragraph{Las rentas} 
Mill prosiguió con la ya clásica concepción de la determinación de las rentas de la tierra según unas reglas diferentes a las que rigen en la formación de las restantes retribuciones de los factores de la producción. Precisó que la causa de la renta de la tierra se encontraba en el monopolio ejercido sobre ella; el monopolio existía siempre que la cantidad de un bien estuviera limitada y no se pudiera aumentar libremente su oferta. El elemento de novedad que introduce MILL es que la persona que ejerce el monopolio sobre un agente natural dispone de él con potestad exclusiva, por disposición de la sociedad.\footnote{Así, \textit{los propietarios de tierras forman la única clase, numerosa o importante, que tiene derecho a una parte en la distribución de los productos por el hecho de ser dueños de algo que ni ellos ni nadie ha producido}. No todas las tierras dan renta, primero, porque en casi ningún país se explotan todas las tierras; siempre suele haber alguna sobrante. Y, segundo, porque hay tierras, de las peores, en las que sólo se puede extraer la reposición de las materias primas y el alimento de quienes las trabajan. En consecuencia, se puede sentar como principio que \textit{la peor tierra que se cultive no da renta}. Y, \textit{lo que una tierra cualquiera produce por encima de la tasa ordinaria de ganancia del capital es lo que la misma tierra produce en exceso de lo producido por la peor tierra en cultivo. El excedente es lo que el campesino puede pagar como renta al terrateniente}. La competencia entre los agricultores es la que permite igualar las tasas ordinarias de ganancia del capital y que el resto vaya a parar a la renta.}

\subsection{Ley de Say}
El núcleo de la cuestión, para Say, era que todo vendedor, por el mero hecho de serlo, es un comprador, ya que, en realidad intercambia mercancías (el dinero también es una mercancía).


En economía, la ley de Say es un principio atribuido a JEAN-BAPTISTE SAY (1803) que indica que la demanda está determinada por la producción, y que solo produciendo se puede generar demanda. Cuantos más bienes –para los que hay demanda– se produzcan, más bienes existirán (oferta) que constituirán a su vez demanda para otros bienes.

Dicha ley suele interpretarse mal en la mayoría de los casos. No es correcto afirmar que, según Say, todo bien producido va a encontrar su demanda, sino que para demandar un bien deben ofertarse otros.
Dicho de otra forma; para que se pueda demandar un bien en concreto, es necesario que antes puedan venderse otros. Por necesidad, la demanda viene limitada por la oferta pero no en un contexto comercial, es decir la ley de Say no implica que todo tipo de producción vaya a encontrar automáticamente su demanda, sino más bien que si los individuos quieren gastar, primero deben producir.
Expresado en palabras del mismo SAY (1803):
“Un producto terminado ofrece, desde ese preciso instante, un mercado a otros productos por todo el monto de su valor. En efecto, cuando un productor termina un producto, su mayor deseo es venderlo, para que el valor de dicho producto no permanezca improductivo en sus manos. Pero no está menos apresurado por deshacerse del dinero que le provee su venta, para que el valor del dinero tampoco quede improductivo. Ahora bien, no podemos deshacernos del dinero más que motivados por el deseo de comprar un producto cualquiera. Vemos entonces que el simple hecho de la formación de un producto abre, desde ese preciso instante, un mercado a otros productos.”

JAMES MILL (padre de J.S. MILL) fue un defensor acérrimo de la ley de Say, y afirmó:
“Si el poder de compra de una nación se mide exactamente por el producto anual (...), cuanto más se incremente el producto anual, más - por ese mismo acto - se expandirá el mercado nacional, el poder de compra y las compras reales de la nación”


\textbf{MUY IMPORTANTE}
\subsubsection{DAVID RICARDO} 
En el Capítulo 21 «Efectos de la acumulación sobre las utilidades y el interés», de su libro Principios de economía política y tributación, Ricardo (ibídem, p. 216-217) cita y se muestra de acuerdo con Say cuando expresa su opinión: "M. Say ha evidenciado en forma muy satisfactoria, sin embargo, que no hay cantidad de capital que no pueda ser empleada en un país, porque la demanda está limitada únicamente por la producción. Ningún hombre produce si no es para consumir o vender, y nunca vende si no es con la intención de comprar alguna otra mercancía, que le pueda ser de utilidad inmediata, o que pueda contribuir a una producción futura. Al producir, entonces, el hombre se transforma necesariamente en consumidor de sus propios productos, o en comprador y consumidor de los productos de alguna otra persona. No cabe suponer que el hombre se mantenga, por largo tiempo, mal informado acerca de las mercancías que él puede producir con más ventaja, para lograr la finalidad que persigue, a saber, la posesión de otros bienes: y, por lo tanto, no es probable que continúe produciendo una mercancía de la cual no existe demanda". En principio, Ricardo está tan de acuerdo con la Ley de Say que ni siquiera se le pasa por la imaginación que alguien desee demandar dinero por el mero hecho de tenerlo, como depósito de valor, hasta el momento futuro de utilizarlo. Para él (cuya opinión -ibídem, p. 217- era que una demanda sólo es efectiva si está respaldada por poder adquisitivo) no hay dilación entre la obtención de un valor y su empleo, ya que en sus razonamientos estaba implícito un peculiar principio de racionalidad económica (aunque en general bastante extendido entre los economistas) según el cual de todo capital disponible hay que obtener un rendimiento de inmediato: “Si se dieran 10.000 £ a un hombre que tiene 100.000 £ al año, no las encerraría en un cofre sino que aumentaría sus gastos en esas 10.000 £, las emplearía para fines de producción o las prestaría a alguna otra persona con el mismo propósito; en cualquier caso, la demanda aumentaría, aun siendo para diversos objetos. Si él mismo aumentó sus gastos, acaso su demanda efectiva sería de construcciones, muebles o algún otro disfrute. Si empleó productivamente sus 10.000 £, su demanda efectiva sería de alimentos, ropa y materias primas, que pondrían a trabajar nuevos obreros: pero de todos modos sería demanda” (ibídem, p. 217). Y por eso no duda en decir algo muy parecido a lo que afirmaba Say: “Las producciones se compran siempre con producciones, o con servicios; el dinero es únicamente el medio por el cual se efectúa el cambio” (ibídem, p. 217-218). No obstante, Ricardo no fue un incondicional seguidor de Say; llegó a conclusiones distintas a las de Say en lo que respecta a la relación entre el tipo de interés, la tasa de beneficios y la acumulación del capital. Say opinaba que "la abundancia de capitales disponibles, en proporción a la magnitud de las actividades que para ellos existía, bajaría la tasa de interés" (citado por Ricardo, ibídem, p. 217). Pero afirmaba que con tipos de interés a la baja la posibilidad de inversión (formación de capital) podría seguir indefinidamente (Schumpeter, 1954, p. 689n). Ricardo pensó que esto era contradictorio. Creía que para proseguir con la acumulación de capital se necesitaba que éste tuviera un empleo productivo, para lo cual se requeriría una adecuada tasa de beneficio. Y sólo había un motivo por el cual la tasa de beneficio no era la suficiente para inducir la inversión: el alza de los salarios. Sólo por el alza de los salarios se podía llegar a una situación en la que "queden tan pocas ganancias al capital, que cese el motivo de acumulación" (Ricardo, 1817, p. 217). A modo de conclusión, y en una apretada síntesis, diríamos que el mensaje que transmitió Ricardo fue: ¡No subáis los salarios a los obreros: expolian nuestros beneficios!

Aceptación de la ley de Say: La producción no sólo genera bienes que satisfacen necesidades, sino que permite comprar otros productos ya que para SAY la producción genera unos medios económicos para comprar otros productos. Es decir, los bienes se pagan con bienes. La ley de
Say genera una gran controversia y no existe consenso sobre su interpretación:
o
RICARDO está de acuerdo con y afirma que “nadie produce con otro fin que el fin de vender la producción para comprar.”
o
En cambio, MALTHUS se opone frontalmente a la ley de Say al señalar que no todas las rentas de la producción serán gastadas. En este sentido, MALTHUS influyó en las ideas de JOHN MAYNARD KEYNES32.

\subsubsection{MALTHUS}
Según SMITH y RICARDO, el progreso económico depende de la acumulación de capital, ya que realizan análisis basados exclusivamente en la oferta.
‒
Para MALTHUS es necesario también tener en cuenta el lado de la demanda, o lo que él llamaba “demanda adecuada”.

MALTHUS rechaza la ley de Say en la concepción de JAMES MILL, donde la ley de Say se amplía en conceptualización para incluir fenómenos intertemporales: si el ahorro es igual a la inversión se cumpliría el equilibrio oferta igual a demanda.
‒
MALTHUS cree que no todas las rentas derivadas de la producción se gastan pues existe una parte del ahorro que no se usa para invertir, sino que se atesora.
o
En este sentido, distingue el hombre frugal de negocios del ávaro que atesora sus ahorros y cierra bajo llave la facilidad de producir.
•
Admite que la inversión requiere ahorro pero precisa que no todos los ahorros son invertidos necesariamente.
‒
Por tanto, para que los recursos se empleen es necesario que la demanda siga la oferta, en caso contrario puede darse una situación de subconsumo.
o
El problema es que los trabajadores cobran el salario de subsistencia: no crean demanda por sí mismos.
•
Por ello, MALTHUS critica el auxilio a los pobres, pues no constituyen la parte más importante de la demanda.
o
La demanda se crea en el sector de los ricos y sobre todo de los terratenientes. En la medida en que estos terratenientes ahorren más que gasten, caerá la demanda efectiva global y provocará un estancamiento a largo plazo.
•
Por ello, hay que estimular la demanda de los terratenientes (por ejemplo, mediante medidas proteccionistas que defiendan la demanda interna de los terratenientes).\footnote{Por ello apoya MALTHUS la Ley de Granos.}

\subsubsection{J.S. MILL} 
MILL intentó coordinar las opiniones de SAY, RICARDO y su padre con las de sus detractores, MALTHUS y SISMONDI, partidarios de las crisis de superproducción (plétoras generales debidas al subconsumo) y de la inadecuación del ahorro y la inversión. 

Mill considera que la Ley de Say es incuestionable, bajo su concepción del dinero neutral. Con estas hipótesis, no manifestadas expresamente, MILL al tratar los cuatro teoremas o proposiciones fundamentales del capital realiza verdaderos esfuerzos para demostrar el principio de la Ley de Say, a saber, que la producción es el elemento más importante de la economía, y ella provee a los consumidores de los medios para sostener una demanda efectiva suficiente para retirar toda la oferta global. 

Se ha mencionado que Mill hace verdaderos esfuerzos, más bien encajes de bolillos, porque dice que si un producto no tiene demanda no se produce; pero su fe ciega en la Ley de Say enseguida le hace decir que un producto puede tener mucha demanda, pero de nada sirve si no hay producción de él. A lo sumo, para lo que sirve la demanda es para orientar la producción y en qué ramas de la industria habrá más o menos trabajadores: 

“Lo que sostiene y emplea el trabajo productivo es el capital que se gasta para ponerlo en actividad y no la demanda de los compradores del producto acabado del trabajo. Demanda de mercancías no equivale a demanda de brazos. Aquélla determina en qué rama particular de la producción se emplearán el trabajo y el capital, determina la dirección del trabajo, pero no el más o el menos del trabajo en sí, o del mantenimiento y el pago del trabajo”(ibídem, p. 92-93). 

Casi no se puede comprender que un MILL que expone el símil de las tijeras para manifestar el sin sentido de intentar demostrar que una causa es la más importante cuando son dos las que intervienen en un asunto; que confía en e la interacción de la oferta y la demanda; y que establece relaciones dialécticas entre causas y efectos, se empecine en demostrar que es la oferta la que crea la demanda, cuando reconoce que, a la inversa, sin demanda no se produciría una mercancía. 

Una explicación posible es acudir al amor filial, ya que su padre fue el primer autor inglés que expresó una tesis similar a la Ley de Say: “No es sino un acto de justicia a dos nombres eminentes llamar la atención sobre el hecho de que el mérito de haber situado este asunto en forma que se pueda apreciar debidamente pertenece, sobre todo en el continente, al juicioso J. B. Say, y en este país a Mr. [James] Mill”. 

No obstante, la realidad es que MILL no estaba tan incondicionalmente de acuerdo con la Ley de Say; al menos, supo ver varias excepciones (eso sí, sin que afectaran de forma definitiva a la Ley): 

“Primera, cuando el trabajo se sostiene, pero no se ocupa en su totalidad, una nueva demanda de algo que puede producir, tal vez estimule al trabajo así sostenido a realizar mayores esfuerzos, cuyo resultado puede ser un aumento de la riqueza, en provecho de los trabajadores mismos y de los demás.” También dice: “Fallando así el fundamento de nuestro teorema, éste falla también, y al surgir la demanda de una mercancía, puede hacer aparecer empleo de esta clase sin despojar al trabajo de una cantidad equivalente de empleo en otra rama de la actividad. Aun en este caso, la demanda  no actúa sobre el trabajo más que a través del capital  existente, pero ofrece un incentivo que hace que el capital ponga en movimiento una cantidad de trabajo mayor que antes” (ibídem, p. 90). 

“La segunda excepción, [...], consiste en el conocido efecto de una ampliación del mercado de una mercancía, que permite aumentar la división del trabajo y en consecuencia una distribución más eficaz de las fuerzas productivas de la sociedad. Esta, como la anterior, es una excepción más aparente que real. El trabajo no se remunera con el dinero desembolsado por el consumidor sino con el capital del productor; la demanda sólo decide la forma en que se ha de emplear el capital y qué clase de trabajo ha de remunerarse; pero si se decide que la mercancía se ha de producir en gran escala, permite que con el mismo capital se produzca mayor cantidad de la mercancía, y, de manera indirecta, siendo la causa de un aumento del capital, puede producir un aumento eventual en la remuneración del trabajador” (ibídem, p. 100). 

También puede ocurrir algún fallo (temporal) en la ley de Say cuando ocurren las llamadas crisis comerciales en economías monetarias. Entonces (ib., p. 487), la utilización del dinero en los intercambios introduce una distorsión, pues disocia el momento de la venta con el de la compra por parte del vendedor; por consiguiente, la demanda de otras mercancías puede no ser inmediata, lo que en terminología moderna se expresaría diciendo que, en ese caso, existe una mayor preferencia por la liquidez. Cuando existen retardos en la demanda de otras mercancías se origina una suboferta (escasez) de dinero y se puede llegar a generalizar una  situación en la que todo el mundo quiere vender (preferencia por el dinero) y nadie quiere comprar (o desprenderse del dinero). El resultado, en efecto, es una crisis económica, con un exceso de mercancías y un descenso de los precios. Sin embargo, Mill piensa que esa situación, aunque pueda prolongarse por un cierto tiempo, no será demasiado duradera y al final le seguirá un periodo de intensa demanda. 

La interpretación de J.S. Mill respecto a la ley de Say es muy interesante por los siguientes aspectos: 
1º. El estar de acuerdo con la ley de Say no es un impedimento para reconocer la existencia temporal de plétoras generales. 
2º. La Ley de Say se considera como una igualdad y no como una identidad. 
3º. Las consideraciones monetarias tienen una gran trascendencia en la determinación de las magnitudes reales de la economía (aunque Mill, por lo general, considere al dinero como neutral). Este aspecto introduce en la economía monetaria moderna, donde la demanda de liquidez es equivalente a la suboferta de dinero. 
4º. La explicación de Mill constituye una incipiente teoría del ciclo económico.


\subsubsection{Economía Monetaria}
\paragraph{ADAM SMITH}




\paragraph{DAVID RICARDO} 
Ricardo se inició en la teoría económica al participar en el debate monetario. Ya desde su estreno como economista demostró énfasis y sugestión en sus escritos, logrando que sus ideas se impusieran, pese a que la calidad de su análisis fuera muy inferior al de Thornton (que se había anticipando a la capacidad de entendimiento de su época y que, además, escribía llanamente, sin excesivo énfasis, aunque, por supuesto, con pretensión de ser convincente). Ya en sus artículos de El precio del oro (1809) y El alto precio de los lingotes (1810), Ricardo se decantaba por los estrictos postulados metalistas y cuantitativistas del dinero, a pesar de recomendar la circulación del papel moneda y descartar las monedas de oro y plata. Los billetes de banco que él recomienda debían estar respaldados por las reservas de oro, aunque no necesariamente en su totalidad, y sobre todo debían ser libremente convertibles en lingotes (Ricardo, 1817, pp. 266 y 269). En estas condiciones el oro monetario y los billetes son exactamente lo mismo, o sea, dinero y no medios de pago crediticios, ya que para Ricardo el dinero debía ser neutral (esto es, el dinero no afecta a lasvariables reales de la economía), pues como él dice (ib., p. 218): “es únicamente el medio por el cual se efectúa el cambio”. Y asimismo dice: “Como el dinero es unbien variable, el aumento de los salarios en dinero será frecuentemente ocasionado por una baja del valor del dinero. En efecto, un aumento de salarios debido a esta causa irá invariablemente acompañado de un aumento en el precio de los bienes; pero en tales casos, se observará que la mano de obra y todos los bienes no han variado con respecto unos a otros, y que la variación ha quedado confinada al dinero [...]. Un aumento en los salarios, debido a una alteración en el valor del dinero, produce un efecto general sobre el precio, y por esa razón no produce ningún efecto real sobre las utilidades“ (ibídem, p. 36). Sin embargo, pese a sus tesis bullionistas, Ricardo (ibídem, p. 268) no culpaba al Banco de Inglaterra de ser el causante de la inflación ni de la crisis de 1797, pues opinaba que había ejercido sus poderes con moderación, aunque reconoce que había emitido más billetes de los que podía garantizar con sus reservas. Con su Plan lingote o Plan Ricardo (1811), que consistía en volver a la convertibilidad de los billetes en oro, pero en lingotes (es decir, se trataba de un «patrón lingote oro»), Ricardo pretendía que la emisión de billetes tuviera un control (el marcado por la exigencia de la convertibilidad), de modo que el valor del billete fuera exactamente igual que el del oro al que representaba (ibídem, p. 269). Ricardo (ibídem, pp., 269-270) también estudió las ventajas e inconvenientes de que el gobierno asumiera el monopolio de la emisión de los billetes para que los beneficios de la emisión no fueran a parar a las manos privadas del Banco de Inglaterra; pero no parece que se decantara por una propuesta en concreto. Por haber asumido aquellos dos postulados del dinero (metalista y cuantitativista) admitió la influencia de la cantidad de dinero en los precios a través de las alteraciones en la demanda (ibídem, pp., 222-223) pero Ricardo no pudo comprender las otras tesis de Thornton (Spiegel, p. 374), a saber: La producción podía fomentarse con una expansión monetaria o crediticia. El aumento de la cantidad de dinero en circulación no tenía por qué afectar necesaria y directamente al nivel de precios (sólo a través de otras variables como la demanda, el retraso de la producción ante el estímulo de la demanda, el alza de los costes, o el tipo de interés). La inflación podía originar un ahorro forzoso con el que se contribuiría a la formación de capital. Los movimientos internacionales de capital financiero podían realizarse por las diferencias existentes en los tipos de interés.


Finalmente, en relación a la teoría monetaria\footnote{RICARDO apoya las tesis de la currency school frente a la banking school.} se decantaba por los postulados cuantitativistas del dinero (la inflación contemporánea era resultado de una emisión excesiva de billetes). RICARDO recomienda un respaldo de los billetes en lingotes de oro.\footnote{Sin embargo, en Inglaterra se acababa de suspender la convertibilidad en oro de los billetes del Banco de Inglaterra. El valor nominal de los billetes emitidos era muy superior al valor del oro que el banco mantenía en reserva.
Una de las causas del gran aumento de los billetes del Banco de Inglaterra fue la necesidad del gobierno británico de recurrir a los anticipos de dicho banco para financiar las guerras contra Francia (1793-1815).
En este debate, los bullionistas como RICARDO querían la restauración del patrón oro en lingotes.}


\paragraph{J.S. MILL}
En lo que respecta a la teoría del dinero, Mill no hizo ninguna aportación original. Básicamente, se acogió a la teoría metalista, a la neutralidad del dinero y a la teoría cuantitativa estricta. En lo que a esto último se refiere, expuso una ecuación de cambios relacionando flujos, el de la totalidad de las mercancías vendidas durante “el proceso” y el de la cantidad de dinero multiplicada por su velocidad de circulación; de modo que el valor del dinero es inversamente proporcional a este último producto, dado un volumen de transacciones (ibídem, p. 433). Sin embargo, realizó una meritoria labor de recopilación, perfeccionamiento y conciliación de todos los logros alcanzados en la primera mitad del siglo XIX

En política monetaria, MILL extiende el mecanismo flujo-especie de HUME. Aumento del tipo de interés vía venta de letras o aumento del tipo de descuento del Banco de Inglaterra. Capital entra en el país. Aumenta stock de oro/oferta monetaria. En tiempos normales banking school. En situaciones de booms de crédito currency school.

\subsection{Comercio Internacional}

\subsubsection{ADAM SMITH}
Muy crítico con los mercantilistas y propenso a la teoría cuantitativa, considera imposible que la balanza comercial se mantenga constantemente favorable ya que la afluencia de dinero provoca un aumento de precios; es una ley natural que los metales preciosos vayan a los países que tienen los precios bajos. SMITH no es partidario del proteccionismo; ni siquiera para la industria naciente ya que los capitales no se orientarían hacia las actividades productivas más rentables. Para SMITH: 
\begin{lnum}
    \item El comercio más provechoso es el interior, ya que fomenta el trabajo productivo originando una mayor producción nacional que la que se obtendría con el mismo capital dedicado al comercio exterior; y,
    \item El comercio exterior cubre la función de dar salida a los excedentes del mercado interior y proporciona oportunidades rentables para utilizar el capital acumulado que no se precise para la producción interior. Pero, como el principio de las ventajas absolutas es el que rige en el comercio internacional y en el nacional: el capital se desplaza de un lugar a otro, según tenga oportunidad de comprar barato o vender caro: la industria nacional debe orientarse hacia la producción que sea más ventajosa, que nunca puede ser aquélla que ya se fabrica fuera más barata.
\end{lnum} 

En su pensamiento late la idea de lo absurdo que es fabricar en casa lo que se puede comprar más barato fuera. Además, defenderá que el comercio debe establecerse libremente, sin estímulos forzados, para que beneficie a ambas partes: \textit{Aquel comercio que, sin fuerza ni violencia, se desarrolla de una manera normal entre los dos países, es siempre ventajoso, aun cuando la ventaja no sea la misma para las dos partes}.

\subsubsection{DAVID RICARDO}
David Ricardo defendió la libertad de comercio exterior; sobre todo para la importación de productos agrícolas con el objeto de mantener baratas las subsistencias e impedir así la elevación de los salarios nominales y, con ellos, la de las rentas y la disminución de la tasa de beneficios. El comercio exterior proporcionaba ventajas en dos órdenes:
\begin{lnum}
    \item \textit{En la distribución.} El comercio exterior no ejerce, en principio, ninguna influencia directa sobre la distribución, porque no modifica la cantidad de trabajo empleada en la producción interior de los bienes, ni altera el importe de los salarios y, por lo tanto, tampoco afecta a la tasa de beneficios y a las rentas. Sin embargo, si el comercio exterior proporciona alimentos más baratos, entonces, indirectamente, sí afecta a la distribución, porque hace disminuir los salarios nominales y, por consiguiente, aumentan los beneficios empresariales; y,
    \item \textit{En la renta real.} Según su teoría del valor, el comercio exterior no altera los valores relativos de la producción interior, pero la mayor disponibilidad de mercancías, que posibilita al comercio exterior, mejora la renta real. Esta ventaja del comercio exterior puede extenderse a todos los ciudadanos, pues, en cuanto que consumidores, aumentan sus posibilidades de disfrutar de más bienes (aunque en principio, como se ha visto, no se altere la distribución, excepto por la disminución del salario nominal).
\end{lnum}

Para demostrar este argumento de la mejora de la renta real, tan atractivo para el público, ideó el principio de las ventajas comparativas.\footnote{Este principio es su gran contribución al análisis económico, puesto que, en determinadas circunstancias, con él se demuestran las ventajas mutuas del comercio; aunque su eficacia práctica quede muy restringida: sólo a los acuerdos comerciales bilaterales cuando existan esas determinadas circunstancias.} 

\paragraph{Teoría de la Ventaja Comparativa}
En primer lugar, RICARDO parte de la base de que la teoría de la ventaja absoluta implicaría la no participación en el comercio internacional de países menos avanzados.\footnote{Sin embargo, según la teoría de la ventaja comparativa de RICARDO para estos países el comercio internacional puede ser beneficioso.} En el capítulo 7 de su obra \textit{Principios de Economía Política y Tributación} (1817), explica por qué es posible que exista comercio internacional cuando un país tiene ventaja absoluta en la producción de todos los bienes.\footnote{La teoría de la ventaja comparativa se atribuye a menudo a DAVID RICARDO. Sin embargo, es posible encontrar enunciados previos de esta teoría, por ejemplo:
\begin{lnum}
    \item En una obra publicada dos años antes por el también británico ROBERT TORRENS (1815);
    \item En una obra publicada anónimamente en 1814 titulada \textit{Considerations on the Importation of Foreign Corn}.
\end{lnum}
En cualquier caso, de acuerdo con GRANCAY y GRANCAY (2015), si bien estos autores mostraron un conocimiento muy decente de lo que hoy se conoce como teoría del comercio internacional, sus enunciados de esta teoría no se consideran satisfactorios ya que son incompletos. Por lo tanto, consideraremos a RICARDO como el primer autor en enunciar la teoría de la ventaja comparativa de forma satisfactoria. Él no acuñó el término ventaja comparativa y probablemente no fuera él el que inventara el principio, pero a día de hoy no se ha encontrado ninguna prueba de ningún otro autor a quien pueda atribuirse este principio.} Éste defiende que las diferencias relativas de los costes de producción debidas a la tecnología son el motor del comercio internacional. Así, según esta teoría un país se especializará y exportará aquellos bienes que produce a un menor coste relativo, es decir, a un menor coste en términos de los otros bienes. Por lo tanto, la variable crucial para explicar la existencia y el patrón del comercio internacional es la tecnología.\footnote{Una diferencia en los costes comparativos de producción –que es condición necesaria en este modelo para que exista intercambio internacional– refleja, de hecho, una diferencia en las técnicas de producción.}\textsuperscript{,}\footnote{Para RICARDO, el comercio internacional y el nacional se rigen por principios diferentes. No se les puede aplicar las mismas reglas porque en el comercio interior los bienes se intercambian por su valor de cambio que depende del trabajo incorporado. En el comercio exterior \textit{se puede cambiar el producto del trabajo de 100 ingleses por el producto de la labor de 80 portugueses, 60 rusos, ó 120 indios orientales}, debido a \textit{la dificultad con que se mueve el capital de un país a otro}, por \textit{la inseguridad real o imaginaria del capital, cuando éste no está bajo el control inmediato de su dueño, aunada a la natural renuencia que siente cada persona a abandonar su país de origen y sus relaciones con nuevas leyes [en] un país extraño}.}

Simplificándolo al máximo y siguiendo el ejemplo de DAVID RICARDO, podemos suponer:
\begin{lnum}
    \item Dos países (Inglaterra y Portugal);
    \item Dos bienes (paño y vino) homogéneos entre sí;
    \item Un factor productivo (trabajo)\footnote{Todos los factores de producción se pueden reducir a uno único, el trabajo. Este supuesto se basa en la Teoría del Valor-Trabajo. Sin embargo, la validez de la teoría clásica del comercio internacional no se basa en la validez de la teoría clásica del valor trabajo, ya que es suficiente con que los costes unitarios de producción sean medibles por una unidad común entre los países que sea constante.}, que será móvil dentro del país, pero inmóvil entre países.
    \item Los países tienen distinta tecnología (i.e. el factor de producción es homogéneo dentro del país, pero heterogéneo entre países) y la dotación factorial en cada uno de ellos puede ser la misma o diferente.
    \item La tecnología se representa mediante una función de producción de coeficientes técnicos fijos (i.e. tecnología lineal).\footnote{Dicho de otra forma, en ambos países, la producción de bienes se realiza de forma que los costes unitarios de la producción de cada bien (expresados en términos de trabajo) sean constantes. Este supuesto implica rendimientos constantes a escala y productividad marginal constante (i.e. PMgL=PMeL).}; y,
    \item Movilidad perfecta de bienes entre países, no existen costes de transporte.
\end{lnum}

Esta teoría afirma, que una condición necesaria para que exista comercio internacional es una diferencia en costes comparativos\footnote{Los costes comparativos se pueden definir de dos maneras:
\begin{lnum}
    \item Como la ratio ente los costes unitarios absolutos de los dos bienes en el mismo país; o,
    \item Como la ratio ente los costes unitarios absolutos del mismo bien en los dos países.
\end{lnum}

Ambos métodos son idénticos, pero en esta exposición, tal y como es práctica habitual, adoptaremos la primera opción.
Si definimos como $a_1$ y $a_2$ los costes unitarios del bien a en los dos países (donde la letra se refiere al bien y el subíndice numérico al país) y los costes unitarios del otro bien como $b_1$ y $b_2$, entonces:
$$\left(\frac{a_1}{b_1}=\frac{a_2}{b_2}\right)\iff\left(\frac{b_1}{a_1}=\frac{b_2}{a_2}\right)\iff\left(\frac{a_1}{a_2}=\frac{b_1}{b_2}\right)\iff \left(\frac{a_2}{a_1}=\frac{b_2}{b_1}\right)\;\;\text{y}\;\;
\left(\frac{a_1}{b_1}\lessgtr\frac{a_2}{b_2}\right)\iff\left(\frac{b_1}{a_1}\lessgtr\frac{b_2}{a_2}\right)\iff\left(\frac{a_1}{a_2}\lessgtr\frac{b_1}{b_2}\right)\iff \left(\frac{a_2}{a_1}\lessgtr\frac{b_2}{b_1}\right)$$}. Sin embargo, esta es únicamente una condición necesaria; la condición suficiente es que la relación real de intercambio se encuentre entre los costes comparativos2\footnote{Se puede demostrar fácilmente que la relación real de intercambio ha de encontrarse estrictamente entre los dos costes comparativos, ya que, si la relación real de intercambio fuera igual al coste comparativo de un país, ese país no tendría interés alguno en comerciar, ya que la ratio de precios interno (dado por el coste comparativo) sería igual al internacional (la relación real de intercambio).}. Cuando ambas condiciones se dan, podemos afirmar que será beneficioso para ambos países especializarse en la producción del bien en el que tienen ventaja comparativa.

\imagenfit{Ricardo_VenComparativa.png}{Ejemplo de ventaja comparativa}{Apuntes Víctor Gutiérrez, Tema 3.A.2. Repositorio TCEE}

Considerando el ejemplo, Portugal es superior a Inglaterra en la producción de ambos bienes, por lo que podría parecer que el comercio internacional no es factible. Sin embargo, vemos como Portugal tiene una ventaja relativa mayor (ventaja comparativa) en la producción de vino, mientras que Inglaterra tiene una desventaja relativa menor (ventaja comparativa) en la producción de paños. Por lo tanto, se da la condición necesaria para que exista comercio internacional.\footnote{Si se diera la condición suficiente se producirá el intercambio de vino portugués por paños ingleses.} Se podría demostrar fácilmente que la relación real de intercambio se ha de hallar comprendida estrictamente entre ambos costes comparativos, ya que, si fuera igual a alguno de ellos, ese país no tendría interés en comerciar debido a que su precio interno (su coste comparativo de producir) sería igual al precio internacional (la relación real de intercambio). En otras palabras, podrían obtener la misma cantidad mediante el comercio internacional que mediante la producción interna. Además, de este análisis resulta obvio que, si la relación real de intercambio se encuentra fuera de este intervalo, uno de los dos países sufriría pérdidas al realizar el intercambio.

Del modelo se obtienen las siguientes implicaciones:
\begin{lnum}
    \item ¿Por qué comercian los países?: El comercio internacional surge por las diferencias en los precios relativos de los bienes causadas por las diferencias entre las tecnologías de los distintos países;
    \item ¿Qué producen?: Cada país produce únicamente el bien en el que tiene una ventaja comparativa, dando lugar a uno de los resultados fundamentales del modelo ricardiano: la especialización completa2\footnote{Nótese que la balanza comercial esta necesariamente en equilibrio, pues el valor de las exportaciones es igual al de las importaciones al producirse mediante trueque a la relación real de intercambio establecida.}\textsuperscript{,}\footnote{En cualquier caso, esto no tiene por qué ser el caso si un país es pequeño en relación al otro, de forma que la producción del bien producido por la economía pequeña no sea suficiente para satisfacer la demanda en el país grande. En este caso, el país pequeño se especializaría completamente, pero el país grande tendría que seguir produciendo parte del bien en el cual no tiene ventaja comparativa. Sin embargo, para analizar este caso es necesario incluir el lado de la demanda y curvas de indiferencia de los consumidores (ausente en el análisis original de DAVID RICARDO, al igual que todo el aparato gráfico desarrollado aquí abajo):
    \imagenfit{OfertaInsuficiente_ExtRICARDO.png}{Desarrollo gráfico. Economía pequeña con oferta insuficiente e impacto sobre la Teoría de la Ventaja Comparativa}{Apuntes Víctor Gutiérrez, Tema 3.A.2. Repositorio TCEE} En el caso representado sí que se produce especialización completa, pero, como decimos, con cambios en las demandas de los bienes, habría cambios en las curvas de indiferencia y se podría dar la otra solución comentada.}. Esto es así debido al supuesto de tecnología lineal;
    \item ¿Qué implicaciones tiene para el bienestar?: En la medida en que cada país se especializa en la producción de aquel bien que hace relativamente mejor que el resto, el comercio será beneficioso para ambos, sin que haya perdedores, ya que la producción total es mayor que en autarquía. De este modo, el comercio internacional aumenta las posibilidades de consumo de los países que participan en él al permitir que cada país obtenga de manera más eficiente el bien en el que no tiene una ventaja comparativa, produciendo el bien para el que tiene una ventaja comparativa e intercambiándolo por el otro bien;\footnote{Fuera de este modelo, RICARDO, en su teoría de la distribución aporta una ganancia adicional del libre comercio} y,
    \item ¿Cuál es la relación real de intercambio? Este modelo, al estar centrado exclusivamente en el lado de la oferta, no permite determinar la relación real de intercambio, es decir, el precio relativo al que se intercambian los bienes. Esta variable ha de ser asumida exógena al modelo pese a su gran relevancia para la existencia de comercio internacional (ya que ha de estar comprendida entre las Relaciones Marginales de Transformación de ambos países para que se cumpla la condición suficiente y exista comercio internacional). Como veremos, esto constituye una de las principales críticas al modelo.
\end{lnum}

Sin embargo, esta teoría presenta un gran inconveniente: no determina la relación real de intercambio (es decir, la relación de precios a la que se intercambian las mercancías) ni la cantidad intercambiada. Esto constituye una limitación de la teoría, ya que un modelo que busca explicar el comercio internacional no debería únicamente explicar las causas y el patrón de comercio, pero también la determinación de los precios.\footnote{Esta limitación puede ser considerada irrelevante, si, tal como sugiere BHAGWATI (1964), estudiamos la teoría ricardiana como un modelo simplificado que sirvió para demostrar que el comercio internacional es beneficioso, es decir, lo consideramos desde el punto de vista normativo y no tanto desde el punto de vista positivo.}

\subsubsection{J.S. MILL} 
Mill obtuvo uno de sus mayores logros en el campo de la teoría económica con su estudio sobre el comercio internacional.\footnote{MILL, modestamente, atribuye todo el mérito a DAVID RICARDO por su brillante exposición del principio de las ventajas comparativas. Él sólo pretende aclarar el alcance de ese principio en lo concerniente a la indeterminación de la relación de precios bajo la cual cada país obtiene simultáneamente una ganancia del comercio exterior. Pero en realidad, acabó construyendo una teoría nueva.} 

MILL se basa en la ley de la oferta y la demanda, que no es sino otra forma de expresar la demanda recíproca, para llegar a fijar la relación de precios internacionales, en la que se estabiliza el intercambio. Para ello realiza unas simplificaciones: supone que el intercambio se limita a dos mercancías y a dos países de similar dimensión y capacidad productiva, y que, en el intercambio, existe \textit{la ecuación de la demanda internacional}; es decir, los valores de las mercancías son iguales (o sea, 1000 € en vino se intercambian por 1.000 € en lana; dicho de otra forma, 1.000 € son 1.000 € ya sean en vino o en lana). También se supone que el Reino Unido tiene ventaja comparativa en la elaboración de paño y Portugal en la de vino. 

El mérito consistió en establecer la condición de equilibrio, cuando se igualan las ofertas y demandas del vino y del paño.\footnote{Esto no significa que no puedan darse muchas situaciones de equilibrio; por ejemplo, una distinta (siempre entre los dos límites de las relaciones de precios) cada vez que se desplace una curva de oferta o de demanda, al cambiar las condiciones bajo las cuales se establecen.} De las argumentaciones de Mill se pueden extraer tres importantes apreciaciones: 
\begin{lnum}
    \item Sus curvas de oferta y demanda eran funciones en las que las cantidades dependían de los precios;
    \item Las condiciones bajo las cuales se establecen la oferta y la demanda podían cambiar dando lugar a una nueva situación de equilibrio. Frente a sus críticos en esta cuestión, manifestó expresamente que su condición no era una proposición idéntica;
    \item Especifica la demanda en términos de elasticidad (en terminología de actualidad).Para él, cuanto más fuerte sea la demanda de un artículo, y más susceptible de aumentarse sea aquélla ante el abaratamiento de su precio, más dinero obtendrá el país que lo produce. \footnote{Algunos autores se asombran, y elevan a la categoría de genio a MILL, por el hecho de que usara razonamientos matemáticos sin los símbolos matemáticos. MARSHALL, al exponer sus teorías, manifestó que sólo hacía un tratamiento en forma de diagrama del problema de MILL de los valores internacionales y que sus curvas eran como poner en forma definida, algo presentado ya por MILL. En el desarrollo de la teoría de Mill, Marshall y más tarde Lerner establecieron la actualmente conocida condición Marshall-Lerner para la estabilidad del equilibrio del intercambio en el comercio exterior. La condición Marshall-Lerner dice que para restablecer el equilibrio de la balanza de pagos por cuenta de renta, ante una depreciación del tipo de cambio, se requiere que la suma de las elasticidades de la demanda exterior de nuestras mercancías y de la demanda interior de importaciones sea mayor que la unidad.} 
\end{lnum}

La determinación de la relación real de intercambio la lleva a cabo MILL a través de las llamadas demandas recíprocas, que es la demanda de cada país de los productos importados (y que, al mismo tiempo, determina la oferta de productos exportados), será aquella que haga que el valor que cualquier país reciba por sus exportaciones sea exactamente igual al que tenga que pagar por sus exportaciones.

\subsection{La política económica}

\subsubsection{ADAM SMITH}
En materia de política económica, ADAM SMITH aboga por un laissez-faire con unas funciones muy concretas reservadas para el sector público
‒
En materia de gastos, ADAM SMITH considera que el sector público sólo debe intervenir en los siguientes casos:
1.
Defensa frente al exterior, salvaguarda del orden público y administración de la justicia, infraestructuras y obras públicas.

2.
Otras: Lucha contra el monopolio, remover obstáculos a la libre competencia sin posiciones de privilegio y que lleve a la natural armonía de intereses, educación pública, regulación de patentes y derechos de propiedad.
‒
En materia de ingresos, ADAM SMITH realiza las siguientes apreciaciones:
o
Principios: Carga en función de la renta, simplicidad y fácil comprensión, costes de administración reducidos.
o
Defiende la recolección de impuestos que recaigan sobre la tierra. La idea subyacente es que las tierras se remuneran a precios supracompetitivos porque están limitadas. Sin embargo, matiza que deben recaer sobre la tierra y no sobre el producto de la tierra.
•
Así, también rechaza impuestos sobre el capital y los salarios porque entiende que son injustos y desincentivarían la producción.

\subsubsection{DAVID RICARDO}

RICARDO, al igual que SMITH defenderá las virtudes de un sistema autorregulado. Así, en materia de política económica se desprenden las siguientes recomendaciones de la obra de DAVID RICARDO:
‒
Al igual que SMITH defiende la financiación de bienes públicos con impuestos a la tierra y a los bienes de lujo e impuestos de suma fija.\footnote{No obstante, en el capítulo XVII de su obra \textit{Principios de Economía Política y Tributación}, DAVID RICARDO observa la emisión de deuda pública como forma alternativa de financiar estos gastos.

RICARDO planteará la posibilidad teórica de que impuestos y deuda pública puedan acabar siendo sustitutivos (equivalencia ricardiana). Los individuos perciben que esa mayor deuda hoy equivale a mayores impuestos futuros para financiar el gasto pasado. Por tanto, ahorrarán en forma de bonos una cantidad equivalente al incremento del gasto. En consecuencia, cae el consumo lo mismo que aumenta el gasto público; el efecto es nulo.

Sin embargo, el propio RICARDO dudaba de que esta proposición tuviera consecuencias prácticas. Siguió la exposición inicial con una afirmación de que los individuos en realidad no evalúan los impuestos de esa manera y, en particular, tienen una visión miope de la trayectoria fiscal.

Más de 150 años después, esta idea fue revivida por THOMAS BARRO en su artículo de 1974 \textit{Are Government Bonds Net Wealth?}. Dada una senda de gasto exógena, la distribución temporal de los impuestos es irrelevante: la asignación y el precio en un equilibrio competitivo queda inalterada. La restricción presupuestaria intertemporal queda inafectada. Menos impuestos hoy implican mayores impuestos mañana y ambos efectos se cancelan.}

Aboga por la libre circulación de bienes siendo partidario de eliminar trabas existentes al comercio internacional. Por ello, aboga por la derogación de la Ley de Granos.\footnote{RICARDO colaboró bastante con MALTHUS, criticando la política monetaria expansiva del Banco de Inglaterra. Sin embargo, con respecto a la Ley de Granos divergieron.}

Y, al igual que los demás contratos, se deberían dejar los salarios a la libre competencia en el mercado y nunca deberían ser controlados por la legislatura. RICARDO, siguiendo a Malthus, está totalmente en contra de las "leyes de pobres", ya que absorben muchos impuestos y recursos por la caridad en cada parroquia, y no fomentan los propios esfuerzos para ganarse la vida. También opina que no es más cierto el principio de gravitación universal que la tendencia de tales leyes a cambiar la riqueza y el poder, en miseria y debilidad [...] y así llegará un momento en que todas las clases sociales se verán infectadas por la plaga de la miseria universal". "Afortunadamente el periodo de vigencia de estas leyes ha sido de prosperidad progresiva [...] Pero de hacerse más lento nuestro progreso, si permanecemos en un nivel estacionario [...] entonces se hará más patente y alarmante la naturaleza perniciosa de estas leyes y, entonces, también su abrogación será obstaculizada por multitud de dificultades adicionales". "El mejor amigodel pobre, así como de la causa de la humanidad, será la persona que pueda señalar un modo de abolir estas leyes con la mayor seguridad, al tiempo que con la menor violencia".

\subsubsection{MALTHUS}
La tesis principal de MALTHUS, a saber, que la población tiende a aumentar más deprisa que las existencias de alimentos, no es suya: puede encontrarse en los escritos de otros autores, entre los que figuran ADAM SMITH y BENJAMIN FRANKLIN. Fue, sin embargo, la presentación del problema demográfico de MALTHUS, la que influyó significativamente en el pensamiento económico posterior.
▪
Son tres los factores que explican los orígenes de la teoría de THOMAS MALTHUS:
i)
La presión de la población sobre las existencias de alimentos en Inglaterra: Hasta 1790 aproximadamente, Inglaterra había tenido en gen medida suficientes alimentos, pero a partir de entonces tuvo que importarlos y los precios subieron notablemente.
ii)
Sensación de que las clases de renta más baja eran cada vez más pobres: Inglaterra estaba urbanizándose como consecuencia de la sustitución de la producción doméstica por la producción fabril y parecía que el crecimiento de las ciudades iba acompañado de un aumento del sufrimiento de la clase de renta más baja.
iii)
Discusión con su padre DANIEL MALTHUS: El padre de THOMAS MALTHUS, aceptó las ideas de WILLIAM GODWIN y del MARQUÉS DE CONDORCET, autores utópicos que defendían que el carácter de una persona no se hereda sino que es modelado por el entorno en el que vive. A GODWIN, en particular, le preocupaban las penalidades, el sufrimiento, la infelicidad y el vicio que creía ver en el mundo que lo rodeaba. Llegó a la conclusión de que el principal responsable era el Estado. THOMAS MALTHUS quería demostrar que las ideas que había aceptado su padre eran incorrectas. Trató especialmente de demostrar que la pobreza y el sufrimiento no se debían a las instituciones sociales y políticas y que los cambios de estas instituciones no erradicarían los males de la sociedad. Cuando mostró su ensayo a sus amigos, estos lo animaron a publicarlo. Lo publicó anónimamente en 1978.
▪
En la 1ª edición de su ensayo sobre la población, MALTHUS llegó a la conclusión de que la posposición del matrimonio sólo engendraría vicio, sufrimiento y degradación del carácter, ya que habría relaciones sexuales prematrimoniales. La introducción de cambios en la estructura institucional no erradicaría, pues, el sufrimiento y el vicio de la sociedad mientras los seres humanos
necesitaran alimentos y tuvieran muchos impulsos sexuales. El espectro de la población que presiona implacablemente sobre las existencias de alimentos ha llevado a los observadores a calificar la economía de ciencia lúgubre, expresión utilizada por primera vez por THOMAS CARLYLE (aunque en un contexto distinto, ya que la economía consideraba que todas las razas tenían la misma capacidad de aprender un oficio, mientras que CARLYLE creía que las razas eran diferentes y que era natural que los negros fueran esclavos).\footnote{De acuerdo con JULIAN SIMON la historia es un caso de aumento de la población y aumento del consumo por persona. Sugiere que ambos han aumentado porque el crecimiento tecnológico ha sido sistemáticamente mayor que el crecimiento de la población. SIMON sugiere que el crecimiento tecnológico depende de la gente y que el crecimiento demográfico aumenta el número de personas y, por tanto también aumenta el crecimiento tecnológico. 

Si esta relación es cierta, el crecimiento de la población nunca será excesivo: el crecimiento tecnológico siempre sobrepasará los rendimientos decrecientes. En este caso, la economía es una \textit{ciencia feliz}.}
▪
Esta tesis desencadenó una considerable controversia y despertó el interés por el problema demográfico. MALTHUS, insatisfecho con su creación inicial, publicó una 2ª edición de su ensayo (en 1803), diferente de la 1ª en cuanto a su fin, su metodología, sus argumentos y sus conclusiones. Ya no intentó criticar las ideas de su padre, GODWIN y CONDORCET sino que decidió formular el problema demográfico de una manera tan científica como lo permitieran los datos existentes.
‒
Mientras que la metodología empleada en la 1ª edición era totalmente deductiva, la 2ª edición era algo inductiva y ahora los argumentos iban apoyados con datos estadísticos.
‒
La 2ª edición era, pues, científica tanto por su método como por su fin. Y lo que es más importante, los argumentos y las conclusiones cambiaron. En la 1ª edición, los mecanismos que frenaban el crecimiento de la población engendraban vicio y sufrimiento, pero en la 2ª edición se introdujo un nuevo freno: la contención moral, es decir, la posposición del matrimonio sin actividades sexuales prematrimoniales. Este nuevo freno destruyó en parte los argumentos de MALTHUS contra los utópicos, pero él ya no se preocupó de refutarlos. Su ensayo sobre la población pasó por 7 ediciones sin apenas cambios a partir de la 2ª edición. La edición de la que se dispone generalmente hoy es la 7ª.
▪
La tesis de MALTHUS sobre la población se aplicó en la teoría y la política económicas clásicas. La doctrina del fondo de salarios, desarrollada por SMITH y ampliada por RICARDO y sus seguidores, implicaba que una subida del salario real del trabajo publicaría un aumento de la población, lo cual acabaría bajando el salario y devolviéndolo a su nivel anterior. Se sostenía pues, que cualquier intento de mejorar el bienestar económico de los grupos de renta más baja de la sociedad fracasaría debido a que el volumen de población aumentaría. Por tanto, aunque los sentimientos humanitarios pidieran medidas sociales para elevar la renta de los trabajadores pobres, la lógica económica sólida decía que esos esfuerzos serían vanos. Los intentos de atenuar las penurias económicas de los grupos de renta más baja de la sociedad fracasarían debido a que el volumen de población aumentaría. Por tanto, aunque los sentimientos humanitarios pidieran medidas sociales para elevar la renta de los trabajadores pobres, la lógica económica sólida decía que esos esfuerzos serían vanos. Los intentos de atenuar las penurias económicas de los grupos de renta baja de Inglaterra por medio de leyes comenzaron hacia 1600; los historiadores económicos las denominan leyes de pobres. Los economistas clásicos esgrimieron la doctrina malthusiana de la población para criticar las leyes de pobres. El análisis de los salarios que lograron combinando la tesis malthusiana con la doctrina del fondo de salarios se ha denominado ley de hierro de los salarios. El modelo malthusiano ha tenido muchas más repercusiones de las que nunca imaginó su autor. Los naturalistas británicos CHARLES DARWIN y A. R. WALLACE, que formularon cada uno por separado lo que se conoce con el nombre de teoría darwinista de la evolución, reconocieron ambos que MALTHUS influyó notablemente en su pensamiento.

\subsubsection{J.S. MILL}
Los desarrollos de MILL conllevan las siguientes implicaciones de política económica:
‒
Ley de Say: admite incumplimientos temporales. Entiende la ley de Say como una ecuación y no como una identidad, considerando posibles excesos de oferta temporales. Sin embargo, rechaza las críticas más frontales a la Ley de Say como las formuladas por MALTHUS.
‒
Influencia pero rechazo de un régimen socialista: Es el clásico más preocupado por los aspectos sociales de la economía y aboga por una serie de reformas, pero rechaza un régimen socialista.
‒
Funciones del sector público:\footnote{MILL se vio influenciado por BENTHAM, que también era partidario de una cierta intervención del Estado.} Distingue entre las funciones necesarias y las funciones facultativas:
o
Funciones necesarias: Regulación del trabajo infantil, determinadas prestaciones y las funciones ya propuestas por ADAM SMITH.
o
Funciones facultativas: Estas funciones quedarían sujetas a la decisión del gobierno y deberán existir cuando no limite la iniciativa privada (se asume que pueden existir distintas opiniones).
Dentro de sus preocupaciones sociales también defendió la igualdad de oportunidades, el sufragio femenino y la condena a la esclavitud. Asimismo, defendió la educación como vía de progreso y como vía de evitar los vaticinios de MALTHUS.

Redistribución de la renta vía imposición: MILL aboga por la instauración de un mínimo exento en el impuesto sobre la renta y en un impuesto de sucesiones con progresividad.
‒
Énfasis en el libre comercio.\footnote{Cuando MILL abandonó finalmente la doctrina del fondo de salarios en 1869, el declive de la economía clásica era casi total. Para entonces se habían abandonado la teoría del valor trabajo ricardiana, la doctrina malthusiana de la población y la doctrina del fondo de salarios. CAIRNES trató de salvar el sistema clásico en vano.}

\clearpage
\section{La crítica de \authorfont{KARL MARX}}
KARL MARX es considerado el último clásico,\footnote{Karl Heinrich Marx (Prusia, 1818-Londres, 14 de marzo de 1883), fue un filósofo, economista, sociólogo, historiador, periodista, intelectual y político comunista de origen alemán. En su vasta e influyente obra abarca diferentes campos del pensamiento en la filosofía, la historia, la ciencia política, la sociología y la economía; aunque no limitó su trabajo solamente a la investigación, pues además incursionó en la práctica del periodismo y la política, proponiendo siempre en su pensamiento una unión entre teoría y práctica (praxis). Su interés fundamental se situó en la búsqueda tanto del conocimiento como del desarrollo de las condiciones para la emancipación de la clase trabajadora respecto de la burguesía y de las leyes de acumulación de los capitales encarnados en ella.\\
Junto a ENGELS, es el padre del socialismo científico, comunismo moderno, marxismo y materialismo histórico. Sus obras más conocidas son el \textit{Manifiesto del Partido Comunista} (1848) y \textit{El capital} (1867).\\
MARX es también citado junto a DURKHEIM y a WEBER como uno de los tres principales arquitectos de la ciencia social moderna; a la vez, junto con NIETZSCHE y FREUD, es visto como uno de los tres maestros del siglo XIX de la \textit{escuela de la sospecha} por RICOEUR. También ha sido descrito como una de las figuras más influyentes de historia. En una encuesta de la BBC de 1999, fue votado como el mayor pensador del Milenio por personas de todo el mundo.\\
El siglo que le tocó vivir a MARX, el siglo XIX, supone en Europa el desarrollo de la industrialización. Ello trajo consigo el aumento de la masa de trabajadores en torno a las ciudades y los núcleos industriales, con fuertes desigualdades sociales y largas jornadas laborales.} aunque escribe contra los ellos y acuña el término que se utiliza hoy en día para referirse a los autores anteriores. A pesar de su categorización, muestra importantes diferencias con los clásicos tradiciones: 
\begin{lnum}
    \item No divide la sociedad en trabajadores, capitalistas y terratenientes, sino en dominantes y dominados, a través de una relación de producción, que se identifica con las clases de la burguesía y el proletariado;
    \item Se considera un economista puro y, por tanto, cree que la economía no puede subsistir como ciencia separada de la sociología, la historia, la antropología, etc. Para ello, desarrolla su teoría del materialismo histórico con la que intenta construir una ciencia social unificada.\footnote{Altamente influenciado por HEGEL (considerado por muchos el \textit{último filósofo}, considerado como el más notable del Idealismo alemán y el último de la Modernidad, llamado incluso la conciencia de la Modernidad), que proporcionaría la última obra holística acerca de la sociedad-política y, en definitiva, el mundo en el que vivimos, desde una perspectiva emanentemente cultural y identitaria. MARX trabajó de manera crítica en su obra sobre la base hegeliana, conformando la narrativa histórica de la identidad/cultura de HEGEL como una narrativa basada en las relaciones de producción.\\
    Como veremos, VILFREDO PARETO se verá fuertemente influenciado por la visión conjunta de MARX (no por su ideología, que se encuentra yuxtapuesta) e intentara de la misma forma proporcionar una Teoría holística de las ciencias sociales. Sin embargo, probablemente por su insuficiente brillantez (PARETO), no culminará su obra y su etapa sociológica será condenada al olvido.} Esta teoría señala que lo que mueve la historia es la lucha de clases, que supone un enfrentamiento por el reparto del excedente;
    \item Mientras que los economistas clásicos consideran que la economía se rige por leyes universales que no se pueden dominar, MARX argumenta que las leyes dependen del momento histórico. En concreto, para MARX la evolución de la sociedad y la progresión del sistema capitalista era importante en su explicación de la transición a otro sistema.\footnote{El historicismo de MARX llega hasta el punto de que a pesar de que MARX condena el capitalismo, lo va a defender frente a modos de producción anteriores. El capitalismo será una etapa transitoria, ya que al final (y comenzando por los países capitalistas más avanzados), el proletariado se impondría a las clases opresoras mediante una revolución que establecería una sociedad sin clases y sin Estado. El capitalismo sería abolido ya que era la institución creada por las clases propietarias para ejercer el dominio sobre todos los demás.\\La aproximación histórica tendrá influencia en la escuela institucionalista y en marxistas y postmarxistas, pero en general poca influencia sobre la corriente económica dominante.} 
\end{lnum}

La obra de MARX es de gran relevancia, no tanto por su peso en la corriente académica ortodoxa,\footnote{Aunque da lugar a una escuela propia con métodos propios: la economía marxista, de gran influencia en economistas postkeynesianos que escriben en el siglo XX como ROBINSON o KALECKI.\\
Pensamiento económico neomarxista (BAUER y SWEEZY): ROBINSON presentó a MARX como un buen complemento de las ideas keynesianas, pero ROBINSON no acepta la Teoría del Valor Trabajo. Las colonias de los países en desarrollo no pueden desarrollarse si no se emancipan del orden económico mundial capitalista.} sino por su visión socio-económica, en el que es capaz de entrelazar la sociología, la política y la economía de una manera única. En este sentido, podemos considerar a MARX como uno de los economistas más importantes de todos los tiempos. Sinteticemos sus aportaciones más importantes. 

\subsubsection{Modos de producción}
MARX parte de que la producción es el primer hecho histórico y, como un sistema económico se caracteriza por cómo la sociedad organiza la producción, las condiciones materiales son los principales determinantes de la realidad social. En este sentido introduce el concepto de \textit{modo de producción}, que es un concepto abstracto, compuesto por tres elementos: 
\begin{lnum}
    \item \textit{Relaciones de producción.} Subraya que el sistema económico de su tiempo queda configurado por relaciones de producción entre los capitalistas (burgueses, propietarios del capital) y los trabajadores (proletariado) donde existe un conflicto permanente. MARX habla del fetichismo de la mercancía, en el sentido de que dichas relaciones contractuales son falsas en la medida en que son controladas totalmente por el capitalista como propietario de los medios de producción;
    \item \textit{Estructura económica}. Es el ecosistema que surge como resultado de todas las relaciones de producción en la economía; y,
    \item \textit{Superestructura jurídica y política}. Es el esquema que se configura a partir de la estructura económica y que tiene como fin perpetuar las relaciones de producción. Está representada por el Estado y por la ideología que lo legitima.
\end{lnum}

\imagenfit{MaterialismoHistorico.png}{Esquema del Materialismo Histórico}{(Original) ESCARTÍN, E. \textit{Apuntes Historia del Pensamiento Económic}, 2002; (Extraído) Víctor Gutiérrez, Tema 3.A.2 Repositorio -- TCEE}

\subsubsection{Teoría del valor: Valor-Trabajo y Ley del Valor}\footnote{La teoría del Valor-Trabajo es usualmente asociada a la economía marxista, aunque esta también aparece en las primeras teorías formuladas por los economistas clásicos, como ADAM SMITH o DAVID RICARDO, y posteriormente retomada en la economía anarquista.\\
SMITH describió el precio de una mercancía en términos de trabajo necesario para adquirirlo, que personifica el concepto de cuanto trabajo es requerido para la obtención de un bien y como una herramienta por ejemplo, puede facilitar su adquisición. La teoría del valor-trabajo es un elemento central de la teoría marxista, que sostiene que la clase trabajadora es explotada bajo el capitalismo y disocia valor y precio. MARX, sin embargo, no se refiere a su propia teoría como Teoría del Valor-Trabajo sino como Teoría del Valor.
El resurgimiento en la interpretación de MARX conocida como Neue MARX-LEKTÜRE también rechaza la economía marxista y la teoría del valor-trabajo, llamándola sustancialista. Esta revisión afirma que la teoría del valor-trabajo es una mala interpretación del concepto de fetichismo de la mercancía con relación al valor y que esta interpretación nunca aparece en la obra de MARX. Esta escuela enfatiza en que \textit{El Capital} es explícitamente una crítica a la economía política, en lugar de una teoría más correcta.}

\paragraph{Teoría del Valor-Trabajo: la plusvalía}

Todas las economías producen productos, aquellos bienes que tienen valor de uso, pero en el capitalismo, los productos adquieren la forma de mercancías, que tienen valor de uso y valor de cambio.\footnote{MARX distingue entre producto y mercancía.} Para MARX el valor de cambio de los bienes está determinado por el coste de producción (en linea con los clásicos), pero el único coste de producción es el del trabajo productivo. Para determinar éste (y, por lo tanto, el valor de la mercancía), MARX distingue entre:
\begin{la}
    \item \textit{Trabajo muerto} (o capital constante). Es el trabajo asociado a una máquina, que no produce ninguna plusvalía. En efecto, MARX niega que la maquinaria sea productiva por sí misma, ya que considera que sólo es productiva por la cantidad de trabajo que se utiliza en ella; y,
    \item \textit{Trabajo vivo}. Que representa el valor de cambio ($c$), y que está formado a su vez por:
    \begin{lnum}
        \item \textit{Trabajo concreto} (o capital variable). Lo constituyen los salarios y representa el valor de uso del trabajo ($w$).
        \item \textit{Trabajo abstracto} (o plusvalor). Es la diferencia entre el valor de cambio y el valor de uso del trabajo ($pl=c−w$). 
    \end{lnum}
\end{la}

La idea de MARX es que, como el capitalista domina las relaciones de producción, puede aprovecharse de esto para obtener una ganancia, es decir, remunerar al trabajo por un valor (trabajo concreto) menor al que realmente produce (trabajo concreto + plusvalor), dando lugar a una plusvalía que procede del trabajo exclusivamente, y no del proceso de intercambio. Así pues, la finalidad de la producción es extraer la máxima plusvalía posible de cada trabajador.\footnote{La existencia de la plusvalía permite a MARX resaltar un elemento implícito esencial de la Teoría de la distribución de RICARDO: el conflicto entre los distintos agentes que participan en el proceso productivo. En efecto, MARX resalta los conflictos distributivos, en el sentido en que existen conflictos de intereses entre propietarios del capital y trabajadores.} 

Si MARX define de esta forma en el Tomo I de su obra el valor implícito de una mercancia, en el Toma III aproxima de una manera (menos elaborada) el valor explícito de las mercancias, es decir, el valor de intercambio (el precio), por las razones que ahora se expondrán.

\paragraph{Ley del Valor}
La Teoría del Valor-Trabajo (en su formulación clásica y en su formulación marxista) tenía evidentes limitaciones al trasladar este precio implícito de los bienes al precio explícito observado en el mercado (como es evidente, bienes con las mismas horas de trabajo, si se quiere medir asi, se intercambian en el mercado con precios muy similares en muchas ocasiones)\footnote{MARX mismo introduce esta limitación a priori de la formulación de su Teoría del Valor-Trabajo:

\begin{adjustwidth}{2cm}{0pt}
\raggedright
\textit{[E]stando dado el valor de la fuerza de trabajo y siendo igualmente grande el grado de explotación de la misma, las masas de valor y plusvalor producidas por diversos capitales estarán en razón directa a las magnitudes de las partes variables de esos capitales, esto es, a sus partes invertidas en fuerza de trabajo viva. Esta ley contradice abiertamente toda la experiencia fundada en las apariencias. Todo el mundo sabe que el dueño de una hilandería de algodón que, si nos atenemos a los porcentajes del capital total empleado, utiliza proporcionalmente mucho capital constante y poco capital variable, no por ello obtiene una ganancia o plusvalor menor que un panadero, quien comparativamente pone en movimiento mucho capital variable y poco capital constante. Para resolver esta contradicción aparente se requieren aún muchos eslabones intermedios.}\\
\raggedleft{\authorfont{KARL MARX}, “\textit{El capital. Tomo I}”, 1867}
\end{adjustwidth}\\
Su colaborador, ENGELS, le advirtió de esta equiparación (Teoría del Valor-Trabajo) como una objeción que debió haber anticipado, pero MARX prefirió no responderla hasta más adelante porque \textit{echaría a perder todo el método dialéctico de exposición}, el cual \textit{constantemente tiende trampas a los filisteos y los economistas vulgares}.\\
De hecho, SMITH, como ya se ha visto, era consciente de la aparente incompatibilidad que presentaba la igualación de la tasa de ganancia entre sectores con la teoría del valor-trabajo, por lo que, de manera indirecta recurrió a relegar esta Teoría del Valor-trabajo a un \textit{estado rudo y primitivo de la sociedad} optando, para la valoración explícita de los bienes, por una teoría del valor como costo de producción (para sociedades capitalistas modernas), donde el valor de las mercancías se determina por la suma de los salarios, beneficios económicos y la renta de la tierra que se destinaban a la producción. Sin embargo, esta explicación es vulnerable al cargo de tautología en el sentido de que explica los precios (de productos) mediante precios (de los insumos), como ya hemos visto; por eso RICARDO intentó minimizar esta inconsistencia argumentando que generalmente la Teoría del Valor-Trabajo era una buena aproximación (de hecho, alrededor del 93\% de las veces, considerando trabajo homogeneo) bajo ciertas condiciones promedias.}. Entonces,\footnote{Ilustramos recurriendo a una obra previa, pero en el Tomo III, como se ha dicho, formaliza esta posición} 

\begin{adjustwidth}{2cm}{0pt}
\raggedright
\textit{[E]l precio se distingue del valor no sólo como lo nominal de lo real; [...] sino porque esta última aparece como la ley de los movimientos que atraviesa la primera. Pero los dos son constantemente diferentes y nunca se equilibran, o sólo se equilibran de manera coincidente y excepcional. El precio de un bien se mantiene constantemente por encima o por debajo del valor de la mercancía, y el valor de la mercancía misma existe sólo en este movimiento hacia arriba y hacia abajo de los precios de la mercancía. La oferta y la demanda determinan constantemente los precios de las mercancías; nunca se equilibran, o sólo por coincidencia; pero el costo de producción, por su parte, determina las oscilaciones de la oferta y la demanda}\\
\raggedleft{\authorfont{KARL MARX}, “\textit{Cuaderno I – Dinero}”, 1857}
\end{adjustwidth}

Por tanto, como vemos, la solución para MARX está en la competencia y la reasignación de capital después de la producción: dado que el beneficio se genera únicamente con el trabajo, una empresa con un uso muy intensivo de capital genera una tasa de ganancia baja, mientras que una empresa con un uso muy intensivo de mano de obra genera una tasa de ganancia más alta; por tanto, los capitalistas de la empresa con una tasa de ganancia baja reasignarán su capital hacia la empresa con una tasa de ganancia alta, la reasignación de capital inducirá una modificación de la distribución de beneficios hasta que las tasas de beneficio de todas las empresas sean todas iguales (hasta aquí bien).\footnote{De hecho, de esta bonita aproximación a los flujos de capital nace una de las ideas más trasgresoras de la ideología marxista, a saber: la distinción subyacente entre el Dinero-Cambio y Dinero-Mercancia (capital). Estos ciclos de $M\to D\to M$, mercancías-dinero-mercancías, recibe el nombre de producción comercial capitalista simple.} Por lo tanto, MARX define una Ley de Transformación o Ley de Valor (sólo aplicable de manera agregada) en la que dice cumplirse lo siguiente:

$$C+Pl+V=Y\quad \text{donde,}
\begin{cases}
    C\equiv\quad \text{Capital Fijo Total}\\[0.5em]
    Pl\equiv\quad \text{Plusvalía total}\\[0.5em]
    C\equiv\quad \text{Capital variable total (salarios, }w)\\[0.5em]
    Y\equiv\quad \text{Producción total (en valor de mercado)}
\end{cases}$$

Además, concluye ofreciendo cuatro argumentos (aproximado, no expplicito) para explicar que esta igualación se cumple en términos agregados y que la desviación observada entre productos particulares (observados de manera empírica), no pone en duda la Teoría del Valor-Trabajo como fundamento del Valor final de mercado:
\begin{lnum}
    \item El valor medio de los bienes de una sociedad ($C+V+Pl$) es igual a la media de los precios ($Y$), por lo que los precios de los bienes giran en torno al valor de los bienes;
    \item El tiempo de trabajo necesario para la producción de un bien sigue influyendo en el precio del bien: un aumento de este tiempo de trabajo aumenta el precio, y una reducción del tiempo de trabajo hace que baje;
    \item Históricamente, antes del modo de producción capitalista, el valor de un bien era efectivamente igual a su precio correspondiente, sólo con la posibilidad de reasignación de capital apareció esta desviación, el valor de un bien es en esencia, por lo tanto, valor trabajo; y,
    \item Los valores se hallan detrás de los precios de producción y, en última instancia, los determinan, es decir que los precios de los bienes sólo se transforman, estos precios son siempre aproximadamente iguales a sus respectivos valores.
\end{lnum}

\paragraph{Problema de la transformación}
Se debe a EUGEN VON BÖHM-BAWERK, la primera formulación (originaria, no matemática) de lo que hoy se conoce como el problema de la transformación:\footnote{\href{https://es.wikipedia.org/wiki/Problema_de_la_transformación}{\textit{Problema de la transformación}}.}

\begin{adjustwidth}{2cm}{0pt}
\raggedright
\textit{Se declaró que el valor era "el factor común que aparece en la relación de intercambio de mercancías" (i. 13). Se nos dijo, en la forma y con el énfasis de una estricta conclusión silogística, sin excepción, que establecer dos mercancías como equivalentes en el intercambio implicaba que existía "un factor común de la misma magnitud" en ambas, al que cada uno de los dos "debe ser reducible" (i. 11). [...] Y ahora, en el tercer volumen, se nos dice breve y secamente que lo que, según la enseñanza del primer volumen, debe ser, no es y nunca puede ser; que las mercancías individuales intercambian y deben intercambiarse entre sí en una proporción diferente de la del trabajo incorporado en ellas, y esto no accidental y temporalmente, sino por necesidad y permanentemente. No puedo evitarlo; no veo aquí ninguna explicación y reconciliación de una contradicción, sino la contradicción desnuda en sí misma. El tercer tomo de Marx contradice al primero. La teoría de la tasa media de ganancia y de los precios de producción no puede conciliarse con la teoría del valor. Esta es la impresión que, creo, debe ser recibida por todo pensador lógico. Y parece haber sido aceptado de manera muy general".}\\
\raggedleft{\authorfont{EUGEN VON BÖHM-BAWERK}, “\textit{La conclusión del sistema marxiano}”, 1896}
\end{adjustwidth}

\subsubsection{Leyes de movimiento capitalista}
MARX se separó notablemente de la economía clásica al subrayar el cambio tecnológico (en terminología clásica \textit{los rendimientos del capital} y la libre competencia, por medio del comportamiento egoísta de los agentes económicos\footnote{SMITH interpretó el progreso en términos de comportamiento humano racional más que en términos de avance técnico. RICARDO tenía una experiencia industrial  muy limitada; por lo que no tuvo  intención de refundir la economía política como una teoría del cambio tecnológico, de hecho, vio el problema económico de la sociedad como un problema agrícola (influenciado probablemente por MALTHUS) que, indudablemente tendía a un estado estacionaria en el mejor de los casos. MILL estuvo más abierto a las perspectivas del cambio tecnológico pero no le otorgó el papel central en su teoría como hizo Marx en la suya.}) como la fuerza motriz de su dinámica social. Para ello describió cinco leyes, o tendencias generales inherentes al capitalismo:\footnote{Cada una de ellas brota de la naturaleza dinámica de la economía y tiene su raíz en el conflicto entre las fuerzas productivas dinámicas y las relaciones de producción estáticas.} 
\begin{lnum}
    \item Ley de acumulación y la tasa decreciente de ganancia. Los capitalistas tratan de obtener las máximas plusvalías posibles y, como éstas proceden del trabajo, tratarán de explotarlo lo máximo posible. En estas circunstancias, las plusvalías pueden ser decrecientes debido a: (i) mayores presiones salariales; y/o, (ii) excesiva acumulación de capital;\footnote{Si un capitalista individual instala un mayor nivel de mecanización, su plusvalía aumenta porque produce a un menor coste (sustituye trabajo por capital), pero puede seguir vendiendo al mismo precio que el resto de empresas, menos mecanizadas. Sin embargo, si otros capitalistas le copian y este proceso se acaba extendiendo, entonces aumentará la competencia y el precio se reducirá y reducirá sus beneficios (y su extracción plusvalía)}
    \item \textit{Ley de concentración}. Como las plusvalías proceden del trabajo, los capitalistas tenderán a crear empresas de gran escala. Es decir, la sustitución del trabajo por capital transforma la industria de pequeña escala en gran escala con una capacidad de producción mayor. Esto dará lugar a sobreproducción, reduciendo los precios hasta el punto en que sólo los productores más eficientes sobrevivan. Así, la industria se irá centralizando progresivamente y el poder económico se concentrará en un número de manos cada vez menor:
    \item \textit{Ley del creciente ejército industrial de reserva}. El capital sustituye al trabajo, lo que lleva a la aparición de paro, dando lugar a un creciente ejército industrial de desempleados;\footnote{Influencia del ludismo, las tecnologías ahorradoras de trabajo generan aumentos de paro.} y,
    \item \textit{Ley de la miseria creciente del proletariado}. Conforme crece el ejército de reserva, los capitalistas compensan parte de la caída de los beneficios reduciendo los salarios, aumentando las jornadas de trabajo, incorporando a mujeres y niños y, en resumen, crece la miseria del proletariado.
    \item \textit{Ley de la crisis}. La miseria creciente del proletariado tiende a la revolución social, esto es, a que el proletariado acabe apoderándose de los medios de producción.
\end{lnum}

Por tanto, la implicación principal de las leyes del movimiento capitalista es la desaparición del capitalismo. Es decir, el capitalismo es un momento histórico que tiene fecha de caducidad. En efecto, aunque los capitalistas aprovechan inicialmente su carácter de clase dominante para apropiarse de la plusvalía de los trabajadores, el proceso de empobrecimiento de éstos acabará desembocando en una crisis que dé lugar a un nuevo modo de producción: el socialismo, donde los medios de producción están en manos de los trabajadores.

MARX nunca se preocupó demasiado en articular una teoría de cómo debería funcionar el nievo sistema económico que sucediera a lo que llamó capitalismo. Su visión era que, bajo la nueva organización social creada por la victoria del proletariado, la productividad sería tal que los problemas de escasez desaparecerían y con ello el objeto de la teoría económica.\footnote{
“\textit{No somos lo suficientemente civilizados para pasar directamente al socialismo, a pesar de que las políticas tienen sus primeros frutos.}”, declaró LENIN. Se refería al hecho de que Rusia seguía siendo una nación principalmente agraria, con una muy baja población urbana y una base industrial débil, además de no poseer el criterio económico necesario para un socialismo completo. }

\subsubsection{Política económica}
En su obra \textit{El Manifiesto Comunista} (1848), MARX expone \textit{a cada uno sus necesidades y de cada uno según sus posibilidades.}, para ello recomienda:\footnote{
¿Quién decide cuánto puede trabajar?¿Quién decide las necesidades de las personas?
La competencia es una gran disciplina contra las desigualdades. Una empresa que gane excesivamente verá cómo le nacen competidores. Un trabajador especializado que cobre grandes salario verá como otros intentan hacer lo mismo que él.
Otra posibilidad es que el Estado puede corregir desigualdades sociales mediante políticas de igualdad de oportunidades. La Economía moderna dice, por ejemplo, que las políticas de renta son, en este sentido, mucho mejores que las políticas de precios o de cuotas.
Los países que han aprovechado estas ventajas son los que han conseguido las sociedades más igualitarias que ha conocido la historia.}
\begin{lnum}
    \item Planificación de la producción;
    \item Expropiación de la tierra;
    \item Abolición del derecho de herencia;
    \item Obligación a trabajar; y,
    \item Educación pública y gratuita.
\end{lnum}

\begin{specialblock}{Las críticas históricas al capitalismo}
    La utopía económica de MARX no es, ni mucho menos, propia. El dilema existente entre la clase obrera empobrecida y la acumulación de riqueza inherente al sistema capitalista ya fue tratada por SMITH, MILL, y por grandes teólogos. (Todo lo que continua es cita literal de: ESCARTÍN GONZÁLEZ, \textit{Historia del Pensamiento Económico, Tema 12 - SMITH})

    Smith fue ante todo un moralista. Continuó con la tradición inglesa de la filosofía moral establecida por SHAFTESBURY y HUTCHESON. Madurando y desarrollando las ideas de ambos autores, llegó a conclusiones propias que transmitió con sencillez e inteligibilidad. Publicada la \textit{Teoría de los sentimientos morales} (1759), sus lectores se identificaron fácil e intuitivamente con los principios expuestos por Smith y que conservó en su \textit{Riqueza de las naciones} (1776). [S]iguiendo ideas que se remontan a SÉNECA y pasan por SAN JUAN CRISÓSTOMO, opinaba que el trabajo realizado en provecho propio colaboraba en cierto modo en satisfacer las necesidades de los demás; el rico, pese a su egoísmo para acumular riquezas, inevitablemente favorece a los pobres, a los que tiene que dar empleo para satisfacer sus propios deseos. 

    \begin{adjustwidth}{2cm}{0pt}
    \raggedright
    \textit{Los ricos escogen del montón sólo lo más preciado y agradable. Consumen poco más que el pobre, y a pesar de su egoísmo y rapacidad natural, y aunque sólo procuren su propia conveniencia, lo único que se proponen con el trabajo de esos miles de hombres a los que dan empleo es la satisfacción de sus vanos e insaciables deseos, dividen con el pobre el producto de todos sus progresos. Son conducidos por una mano invisible que los hace distribuir las cosas necesarias de la vida casi de la misma manera que habrían sido distribuidas si la tierra hubiera estado repartida en partes iguales entre todos sus habitantes; y así, sin proponérselo, sin saberlo, promueven el interés de la sociedad y proporcionan medios para la multiplicación de la especie.}\\
    \raggedleft{\authorfont{ADAM SMITH}, “\textit{Teoría de los sentimientos morales}”, 1759}
    \end{adjustwidth}

    [Sin embargo, para SMITH], la principal institución económica que no actuaba con la rectitud de conducta necesaria para lograr el interés general de la sociedad, puesto que violaba los principios de equidad y de justicia, era el monopolio. El monopolio, o mejor dicho, lo que hoy en día llamamos oligopolio, en la época de Smith ya estaba bastante extendido y era la tendencia a la que aspiraban los hombres de negocio. Smith, con su crítica, intentó la erradicación de las prácticas desleales en la economía. [Y su animadversión recurrente, fruto de su época, a los Aristócratas ociosos que reciben rentas de la tierra por "derecho natural"]

    Tan lapidarias como las anteriores sentencias son las siguientes: 

    \begin{adjustwidth}{2cm}{0pt}
    \raggedright
    \textit{(...) los comerciantes y manufactureros, con sus protestas y razonamientos capciosos, les convencen fácilmente de que el interés privado de una parte de la sociedad coincide con el general de toda ella\\
    “Los comerciantes y los fabricantes [...] como su inteligencia se ejercita por regla general en los particulares intereses de sus negocios específicos, más bien que en los generales de la sociedad, su dictamen, aun cuando responda a la mejor buena fe (cosa que no siempre ha ocurrido) se inclina con mayor fuerza a favor del primero de esos objetivos que del segundo}\\
    \raggedleft{\authorfont{ADAM SMITH}, “\textit{Teoría de los sentimientos morales}”, 1759}
    \end{adjustwidth}
    
    Como se aprecia, el primer crítico de esta teoría fue su propio creador, aunque en puridad de términos, más bien deberíamos decir que SMITH no se critica a sí  mismo sino que se traslada del plano de la realidad al plano de la fantasía.

    [MARX], el radical, se enfrentó a un dilema similar al de SMITH, pero su propuesta fue mucho más drástica. No se contentó con transformar alguna faceta del capitalismo, como la sustitución del monopolio y oligopolio por la libre competencia. Su idea consistió en desmoronar directamente toda la estructura del sistema capitalista; para ello sólo bastaba minar lo que era considerado un pilar fundamental: la propiedad privada (a diferenci de SMITH\footnote{Smith puso todas sus esperanzas en el naciente capitalismo industrial para lograr el progreso de la sociedad. Consideró que es consustancial de los pueblos civilizados la institución de la propiedad privada y, por consiguiente, la apropiación de las tierras, del capital y del propio trabajo; de ahí que el  precio de los bienes tenga que resolverse en tres partes para distribuirse entre cada uno de los propietarios del respectivo factor productivo\\
    SMITH dice que cuando en una sociedad no existe la acumulación de capital ni la apropiación de la tierra, se trata de sociedades rudas y primitivas.}). 
    
    Mas se equivocó, porque ésta, en realidad, era un falso pilar, pues no pertenecía a la estructura, sino a la superestructura jurídica. El verdadero pilar son los beneficios, con independencia del régimen de propiedad. La acumulación de capital, la industrialización, el desarrollo económico, en suma, el progreso social descansa en el beneficio. Y éste puede obtenerse tanto con propiedad privada como con la pública, y lo mismo ocurre con la decisión de destinarlo a la acumulación de capital o al consumo.\\
    
    La idea más subversiva contra el capitalismo se encuentra en la \textit{Utopía} de TOMÁS MORO. Este autor sí supo poner la dinamita en un pilar básico de la estructura del sistema capitalista: el dinero. Al suprimir el dinero, eliminaba la forma más eficaz de obtener el beneficio y de acumular el capital. No es de extrañar, por consiguiente, que su sociedad ideal fuera estacionaria; MORO en ningún momento describe una verdadera dinámica que entrañe cambios en la sociedad de los utopienses.
    
\end{specialblock}

\textbf{NOTA.-}  

En el Epígrafe 3 se han transcrito frases que
avalan esta poca confianza de Smith en la consecución
de los intereses generales mediante la actuación egoísta
de empresarios, manufactureros y comerciantes. Tan
lapidarias como las anteriores sentencias son las
siguientes: “los comerciantes y manufactureros, con sus
protestas y razonamientos capciosos, les convencen
fácilmente de que el interés privado de una parte de la
sociedad coincide con el general de toda ella” (Smith,
1776, p. 124). “Los comerciantes y los fabricantes [...]
como su inteligencia se ejercita por regla general en los
particulares intereses de sus negocios específicos, más
bien que en los generales de la sociedad, su dictamen,
aun cuando responda a la mejor buena fe (cosa que no
siempre ha ocurrido) se inclina con mayor fuerza a
favor del primero de esos objetivos que del segundo”
(ib., p. 240).

\clearpage
\section{Conclusión}

\subsection{Recapitulación}



\subsection{Extensiones y relación con otras partes del Temario}



\subsection{Opinión}

\subsubsection{Idea final (Salida o cierra)}

\clearpage
\pagenumbering{gobble}
\MyBibliography
\MyBibItem{Autor}{Título}{Año}{DOI -- LINK}


% --- Anexos (no aparecen en el índice) ---
\clearpage
\anexo{\textbf{Preguntas Test}}
%\input{\TemaCode/questions.tex}

\clearpage
\anexo{La economía preclásica}\label{anx:preclasicos}

Como el pensamiento económico antes del siglo XVIII ha resultado muy influyente en la elaboración de la ciencia económica, sería perjudicial no decir nada sobre el pensamiento económico anterior.
Si no elegimos el siglo XVIII como origen de la economía, de nuevo se nos presenta el problema del punto de partida, ¿cuándo comienza el ser humano a pensar en temas económicos?
Para analizar esta primera etapa del pensamiento económico, seguiremos a LANDRETH y COLANDER (2010) y dividiremos este periodo en dos partes:
▪
Comienzos del periodo preclásico (800 a.C–1500 d.C.): Como abarcan 2.300 años (alrededor de 12 veces los casi 250 años que han transcurrido desde 1776) hay que hacer una gran selección. Dividiremos los comienzos del período preclásico en cuatro subperíodos:
‒
Los inicios del pensamiento económico oriental: GUAN ZHONG (siglo VII a.C.)
‒
El pensamiento griego: HESÍODO, JENOFONTE, PLATÓN y ARISTÓTELES (siglos VIII – IV a.C.).
‒
El pensamiento árabe-islámico: AL-GHAZALI y IBD KHALDUN (siglos XI – XV d.C.).
‒
El pensamiento económico de los escolásticos: SANTO TOMÁS DE AQUINO (siglo XIII d.C.)
▪
Era preclásica (1500–1776): El período comprendido entre 1600 y 1750 se caracteriza por el aumento de la actividad económica. El feudalismo, con sus feudos autosuficientes desde el punto de vista económico, social y político, estaba dejando paso a un creciente comercio, al auge de las ciudades fuera del feudo y a la expansión del estado-nación. La actividad individual estaba menos controlada por la costumbre y la tradición de la sociedad feudal y por la autoridad de la Iglesia. La producción de bienes para el mercado estaba cobrando más importancia y la tierra, el trabajo y el capital comenzaban a comprarse en los mercados. Estas tendencias prepararon el terreno para la Revolución Industrial.

Durante este período el pensamiento económico dejó de ser una simple aplicación de ideas sobre los individuos, los hogares y los productores para convertirse en una visión más compleja de la economía como un sistema que tenía sus propias leyes e interrelaciones. Dividimos nuestro análisis de este período en tres grandes apartados:
‒
Mercantilismo.
‒
Precursores del pensamiento económico clásico.
‒
Fisiocracia

\textsc{Comienzos del periodo preclásico}

\underline{El pensamiento económico oriental}\\
Los historiadores occidentales del pensamiento económico tienen tendencia a concentrarse en los autores occidentales. Por ejemplo, J.A. SCHUMPETER (1954) sostenía que no había podido encontrar ningún primer escrito económico no occidental que tuviera contenido analítico y afirmó que “no ha llegado hasta nosotros ningún análisis de temas estrictamente económicos que pueda llamarse ‘científico’ en el sentido en que utilizamos nosotros el término”. SCHUMPETER también señaló el curioso vació que existía, a su juicio, en la literatura económica entre los escritos de los griegos y Santo Tomás de Aquino, un periodo de casi mil años durante el cual parecía que no se había escrito de economía nada que mereciera la pena.
Desde SCHUMPETER, algunos estudiosos han puesto en cuestión sus conclusiones y han comenzado a encontrar algunos escritos interesantes sobre economía. Nosotros examinaremos brevemente las obras de GUAN ZHONG, autor chino del siglo VII a.C., lo que nos valdrá para demostrar que es probable que se hayan realizado análisis de la actividad económica en diversas épocas y lugares. La obra de GUAN ZHONG, «Guan Zi» (traducido literalmente quiere decir «Los escritos del maestro Guan»), contiene algunas ideas fundamentales para el pensamiento económico:

Precursor de la ley de la oferta y la demanda: GUAN ZHONG sostenía que cuando un bien era abundante, se volvía ligero y su precio bajaba. Cuando “se guardaba bajo llave”, se volvía pesado y su precio subía. Había movimientos de entrada y salida de bienes de los mercados dependiendo de su peso y una tendencia clara hacia un único precio: el equilibrio. La teoría de lo ligero y lo pesado es, pues, una formulación de la ley de la oferta y la demanda.

Precursor de la teoría cuantitativa del dinero: GUAN ZHONG también aplicó esta teoría de lo ligero y lo pesado al dinero para desarrollar una teoría cuantitativa del dinero, afirmando que cuando el dinero era pesado, su precio subía (los precios de los bienes bajaban) y cuando era ligero, su precio bajaba (los precios de los bienes subían).

Precursor de la política fiscal anticíclica: Afirmó que para detener las fluctuaciones generadas por los cambios en los precios, aconsejaba que el estado comprara bienes cuando el dinero fuera pesado (manteniendo así alto el nivel de precios) y vendiera cuando fuera ligero (manteniendo así bajo el nivel de precios). Eso no sólo contribuiría a estabilizar el nivel de precios sino que también generaría ingresos para el Estado.

\underline{El pensamiento griego}
Cabría pensar que el pensamiento económico griego ha sido analizado exhaustivamente y que los estudiosos están plenamente de acuerdo sobre la importancia relativa de los diversos autores, pero no es así. No hay una idea clara sobre que autores han sido más relevantes. Aquí mencionaremos a cuatro: HESÍODO, JENOFONTE, PLATÓN y ARISTÓTELES.
HESÍODO (alrededor del siglo VIII a.C.)
La tradición admitida a día de hoy es considerar que el poeta griego HESÍODO es el primer autor occidental en hablar de economía. Según HESÍODO, la escasez no se debe a que los recursos son limitados y los deseos del hombre son ilimitados sino que es uno de los males que salieron de la caja de Pandora cuando esta la abrió57. HESÍODO pasó su infancia y adolescencia dedicándose a la agricultura y al pastoreo, junto a su hermano PERSES, con quien acabó enfrentado debido a la herencia que recibieron al fallecer su padre. Al parecer, PERSES había dilapidado su parte rápidamente y entabló un pleito judicial con HESÍODO. El tribunal de justicia dio la razón a PERSES y Hesíodo tuvo que darle parte de lo que había recibido como herencia. Con posterioridad, Perses volvió a quedar en situación económica precaria e intentó recurrir a la ayuda de su hermano al que incluso volvió a amenazar con ir a los tribunales, pero este rehusó ayudarlo. En su lugar, para guiarlo por un buen camino y que llevara a cabo una buena gestión, HESÍODO compone el poema «Trabajos y días». Este poema gira en torno a dos verdades generales: el trabajo es el fundamento de la vida en sociedad y solamente aquel que trabaja tendrá la posibilidad disfrutar de una vida decente y feliz.
JENOFONTE (siglos V–IV a.C.)
Discípulo de Sócrates, escribe en forma de diálogos socráticos como haría PLATÓN. A nivel económico escribe un libro de consejos prácticos para los responsables de la gestión de un dominio agrícola y en él trata temas referidos a la gestión eficiente del hogar. Titula este libro «Οίκονομικός» (Oeconomicus) que en griego quiere decir «la ley (νόμος) de la casa (Οἶκος)». JENOFONTE, que escribió 4 siglos después que HESÍODO, llevaría los conceptos de gestión eficiente mucho más allá que HESÍODO y los aplicaría al hogar, al productor, al ejército y a la administración pública. Esto le permitió comprender que es posible mejorar la eficiencia practicando la división del trabajo. Esto tendría influencia en el pensamiento económico posterior, desde otros griegos como ARISTÓTELES, hasta los escolásticos y finalmente ADAM SMITH, quien reconoció especialmente la influencia de la eficiencia de la economía y de la sociedad en «La Riqueza de las Naciones».
A partir de este momento, las ideas económicas aparecerán progresivamente de forma más estructurada, en el ámbito del pensamiento antiguo griego, es decir, entorno a la filosofía. Se encuentran trazos sobre temas económicos en los grandes filósofos de la época. Aquí nos centraremos en dos de los más remarcables, PLATÓN y ARISTÓTELES.
PLATÓN (siglos V–IV a.C.)
También discípulo de SÓCRATES, escribe en una época en la que los griegos dudan de la pertinencia de sus instituciones. El hombre griego, que normalmente da todo por la ciudad y construye su vida para servirla, cada vez acepta menos la organización institucional de las ciudades. Uno de los motivos de estas tensiones es la existencia de desigualdades sociales, que crean revueltas y saqueos continuos a menudo seguidos del advenimiento de tiranos. Este contexto no es favorable a la paz civil ni al crecimiento económico.
Esta situación conduce a PLATÓN a imaginar una sociedad ideal y pacífica en dos libros: «Las Leyes» y «La República». A menudo presentado como un precursor de comunismo, PLATÓN introduce en las relaciones económicas la lógica de apego patriótico a la ciudad, que es el fundamento de la vida política griega. De sus textos se deducen dos ideas claves:
▪
La voluntad de una limitación de la población58: Una sociedad armoniosa supone pocos hombres. Esto es así por motivos políticos (más fácil tomar decisiones) y económicos (si la ciudad es más grande, tiene más habitantes con más necesidades, por lo que es necesario aprovisionarse más lejos y eso tiene un coste).
▪
La ciudad determina todos los parámetros de la organización social y, en particular, las condiciones que rigen la circulación de la moneda: Para PLATÓN la moneda es un nomisma, es decir, el fruto de una convención social. Un objeto deviene monetario porque el poder político así lo quiere y su valor está garantizado por ese poder. De este modo, PLATÓN inicia un debate sobre la relación profunda entre la moneda y el poder político.
ARISTÓTELES (siglo IV a.C.)
Discípulo de PLATÓN, su relevancia trasciende sus aportaciones al pensamiento filosófico debido a su influencia en las ideas económicas durante el periodo del escolasticismo.
ARISTÓTELES aborda 5 problemas:
1.
El origen del precio y el valor de los objetos: ARISTÓTELES considera que el valor de un bien no viene de sí mismo. Identifica un valor de uso (correspondiente al servicio prestado) y un valor de cambio (formado en los mercados). El valor de cambio son el centro del problema económico. ¿En qué medida es legítimo que el precio (medida concreta del valor de cambio) se desvíe del valor de uso? Introduce así una noción que será posteriormente desarrollada en la Edad Media, el problema del «precio justo». Para que el precio de mercado sea justo ARISTÓTELES identifica dos condiciones: el comprador debe tomar su decisión de forma libre y el vendedor no debe estar en situación de poder en la determinación del precio.
2.
La naturaleza de la moneda: Si bien ARISTÓTELES está de acuerdo con PLATÓN en que el poder liberatorio de la moneda depende de una convención social, no considera que esta sea su única razón de ser. ARISTÓTELES considera que la moneda provee de un servicio independiente a la voluntad del Estado y que por tanto, tiene una existencia autónoma. Concretamente, considera que ayuda a resolver el problema de las doble coincidencia de necesidades, ya que en una economía de trueque es necesario que comprador y vendedor necesiten aquello que les ofrezca la contraparte, pero en una economía monetaria esto no es necesario debido al poder de compra universal de la moneda que sirve como medio de cambio. Como consecuencia, independientemente del Estado, el mundo del comercio es libre de elegir un objeto material y utilizarlo como medio de cambio.
3.
Condena del préstamo a interés: ARISTÓTELES parte de la idea de que el dinero no crea dinero. Por lo tanto, aquel que pude prestado lo hace en general en situaciones de urgencia para solucionar compromisos previos tomados a la ligera. ARISTÓTELES concibe el préstamo como un mecanismo que conduce inexorablemente a la ruina de los prestatarios y a la pérdida de su libertad. Este pensamiento influirá en la condena del préstamo de los pensadores religiosos de la Edad Media. La finanza islámica continúa funcionando bajo esta condena y está obligado a inventar pretextos para eludirlo.
4.
Condena de la búsqueda del beneficio monetario: Según ARISTÓTELES, las necesidades de los individuos son moderadas, pero sus deseos son ilimitados. De ahí que la producción de mercancías para satisfacer las necesidades sea correcta y natural, mientras que la producción de bienes para intentar satisfacer los deseos ilimitados no sea natural. ARISTÓTELES admitía que cuando se producen bienes para venderlos en un mercado, puede ser difícil saber si esta actividad satisface necesidades o deseos desmesurados; pero suponía que si un intercambio de mercado no se realiza mediante trueque, se realiza para satisfacer necesidades naturales y no se pretende obtener ningún beneficio económico. Sin embargo, la utilización del dinero induce a pensar que el objetivo del intercambio es un beneficio monetario, que es algo que ARISTÓTELES condenaba. Una de las observaciones más interesantes de ARISTÓTELES es que el problema de la escasez puede resolverse reduciendo el consumo, cambiando las actitudes humanas. Ésta es una poderosa idea para los utópicos y los socialistas que confían en poner término a los conflictos sociales eliminando los conflictos inherentes a la escasez.
5.
Defensa de la propiedad privada: PLATÓN había afirmado que la clase gobernante de su sociedad ideal, los soldados y los filósofos, no debía poseer propiedad privada sino propiedad comunitaria para evitar conflictos por la propiedad que pudieran desviar su atención de cuestiones más importantes. Sin embargo, ARISTÓTELES creía que la propiedad privada cumplía una función útil en la sociedad y que no debía tomarse ninguna medida para limitar la cantidad de propiedad privada. La incoherencia en la que incurrió aparentemente al condenar la búsqueda del beneficio económico y defender al mismo tiempo el derecho a la propiedad privada preocupó a los filósofos morales hasta el siglo XVI.
ARISTÓTELES coincide con PLATÓN y con casi todos los demás pensadores griegos en la necesidad de ver la actividad económica desde una perspectiva más amplia y no compartimentar el estudio.

\underline{El pensamiento árabe-islámico}\\
Los estudiosos árabe-islámicos fueron mucho más que meros traductores del pensamiento griego. Se sabía perfectamente que las ideas griegas se tradujeron del árabe al latín para uso de los escolásticos, no del griego. Ahora está reconociéndose que en muchas disciplinas los árabes hicieron importantes aportaciones propias.
Tanto los autores árabe-islámicos, como los escolásticos que los siguieron, escribieron en un entorno muy diferente al nuestro. Los economistas modernos abstraen las actividades económicas de la totalidad de la vida humana, lo cual quizá sea correcto en el siglo XXI para analizar las complejas economías de los países desarrollados, en las que las actividades económicas son extraordinariamente importantes. Sin embargo, los autores árabe-islámicos consideraban todos los aspectos de la actividad humana y especialmente las consecuencias de esta actividad –de la cual la actividad económica era una pequeña parte– para la salvación del individuo. No se realizaba un análisis económico formal independiente como hoy; los estudiosos islámicos medievales examinaban, por el contrario, las cuestiones económicas en el contexto más amplio de sus ideas religiosas.
Como se consideraba que todas las actividades humanas estaban interrelacionadas y sometidas al gobierno de la ley divina, era difícil formular un modelo económico analítico. Sólo se avanzó en el conocimiento de la economía cuando los estudiosos decidieron examinar, no el mundo musulmán ideal, con sus numerosas limitaciones de la actividad económica, sino el mundo musulmán real de su tiempo. Uno de los primeros temas de interés de estos intentos de estudiar la actividad económica fue la tributación.

\underline{El pensamiento económico de los escolásticos}\\
Los pensadores chinos, griegos, árabe-islámicos y escolásticos no analizaron la economía como una disciplina independiente; estaban interesados en cuestiones mucho más amplias y filosóficas.
A mediados del siglo XV, las ideas escolásticas sobre la vida virtuosa chocaban con la práctica económica vigente y los juicios éticos de la Iglesia parecían fuera de lugar en las economías en desarrollo de Europa occidental. No obstante, la doctrina escolástica sí aportó ideas sobre el funcionamiento de la creciente economía de mercado y contribuyó a sentar las bases para el desarrollo de un enfoque más analítico.
Tuvieron que ocurrir algunas cosas antes de que la economía de mercado pudiera desarrollarse plenamente y generar la enorme oleada de bienes inherente a los recursos naturales que había para utilizar y los conocimientos y la tecnología que había para explotar. Uno de los cambios más cruciales fue la gran transformación de la estructura institucional de Europa occidental. La libertad fue el elemento clave en este cambio: liberación de las ataduras de la fría mano de la tradición que ahogaba en cambio, liberación de la ideología de las enseñanzas religiosas que no veían con buenos ojos la actividad económica, liberación del poder político y económico de la Iglesia que se oponía a la aparición de nuevos intereses económicos y liberación del gobierno que creaba y apoyaba el monopolio y se dedicaba a otras actividades que retrasaban el avance económico. Retrospectivamente, la dotrina escolástica representa un lento repliegue en favor de una aceptación mayor de las actividades económicas. La economía tuvo que liberarse de la iglesia tanto en el terreno intelectual como en el práctico.

\textsc{La era preclásica}

\underline{El mercantilismo}\\
El mercantilismo es el nombre que se da a la literatura y la práctica económicas del período que va de 1500 a 1750. Aunque hubo literatura mercantilista en todas las economías en desarrollo de Europa occidental, fueron los ingleses y los franceses los que hicieron las aportaciones más significativas.
Mientras que los autores de la literatura económica del escolasticismo son clérigos medievales, la teoría económica del mercantilismo es obra de comerciantes. Su literatura se refiere a cuestiones de política económica y está relacionada normalmente con el interés específico que el comerciante-autor trataba de promover. Por este motivo a menudo se veían con notable escepticismo los méritos analíticos de algunos argumentos y sobre la validez de sus conclusiones. Pocos autores podían asegurar que veían las cuestiones desde una distancia suficiente como para que su análisis fuera objetivo. Sin embargo, durante todo el período mercantilista aumentó tanto la cantidad de literatura económica como su calidad. La literatura mercantilista escrita entre 1650 y 1750 es claramente de mayor calidad; diseminados por toda ella se encuentran casi todos los conceptos analíticos en los que ADAM SMITH basó su obra Wealth of Nations, que publicó en 1776.
Se ha dicho que la era del mercantilismo fue un período en el que todo el mundo era su propio economista. Cada autor tendió a centrar la atención en un tema y ninguno fue capaz de sintetizar estas aportaciones lo suficiente para influir en el desarrollo posterior de la teoría económica, quizá porque la economía como disciplina intelectual aún no se había hecho un hueco en la universidad; era estudiada principalmente por hombres de negocios que escribían panfletos sobre los problemas económicos que les preocupaban.
Como mejor se entiende el mercantilismo es como una reacción intelectual a los problemas de la época. En este período de declive del feudo y de auge del estado-nación. Los mercantilistas trataron de averiguar cuáles eran las mejores medidas para aumentar el poder y la riqueza de la nación. Al igual que MAQUIAVELO, (estadista italiano teórico político y autor de El príncipe (1513)) había aconsejado a los gobernantes medidas políticas que les resultaran útiles, los mercantilistas les aconsejaron medidas económicas que consolidaran y aumentaran más el poder y la prosperidad de las economías en vías de desarrollo.

El mercantilismo se puede basar en las siguientes ideas:
1.
Los mercantilistas se basaban en el supuesto de que la riqueza total del mundo era fija. Basándose en este mismo supuesto, los escolásticos habían afirmado que cuando los individuos comerciaban, el beneficio que obtenía uno de ellos era necesariamente una pérdida para el otro (el comercio es un juego de suma cero). Los mercantilistas aplicaron este razonamiento al comercio entre naciones y llegaron a la conclusión de que la riqueza y el poder económico de una nación aumentaban a expensas de otras, por lo que pusieron el énfasis en el comercio internacional como medio para aumentar la riqueza y el poder de una nación y, especialmente, en la balanza comercial entre las naciones.
2.
En línea con la idea anterior, un país debía fomentar las exportaciones y desincentivar las importaciones por medio de aranceles, contingentes, subvenciones, impuestos, etc. con el fin de lograr la llamada balanza comercial favorable. Con ello, se provocaría la entrada de metales preciosos
en la economía nacional\footnote{Los primeros mercantilistas sostenían que había que conseguir una balanza comercial favorable con cada una de las demás naciones. Sin embargo, algunos autores posteriores mantenían que sólo era importante la balanza comercial global con todas ellas. Así, por ejemplo, Inglaterra podía tener una balanza comercial desfavorable con la India, pero como podía importar de la India materias primas baratas que podía utilizar para producir bienes en Inglaterra para la exportación, teniendo en cuenta todas las demás naciones podía muy bien tener una balanza comercial favorable.}. Para muchos de los primeros mercantilistas, la riqueza de una nación no quedaba definida por su producción ni por el consumo sino por la cantidad de metales preciosos que poseían\footnote{Una cuestión relacionada con ésta es la exportación de metales preciosos (bullion). Los primeros mercantilistas recomendaban que se prohibiera estrictamente la exportación de metales preciosos. Los autores posteriores sugirieron que la exportación de lingotes podía mejorar la balanza comercial global si se utilizaban para comprar materias primas para producir bienes para la exportación.}. Para ello, era vital estimular la producción por medio de la intervención gubernamental en la economía nacional y la regulación del gobierno exterior.
3.
El objetivo de la actividad económica era, según la mayoría de los mercantilistas, la producción y no el consumo como afirmaría más tarde la escuela clásica. Para los mercantilistas, la riqueza de la nación no era la suma de la riqueza de sus miembros. Abogaban por aumentar la riqueza de la nación fomentando simultáneamente la producción, aumentando las exportaciones y manteniendo bajo el consumo interior. Por tanto, la riqueza de la nación se basaba en la pobreza de la mayoría de la población.
a.
Aunque ponían mucho énfasis en la producción no consideraban deseable que hubiera abundante oferta de bienes dentro de un país. Un elevado nivel de producción, junto con un bajo nivel de consumo interior, permitirían aumentar las exportaciones, lo cual aumentaría la riqueza y el poder de la nación. Los mercantilistas era partidarios de que los salarios fueran bajos con el fin de dar a la economía nacional ventajas competitivas en el comercio internacional. Creían, además, que unos salarios superiores al nivel de subsistencia llevarían a reducir el esfuerzo: inducirían a los trabajadores a trabajar menos horas al año, por lo que la producción nacional disminuiría. Por tanto, cuando el objetivo de la actividad económica se define en función del producto nacional y no del consumo nacional, la pobreza del individuo beneficia a la nación.
4. Los primeros mercantilistas estaban muy impresionados por la importancia de la enorme entrada de metales preciosos del Nuevo Mundo en Europa, especialmente en España. Sin embargo, los mercantilistas posteriores no eran de esa opinión… LANDRETH y COLANDER pág. 46
La evaluación de autores anteriores plantea algunas cuestiones difíciles pero interesantes. Siempre existen diferencias de opinión sobre lo que quiso decir realmente un autor cuando lo dijo. ADAM SMITH, otros economistas clásicos y la línea ortodoxa del pensamiento económico desde 1776 hasta KEYNES apeas encontraron nada que mereciera la pena en una gran parte de la literatura mercantilista. Sinn embargo, J.M. KEYNES analiza el pensamiento mercantilista en un apartado de su General Theory (1936) titulado “Notas sobre el mercantilismo”. En este apartado les atribuye el mérito de haber sabido encontrar una política aceptable para estimular el desarrollo económico. Simpatizó con sus teorías subconsumistas y declaró acertada su creencia de que los aumentos de la cantidad de dinero incrementaban la producción. KEYNES afirmó que los mercantilistas sostenían que una balanza comercial favorable aumentaba el gasto interior y, por tanto, el nivel de renta y empleo.

\underline{Precursores del pensamiento económico clásico}\\
\textbf{Buscar sobre escuela de Salamanca}

\footnote{SANCHO DE MONCADA (1580-c.1638) fue un clérigo, universitario y economista mercantilista español que estudió y ocupó cátedras en la Universidad de Toledo. Su obra principal es Restauración política de España (1619) , que incluyen desde teoría política hasta pedagogía real, pero cuya parte esencial es el análisis de la situación económica y sus propuestas de solución, con implicaciones demográficas, monetarias y hacendísticas. Se ha considerado que su pensamiento económico, así como el de JUAN DE MARIANA, son precedentes de algunas ideas de la Escuela Austriaca de Economía, en su consideración de que la falta de información hace imposible que el gobierno organice coactivamente la sociedad civil.
https://es.wikipedia.org/wiki/Sancho_de_Moncada}

\underline{La fisiocrcacia}\\

\textsc{El paso a la era clásica}

El año 1776 es considerado un hito en la historia del pensamiento económico. Es en realidad, el punto donde arranca, y ya sin solución de continuidad, el pensamiento económico de la cultura occidental que, junto a la globalización del mundo de los negocios, se va haciendo dominante a escala mundial.
La actual sociología del conocimiento ha establecido en ese año la divisoria entre la economía científica y la especulativa. Pero se debe reconocer que el consenso de esa fecha está influenciado en demasía por una visión anglosajona y que anteriormente la Escuela de Salamanca y la Fisiócrata ya sentaron los precedentes para afrontar científicamente el estudio de los temas económicos. Las causas por las que hasta hace poco no se había considerado la existencia de una economía científica antes de 1776 fueron las explicaciones parciales de las actividades económicas, la ausencia de desarrollo posterior de las teorías, el olvido histórico en el caso de la Escuela de Salamanca y en el caso de la Escuela Fisiócrata, la brevedad de su duración (que casi fue la de una moda).
En ese año se publicó Indagación acerca de la naturaleza y causas de la riqueza de las naciones de ADAM SMITH.

\anexo{La equivalencia Ricardiana}\label{anx:equivalencia}
La equivalencia ricardiana, o la proposición de equivalencia Barro-Ricardo,\footnote{Posteriormente, Robert Barro publicó un artículo intitulado \textit{Are Government Bonds Net Wealth?} en el Journal of Political Economy (Vol. 82, 1974). Este modelo supone que las familias actúan como dinastías que viven hasta el infinito, debido al altruismo intergeneracional, que los mercados de capitales son perfectos (en el sentido de que todos pueden prestar y endeudarse a la misma tasa de interés) y que la senda de los gastos del gobierno está dada. En estas condiciones, si el gobierno financia los gastos mediante emisión de bonos de deuda, las familias dejarán donaciones a sus hijos justo lo suficientemente grandes como para compensar los mayores impuestos que se necesitarán para pagar esos bonos. Este artículo es una contribución importante a la Nueva Macroeconomía Clásica, construida en torno a la hipótesis de las expectativas racionales.
} es una teoría económica que sugiere que el déficit fiscal no afecta a la demanda agregada.

La argumentación en que se basa la teoría es la siguiente: el gobierno puede financiar su gasto mediante los impuestos cobrados a los contribuyentes actuales o mediante la emisión de deuda pública. No obstante, si elige la segunda opción, tarde o temprano tendrá que pagar la deuda subiendo los impuestos por encima de lo que estos se ubicarían en el futuro si otra fuera la elección. La elección es entre pagar impuestos hoy o pagar impuestos mañana. Supóngase, por ejemplo, que el gobierno decide financiar un gasto adicional a través de déficit, esto es, mediante cobrar impuestos mañana. RICARDO argumentaba que aunque los ciudadanos tienen más dinero hoy, ellos se darían cuenta de que tendrían que pagar impuestos mayores en el futuro y, por lo tanto, ahorrarán un dinero adicional para poder pagar los impuestos futuros. Este mayor ahorro por parte de los consumidores compensaría exactamente el gasto adicional del gobierno, de modo tal que la demanda agregada permanecerá inmodificada.

La teoría de la equivalencia ricardiana sugiere que los intentos del gobierno de influir sobre la demanda agregada mediante la política fiscal están condenados al fracaso. Esta idea se opone frontalmente a la teoría keynesiana, que afirma que la política fiscal, debido a los efectos del multiplicador de la renta, será efectiva logrando que los incrementos de déficit público logren incrementos mayores en proporción de la demanda agregada del gobierno actual.

Sin embargo, la investigación empírica rechaza la equivalencia ricardiana en su forma pura, aunque algunos estudios han encontrado efectos ricardianos en el comportamiento del ahorro. Para una revisión técnica de la literatura, véase BRIOTTI, \textit{Economic Reactions to Public Finance Consolidation: a Survey of the Literature}, 2005.



\end{document}